\documentclass[11pt]{article}

\usepackage[margin=1in]{geometry}
\usepackage{amsmath,amssymb,amsthm}
\usepackage{mathtools}
\usepackage{enumitem}
\usepackage[hidelinks]{hyperref}
\usepackage[expansion=false]{microtype}
\usepackage{xr}


\newtheorem{theorem}{Theorem}[section]
\newtheorem{lemma}[theorem]{Lemma}
\newtheorem{proposition}[theorem]{Proposition}
\newtheorem{corollary}[theorem]{Corollary}
\newtheorem{conjecture}[theorem]{Conjecture}
\theoremstyle{definition}
\newtheorem{definition}[theorem]{Definition}
\newtheorem{hypothesis}[theorem]{Hypothesis}
\newtheorem{remark}[theorem]{Remark}

\DeclareMathOperator{\lcm}{lcm}
\newcommand{\NN}{\mathbb{N}}
\newcommand{\ZZ}{\mathbb{Z}}
\newcommand{\QQ}{\mathbb{Q}}
\newcommand{\PP}{\mathcal{P}}
\newcommand{\der}{{}'}
\newcommand{\dd}{\,\mathrm{d}}

\title{The Pythagorean layers of Erd\H{o}s Problem~\#307:\\ a density theorem and the square sieve at the mass--two wall}
\author{Vico Bonfioli\thanks{\texttt{vicobonfioli@gmail.com}.
  See the acknowledgments for a candid note on the use of AI assistance in this work.}}
\date{Version 1.75.0\quad 2 September 2026}

\externaldocument[K-]{erdos307-core}
\externaldocument[A-]{erdos307-computational}
\externaldocument[B-]{erdos307-obstructions}
\begin{document}
\maketitle

\begin{abstract}
This note records the analytic route to Erd\H{o}s Problem~\#307 and how far it reaches. The
deviation ladder gives a second derivation of the barrier from the frontier side rather than from
the cycle. The main positive result is that both Pythagorean layers have density zero on the
squarefree stratum, with a rate. The square sieve at the mass--two wall is developed and does not
attain that rate: the loss is located in the character expansion, where the required uniformity in
the modulus is shown to be unavailable, and not in the arithmetic. Statements whose proofs rest on
inputs absent from the formalisation, Siegel--Walfisz and Hal\'asz among them, are labelled as such
and carry no machine check.
\end{abstract}
\section{The deviation ladder, the plus layer, and a second route to the barrier}\label{sec:ladderplus}

The material below accreted onto the ``note added'' above because it grew out of the imported
union criterion, and it stayed there. None of it is prior work: it is the deviation ladder, the
function--field places, the mod--$8$ law, the plus layer, the slot anatomy, and an independent
derivation of the barrier constant, and it includes two results carrying Lean formalisations
(Propositions~\ref{prop:noplace} and~\ref{prop:plusthin}). It is given its own section and
subsections here so that it can be found.

\subsection{The deviation ladder}\label{ssec:ladder}

\begin{proposition}[The deviation ladder]\label{prop:ladder}
Let $(a,b)$ be a Pythagorean pair with $a<b$ and parameter $k$. Writing
$\sigma(m)=\sum_{p\mid m}1/p$, one has $k=\sigma(a)-b/a=a/b-\sigma(b)\in\mathbb{Z}$, and:
\begin{enumerate}[label=\textup{(\arabic*)},leftmargin=2.2em]
\item $k\le0$, with $k=0$ exactly the \textup{(}open\textup{)} two--cycle case. Indeed $b'=a-bk>0$ and
$b>a\ge1$ force $k\le0$ \textup{(}if $k\ge1$ then $b'\le a-b<0$\textup{)}. Moreover
$|k|<\min(b/a,\sigma(b))\le\sigma(b)\le\log\log b+O(1)$, so the ladder has finitely many rungs; and,
since $\sigma(b)$ barely exceeds $2$ at the relevant scales, in practice only one or two.
\item \emph{(Parity.)} For odd squarefree $m$, $m'\equiv\omega(m)\pmod2$; for even squarefree $m$, $m'$
is odd. Hence if exactly one member is even, $k\equiv\omega(\text{odd member})\pmod2$; if both are odd,
$\omega(a)\equiv\omega(b)\pmod2$ and $k\equiv\omega(a)+1\pmod2$. The case $k=0$ recovers the parity
dichotomy of Theorem~\ref{K-thm:parity}.
\item Each rung is the system $a'=b+ak$, $b'=a-bk$, equivalently $(\sigma(a)-k)(\sigma(b)+k)=1$ (a
shifted copy of the cross--match $st=1$) so every rung is exactly as rigid as $k=0$. Since
$\sigma(a)+\sigma(b)=a/b+b/a$ is independent of $k$, Corollary~\ref{K-cor:emptytest} holds on \emph{every}
rung: the whole ladder is empty below $10^{112}$. (We make no quantitative claim that $k=0$ is easier
than the rest.)
\item On the rung $k=-1$ the equations read $a'=b-a$, $b'=a+b$ and $a'^2+b'^2=2(a^2+b^2)$: the
derivative acts on $z=a+bi$ as the orientation--reversing $\sqrt2$--similarity $z\mapsto(-1+i)\,\bar z$.
\item By uniqueness of the splitting \textup{(}Bado, Thm.~7.5\textup{)}, each squarefree $N$ with
$N'^2-4N^2=\square$ carries, under the convention $a<b$, a well--defined invariant
$k(N)\in\mathbb{Z}_{\le0}$; \#307 holds iff $k(N)=0$ for some $N$, and by~(3) the invariant lives on a
domain that is empty below $10^{112}$.
\end{enumerate}

\textup{Formalised.} \texttt{Erdos307.k\_nonpos} is the bound $k\le0$ of \textup{(1)}; the
finiteness of the ladder, $|k|<\min(b/a,\sigma(b))$, is not formalised. Further, \texttt{rung\_parity\_mixed} and
\texttt{rung\_parity\_odd} are \textup{(2)}, \texttt{rung\_cross} and \texttt{rung\_mass\_indep} are
\textup{(3)} \textup{(}the second being the invariance that carries
Corollary~\ref{K-cor:emptytest} onto every rung\textup{)}, and \texttt{rung\_minus\_one} is
\textup{(4)}. Corollary~\ref{cor:kdetermined} is \texttt{k\_determined}, with the confinement to
$k\in\{0,-1\}$ entering as a hypothesis rather than proved, and Bado's uniqueness of the splitting
in~\textup{(5)} cited \textup{(\texttt{lean/Erdos307/Ladder.lean})}.
\end{proposition}

The whole ladder is one equation in $\mathbb{Z}[i]$, and the rung index is a Gaussian multiplier.

\begin{proposition}[Gaussian multiplier form of the ladder]\label{prop:multiplier}
For coprime squarefree $a,b$ put $z=a+bi\in\mathbb{Z}[i]$ and let $\mathcal{D}z:=a'+b'i$ be the
coordinatewise arithmetic derivative. Then, unconditionally,
\[ z\cdot\mathcal{D}z\;=\;(aa'-bb')\;+\;i\,(ab)'. \]
Consequently $(a,b)$ is a Pythagorean pair if and only if $\mathcal{D}z=\mu\bar z$ for a Gaussian
integer $\mu$ with $\operatorname{Im}\mu=1$; that multiplier is $\mu=k+i$ with $k$ the deviation of
Proposition~\ref{prop:ladder}, and
\[ k=\frac{aa'-bb'}{a^2+b^2},\qquad
   a'^2+b'^2=(k^2+1)(a^2+b^2)=N(\mu)\,(a^2+b^2),\qquad \gcd(a',b')\mid k^2+1. \]
Hence \#307 asks for a Pythagorean pair whose multiplier is a \emph{unit}. Since
$\{k+i:k\in\mathbb{Z}\}\cap\mathbb{Z}[i]^{\times}=\{i\}$, the two--cycle rung $k=0$ is the unique rung
on which $\mathcal{D}$ acts as an isometry ($a'^2+b'^2=a^2+b^2$); every other rung is an expanding
anti--similarity of ratio $\sqrt{k^2+1}>1$, and so cannot be periodic. The condition
$\operatorname{Im}\mu=1$ is the Pythagorean equation and $\operatorname{Re}\mu=0$ is \#307: one
coordinate of a single Gaussian integer each.
\end{proposition}
\begin{proof}
Leibniz gives $\operatorname{Im}(z\,\mathcal{D}z)=a'b+ab'=(ab)'$ and
$\operatorname{Re}(z\,\mathcal{D}z)=aa'-bb'$, which is the identity (verified symbolically and on
$7{,}501$ coprime squarefree pairs). If $(a,b)$ is Pythagorean then $(ab)'=a^2+b^2=N(z)$, so
$z\,\mathcal{D}z=(aa'-bb')+i\,N(z)$; dividing by $N(z)=z\bar z$ gives
$\mathcal{D}z/\bar z=(aa'-bb')/N(z)+i$, which lies in $\mathbb{Z}[i]$ with real part the integer $k$ of
Proposition~\ref{K-prop:pairform}. Conversely $\mathcal{D}z=\mu\bar z$ with $\operatorname{Im}\mu=1$
forces $(ab)'=\operatorname{Im}(z\,\mathcal{D}z)=\operatorname{Im}(\mu)N(z)=a^2+b^2$. Substituting
$a'=b+ak$, $b'=a-bk$ gives $\mathcal{D}z=(k+i)\bar z$ identically, whence the norm identity
$|\mathcal{D}z|^2=|\mu|^2|z|^2$. Finally $g=\gcd(a',b')$ divides both
$k(b+ak)+(a-bk)=a(k^2+1)$ and $(b+ak)-k(a-bk)=b(k^2+1)$, and $\gcd(a,b)=1$.
\end{proof}

The ladder of Proposition~\ref{prop:ladder} has a twin, visible only after extending the derivative
to $\mathbb{Q}$ by the quotient rule $(a/b)'=(a'b-ab')/b^2$ \cite{UfnarovskiAhlander2003}.

\begin{proposition}[The co--ladder and the arithmetic Riccati equation]\label{prop:coladder}
For coprime squarefree $a\neq b$ the following are equivalent:
\textup{(R)} $u=a/b$ satisfies the arithmetic Riccati equation
\[ u'\;=\;1-u^2; \]
\textup{(W)} $a'b-ab'=b^2-a^2$; \textup{(L)} there is $k\in\mathbb{Z}$ with $a'=b+ak$ \emph{and}
$b'=a+bk$; the ladder of Proposition~\ref{K-prop:pairform} with the sign of one equation flipped.
The deviation cancels in the quotient: on \textup{(L)}, $a'b-ab'=(b+ak)b-a(a+bk)=b^2-a^2$
identically, so the Riccati equation is the $k$--free form of this \emph{second} ladder. On it:
\begin{enumerate}[label=\textup{(\roman*)},leftmargin=2.2em]
\item \emph{(Positivity barrier.)} $k\ge0$ always: with $a<b$, $k\le-1$ would give
$b'=a+bk\le a-b<0$. Hence $\sigma(a)+\sigma(b)=a/b+b/a+2k>2$, so every co--ladder pair has
$\omega(ab)\ge59$ and $ab>7.9\times10^{112}$: the co--ladder is empty at every reachable scale
\emph{by positivity alone}; a second transversal relaxation of the two--cycle carrying the full
barrier through a mechanism independent of the minus square. Positivity separates the two cases
further. If $k\ge1$ then $\sigma(a)=k+b/a>2$, since $a<b$, so the \emph{smaller} member alone
carries mass exceeding $2$: $\omega(a)\ge59$ and $a\ge\Pi_{59}>8.77\times10^{112}$, against
$\min(a,b)\ge2.09\times10^{56}$ for a two--cycle. A co--ladder point with $k\ge1$ is therefore
larger on its smaller member by a factor exceeding $4\times10^{56}$, and $k=0$ is the only case the
barrier permits below $\Pi_{59}$: the relaxation buys nothing in the range where anything could be
found, which is the conclusion Proposition~\ref{prop:nogain} reaches for the mass condition, here
for the ladder.
\item \emph{(The cycle is the intersection.)} A two--cycle is exactly a common point of the two
ladders at $k_{\mathrm{sym}}=k_{\mathrm{anti}}=0$: \#307 asks for $u\in\mathbb{Q}_{>0}\setminus\{1\}$
satisfying the Riccati equation together with the Pythagorean equation of
Proposition~\ref{K-prop:pairform}. In multiplier form the two ladders read
$\mathcal{D}z=kz+i\bar z$ (symmetric; norm $(k^2+1)N(z)+4kab$) and $\mathcal{D}z=(k+i)\bar z$
(antisymmetric; norm $(k^2+1)N(z)$), meeting at $\mathcal{D}z=i\bar z$.
\item \emph{(Inversion symmetry.)} $(1/u)'=1-(1/u)^2$ is automatic from \textup{(R)}; the
Riccati layer is symmetric in $a\leftrightarrow b$, as \textup{(L)} is.
\end{enumerate}
\end{proposition}
\begin{proof}
\textup{(W)}$\Leftrightarrow$\textup{(L)}: the equation rearranges to $b(a'-b)=a(b'-a)$, and
$\gcd(a,b)=1$ forces $a\mid a'-b$, $b\mid b'-a$ with equal quotients; the converse is the displayed
cancellation. \textup{(R)}$\Leftrightarrow$\textup{(W)} is the quotient rule. In (i), summing
$\sigma(a)=b/a+k$ and $\sigma(b)=a/b+k$ gives the mass bound (strict, as $a\neq b$), and mass
$>2$ forces $\ge59$ primes as in Theorem~\ref{K-thm:barrier}. In (iii),
$(1/u)'=-u'/u^2=(u^2-1)/u^2$. All identities were verified symbolically; exhaustively, no
co--ladder pair with $ab\le10^7$ exists ($3{,}918{,}188$ candidate pairs through the $a\mid a'-b$
prefilter; the barrier's prediction), and a first small--box probe over $\mathbb{Z}[i]$
(canonical $\pi'=1$) also finds none, though there the positivity mechanism is unavailable and the
layer's status is genuinely open (\texttt{code/riccati\_coladder.py}).
\end{proof}

The two ladders embed in a plane, and the barrier becomes a half--plane law.

\begin{proposition}[The deviation half--plane]\label{prop:halfplane}
For coprime squarefree $a\neq b$ call the pair \emph{integral} if both deviations
$s=(a'-b)/a$ and $t=(b'-a)/b$ are integers, and place it at $(s,t)\in\mathbb{Z}^2$; the
antisymmetric ladder is the line $t=-s$, the co--ladder is $t=s$, and a two--cycle is the origin.
Then
\[ \sigma(ab)\;=\;\frac ab+\frac ba+s+t\;>\;2+(s+t), \]
so \emph{every} stratum in the closed half--plane $s+t\ge0$ carries the full barrier:
$\omega(ab)\ge59$ and $ab>7.9\times10^{112}$. Both ladder barriers are the two central lines of
this one inequality, and \#307 sits at the origin; on the boundary of the protected region, at
the crossing of its two distinguished lines. The open half--plane $s+t\le-1$ is unprotected, and is
populated from $ab=30$ on: over $ab\le10^7$ exactly $21$ strata occur, all obeying the law, the
nearest being $(0,-1)$ and $(-1,0)$ at $(a,b)=(6,5)$ and $(5,6)$; off--axis strata occur
($(1,-13)$ at $(690,53)$), so integrality does not force one exact side. The $(0,-1)$ stratum with
$b=a'$ is the chain equation
\[ a''+a'=a,\qquad\text{equivalently}\qquad \sigma(a)\bigl(1+\sigma(a')\bigr)=1, \]
a companion of the cross--match $\sigma(a)\sigma(b)=1$ with the same shape and one shifted factor.
Its solutions up to $10^7$ are $\{6,\,42,\,15590,\,47058\}$ and they split in two. Every primary
pseudoperfect number whose predecessor is \emph{prime} solves it, trivially: then $a'=a-1$ and
$a''=1$ ($6,42,47058$; $1806$ fails, its predecessor $1805=5\cdot19^2$ being composite, and $2$
fails as $a''=0$). The remaining solution $15590=2\cdot5\cdot1559$ is not on the pseudoperfect line
at all: $15590'=10923=3\cdot11\cdot331$ and $10923'=4667$, a genuine two--step closure. It is the
first solution of the chain equation whose derivative is composite; the sequence
$6,42,15590,47058$ does not appear in the OEIS. All census claims verified exhaustively
(\texttt{code/deviation\_halfplane.py}).
\end{proposition}
\begin{proof}
$\sigma(a)=a'/a=b/a+s$ and $\sigma(b)=a/b+t$ sum to the display, with strict inequality by AM--GM
($a\neq b$); mass $>2$ forces $\ge59$ primes as in Theorem~\ref{K-thm:barrier}. The ladder
identifications are Propositions~\ref{K-prop:pairform} and~\ref{prop:coladder}; the census is a
finite computation.
\end{proof}

\begin{corollary}[Below mass $5/2$ the rung is not a choice but a parity]\label{cor:kdetermined}
Let $(a,b)$, $a<b$, be a Pythagorean pair with $\sigma(ab)<5/2$. Then $k\in\{0,-1\}$, and $k$ is
\emph{determined} by the parity of $\omega$:
\[
\begin{array}{ll}
\text{exactly one member even:} & k=0\iff\omega(\text{odd member})\ \text{even};\\
\text{both members odd:}        & k=0\iff\omega(a)\ \text{odd}.
\end{array}
\]
Hence, at these masses, the single--integer test of Proposition~\ref{K-prop:pyth} together with one
parity check is \emph{equivalent} to \#307, not merely necessary for it: a squarefree $N$ with
$N'^2-4N^2=\square$ whose splitting has the stated $\omega$--parity \emph{is} a two--cycle. In
particular the finite verifications of Proposition~\ref{K-prop:close59} and Remark~\ref{A-rem:tier60} are
decisive in both directions; a surviving candidate would have been a solution, not a candidate,
and the whole level--$59$/$60$ regime lies in this range, its masses being confined to
$(2,\,2.00235]$, where $r=b/a\le1.0497$.
\end{corollary}
\begin{proof}
$\sigma(ab)=r+1/r$ with $r=b/a>1$, and $r+1/r<5/2\iff r<2$. By
Proposition~\ref{prop:ladder}(1), $k\le0$ and $|k|<b/a=r<2$, so $k\in\{0,-1\}$; a set of two
consecutive integers is separated by parity, and Proposition~\ref{prop:ladder}(2) supplies that
parity from $\omega$. The stated equivalences are the two cases of that law at $k\equiv0$. The
mass range for $|U|\in\{59,60\}$ is Theorem~\ref{K-thm:barrier} and Proposition~\ref{K-prop:close59}.
\end{proof}

\begin{corollary}[Diagonal stratification: the cost of the last lattice step]\label{cor:diagonals}
Let $m(x)$ be least with $\sum_{i\le m(x)}1/p_i>x$. The diagonal $s+t=-j$ carries the barrier
$\sigma(ab)>2-j$, hence $\omega(ab)\ge m(2-j)$ and $ab\ge\Pi_{m(2-j)}$:
\[
\begin{array}{llll}
j\ge2: & \text{mass}>0, & \text{no barrier}, & ab\ge2;\\
j=1:   & \text{mass}>1, & \omega\ge3,        & ab\ge2\cdot3\cdot5=30;\\
j=0:   & \text{mass}>2, & \omega\ge59,       & ab\ge\Pi_{59}>8.77\times10^{112}.
\end{array}
\]
The deviation plane is thus stratified by mass, one diagonal at a time, and \#307 sits on the first
diagonal carrying a nontrivial barrier. The $j=1$ floor is \emph{attained}: the census minimum on
that diagonal is exactly $30$, at $(a,b)=(5,6)$ and $(6,5)$; so the barrier is sharp one step off
the origin. The final lattice step therefore costs
\[ \Pi_{59}/30\;\approx\;2.9\times10^{111}, \]
one hundred and eleven orders of magnitude for a single unit of $s+t$. This is the barrier of
Theorem~\ref{K-thm:barrier} in its most compressed form: not a bound on solutions, but the price of
the last step of an integer lattice walk, paid because the diagonal index is a mass threshold and
the mass thresholds $1$ and $2$ sit at $3$ and $59$ primes.
\end{corollary}

\emph{The populated strata calibrate the heuristic model.} The unprotected strata are the first family of \#307--type exact coincidences where the
average--case model can be tested against exhaustive truth. Pricing the second exactness condition
of a candidate pair at $1/b$ and re--running the census enumeration with identical filters gives
expected counts $11.4,\,18.3,\,26.1,\,34.2$ at $ab\le10^4,\dots,10^7$ against observed
$6,12,18,23$: the ratio is stable near $2/3$ over four decades; the model has the right order of
growth and a constant local correction, i.e.\ it is \emph{calibrated} at the one--condition level.
Composing two levels, however, already breaks the naive pricing: for the chain equation $a''+a'=a$
the crude spread price predicts $E=48.8$ against the observed $4$; a twelvefold bias at
reachable height, because the composition requires mass--window conditioning between the levels
(the stratified--versus--naive lesson of Section~\ref{A-sec:anatomy} of the computational companion, here measured rather than
inferred). The reliability of averaging thus degrades measurably with each composed exact
condition, and the two--cycle at the origin composes many. This is the quantitative face of the
average--to--instance gap: the model earns trust stratum by stratum, and loses it condition by
condition (\texttt{code/strata\_calibration.py}).


The Pythagorean layer behaves very differently over function fields, and refines into two arithmetic
sub--layers over $\mathbb{Z}$.

\begin{proposition}[The Pythagorean layer over function fields]\label{prop:ff-pyth}
Over $\mathbb{F}_q[t]$ with $q$ odd, $a^2+b^2=(ab)'$ has no coprime squarefree monic solutions:
$\deg(a^2+b^2)=2\max(\deg a,\deg b)>\deg a+\deg b-1\ge\deg((ab)')$, the leading coefficient $1$ or
$2\ne0$ preventing cancellation. Over $\mathbb{F}_2[t]$ the layer is nonempty: up to degree $8$
(exhaustively, $23{,}548$ coprime squarefree pairs) its only solutions are the three pairs of the
Artin--Schreier tower $x_{j+1}=x_j(x_j+1)$, namely $(t,t+1)$, $(t^2+t,t^2+t+1)$, $(t^4+t,t^4+t+1)$;
each has $a'=(a+1)'=1$ and $(a(a+1))'=(a+1)+a=1$, and the climb halts at level $4$ (there
$(x_3+1)'=t^2+t\ne1$). These pairs lie off the cycle stratum $k=0$, on the side the integer
normalisation of Proposition~\ref{prop:ladder}(1) forbids. Cycles ($k=0$) remain impossible over every
$\mathbb{F}_q[t]$ (Theorem~\ref{B-thm:ff}): the function--field world separates the Pythagorean layer
from the cycle layer, sharpening the integer obstruction to ``the archimedean order forces $k\le0$.''
\end{proposition}

The function--field immunity is more fragile than the characteristic: it is destroyed by inverting
a single place.

\subsection{Function--field places}\label{ssec:ffplaces}

\begin{proposition}[Inverting one place destroys the function--field obstruction]\label{prop:ffsunit}
Over the $S$--integers $\mathbb{F}_q[t,t^{-1}]$, whose unit group $\mathbb{F}_q^\times\cdot t^{\mathbb{Z}}$
is infinite, unit--valued derivative two--cycles exist. Over $\mathbb{F}_3$ the smallest has degree
$3$ and only \emph{two} prime classes:
\[ a=t^3+t^2+1=(t-1)(t^2-t-1),\qquad b=t^3-t^2-1=(t+1)(t^2+t-1), \]
with $D(a)=b$ and $D(b)=a$ under the assignment $(u_1,u_2)=(t,1)$ on each side, $a,b$ coprime
squarefree and both coprime to $t$.
\end{proposition}
\begin{proof}
$D(a)=t\,(t^2-t-1)+1\cdot(t-1)=t^3-t^2-1=b$ and $D(b)=t\,(t^2+t-1)+1\cdot(t+1)=t^3+t^2+1=a$, both
exact in $\mathbb{F}_3[t]$; the factorisations are as displayed and share no factor. The values
$t$ and $1$ are units of $\mathbb{F}_q[t,t^{-1}]$, so this is a unit--valued Leibniz derivative in
the sense used throughout Section~\ref{B-sec:ff} of the obstructions companion. Verified independently by exhaustive search and by
direct recomputation (\texttt{code/hunt/ff\_twist.rs}).
\end{proof}

\begin{remark}[What the function--field theorem is really about]\label{rem:ffplaces}
Proposition~\ref{prop:ff-pyth} is a \emph{degree} argument: with $\pi'=1$ one has
$\deg D(a)\le\deg a-1$, so $D(a)=b$, $D(b)=a$ forces $\deg b\le\deg b-2$. Sign twists cannot
disturb it, because $\mathbb{F}_q^\times$ consists of constants and degree is blind to them,
exactly as, over $\mathbb{Z}$, twisting by $\pm1$ retains the full barrier
(Remark~\ref{B-rem:ordered}). But allowing the values $t^{k}$, units once $t$ is inverted, makes
$\deg D(a)\le\deg a-1+K$, and the two--cycle inequalities collapse to $0\le2K-2$: the obstruction
is vacuous from $K=1$, and cycles duly appear at once. A census over $\mathbb{F}_3$ confirms both
halves: at $K=0$ the count is $0$ at every degree tested, and at $K=1,2$ it is
$4,16$ / $4,48$ / $12,122$ for degrees $\le4,5,6$.
The moral sharpens Section~\ref{B-sec:ff} of the obstructions companion's thesis rather than softening it. What immunises
$\mathbb{F}_q[t]$ is not the characteristic, nor the finiteness of its unit group (cycles exist
over $\mathbb{Z}[\sqrt{-2}]$, whose unit group has order $2$) but the presence of a
\emph{single} distinguished place carrying a degree function that the derivative strictly lowers.
Adjoining a second place destroys it. This is the exact structural analogue of the claim made for
$\mathbb{Z}$: one archimedean place, one mass function the derivative must respect, and a barrier
that survives every twist available in the ring.
\end{remark}

\begin{proposition}[No place of $\mathbb{Q}$ carries the function--field descent]\label{prop:noplace}
The mechanism isolated in Remark~\ref{rem:ffplaces} (a place whose valuation the derivative
strictly lowers) is unavailable over $\mathbb{Q}$, and this is a classification statement rather
than a failure of ingenuity.
\begin{enumerate}[label=\textup{(\roman*)},leftmargin=2.2em]
\item \emph{Archimedean place.} $|D(n)|=n\,\sigma(n)$, which exceeds $n$ for every $n$ of mass
$>1$; the derivative \emph{expands}, first at $n=30$, and on $1.79\%$ of squarefree $n\le10^6$.
The absolute value is in any case not ultrametric, so the second half of the degree argument
($\deg(x+y)\le\max$) is unavailable too.
\item \emph{$p$--adic places.} For squarefree $n$ with $p\mid n$, cofactor survival gives
$v_p(D(n))=0<1=v_p(n)$: the derivative lowers $v_p$ there. But it \emph{raises} it elsewhere,
there are squarefree $n$ with $p\nmid n$ and $p\mid D(n)$ (e.g.\ $p=2$, $n=15$; $p=3$, $n=14$;
$p=5$, $n=6$), in comparable abundance to the lowering cases at every $p$ tested. So no $p$--adic
place contracts uniformly.
\item By Ostrowski's theorem \textup{(i)} and \textup{(ii)} exhaust the places of $\mathbb{Q}$.
\end{enumerate}
Hence no valuation--theoretic descent of the type that proves Proposition~\ref{prop:ff-pyth} can
exist over $\mathbb{Q}$: the function--field argument fails to transfer for a structural reason
internal to the classification of absolute values, not for want of a cleverer place.
(Counts and witnesses: \texttt{code/no\_place.rs}.)

\textup{Formalised.} \texttt{Erdos307.not\_descent\_of\_strictMono} gives \textup{(i)} in a form
stronger than stated, killing \emph{every} strictly monotone valuation at once and not only the
absolute value, on the witness $30\mapsto31$ \textup{(}\texttt{der\_thirty}\textup{)};
\texttt{padic\_lowers\_on\_support} is the lowering half of \textup{(ii)} and is Rigidity, while
\texttt{raises\_two}, \texttt{raises\_three}, \texttt{raises\_five} check the three raising
witnesses outright. Ostrowski enters \texttt{no\_place\_descent} as an explicit disjunction
hypothesis rather than silently. $0$ \texttt{sorry}
\textup{(\texttt{lean/Erdos307/NoPlace.lean})}.
\end{proposition}

\subsection{The mod--$8$ law and the plus layer}\label{ssec:pluslayer}

\begin{proposition}[Deviation calculus for $k$--cycles]\label{prop:kdev}
For pairwise coprime squarefree $a_1,\dots,a_k$ with product $N$ and deviations
$\kappa_i=(a_i'-a_{i+1})/a_i$ (indices mod $k$): the ratios $a_{i+1}/a_i$ have product $1$
(telescoping), and $\sigma(N)=\sum_i a_{i+1}/a_i+\sum_i\kappa_i$. On the relaxed layer
$\sum_i\kappa_i=0$ (which contains the cycles $\kappa_1=\dots=\kappa_k=0$) AM--GM gives
$\sigma(N)\ge k$, so the distinct--prime barrier $m_k$ of Proposition~\ref{K-prop:kcycles} extends to the
relaxed problem for every $k$, and to genuine $k$--cycles for $k\le3$; for $k\ge4$ it inherits the
pairwise--coprimality caveat of Proposition~\ref{K-prop:kcycles} (the bound then holds for the multiset
reciprocal sum). At $k=2$, $\kappa_1=-\kappa_2$ is the deviation of Proposition~\ref{prop:ladder} and
$\sum\kappa_i=0$ is $a^2+b^2=(ab)'$.
\end{proposition}

\begin{proposition}[Split discriminant]\label{prop:splitdisc}
For coprime squarefree $a,b$ with $N=ab$,
\[ N'+2N-(a+b)^2=N'-2N-(a-b)^2=(ab)'-a^2-b^2. \]
Hence $N$ is the product of a Pythagorean pair iff $N'-2N$ and $N'+2N$ are \emph{both} perfect squares
$d^2,s^2$ (their product is Bado's discriminant $N'^2-4N^2$), with $a,b=(s\pm d)/2$. The minus factor
carries the obstruction: $N'-2N=(a-b)^2\ge0$ forces $\sigma(N)\ge2$, hence the $59$--prime barrier,
indeed $0$ of $103{,}596$ squarefree $N<2\times10^5$ satisfy $N'-2N=\square$
(Corollary~\ref{K-cor:emptytest} in one line). The plus factor is satisfiable at small scale ($114$ such
$N<2\times10^5$). In characteristic $2$ the obstructing factor vanishes ($a-b=a+b$) and the criterion
degenerates to $(a+b)^2=N'$: the permissiveness of Proposition~\ref{prop:ff-pyth} and the barrier of
Corollary~\ref{K-cor:emptytest} are the absence and presence of one and the same algebraic factor.
\end{proposition}

Both squares of the split criterion are pinned modulo $8$ by an elementary law that seems worth
recording on its own.

\begin{proposition}[The mod--$8$ law of the derivative]\label{prop:mod8}
For odd squarefree $m$ write $S(m)=\sum_{p\mid m}p$. Then
\[ m\,m'\;\equiv\;S(m)\pmod 8. \]
Consequently: \textup{(a)} modulo $2$ this is the parity law $m'\equiv\omega(m)$ used in
Theorem~\ref{K-thm:parity}; \textup{(b)} an all--odd two--cycle $a'=b$, $b'=a$ forces
\[ S(a)\;\equiv\;S(b)\;\equiv\;ab\pmod 8 \]
the two prime sums are congruent to each other and to the product; refining
Theorem~\ref{K-thm:parity}(ii); \textup{(c)} in the mixed case $a=2k$ ($k$ odd), the odd member obeys
$b\equiv k\bigl(1+2S(k)\bigr)\pmod{16}$ and $S(b)\equiv2kb\pmod8$, and the even member the companion
law $(2k)\,(2k)'\equiv2+4S(k)\pmod{16}$; \textup{(d)} for odd squarefree $N$ the plus quantity
satisfies $N'+2N\equiv N\bigl(S(N)+2\bigr)\pmod8$, so an odd plus--hit must have
$N(S(N)+2)\equiv1\pmod8$ when $\omega(N)$ is odd and $\equiv0,4\pmod8$ when $\omega(N)$ is even.
None of these congruences is an obstruction, consistently with Proposition~\ref{A-prop:nolocal},
but they thin every hunt over the odd sector by a fixed factor.

\textup{Formalised.} \texttt{Erdos307.mod8\_law} is the law itself, resting only on
\texttt{odd\_sq\_mod8}; \texttt{prime\_sums\_congruent} is \textup{(b)}; \texttt{parity\_law} is
\textup{(a)} and is stated over $\NN$ with $\bmod$, so it needs no \texttt{ZMod} API, and
\texttt{omega\_even\_of\_even\_derivative} is the step Theorem~\ref{K-thm:parity}\textup{(i)} uses.
Part~\textup{(b)}'s second half, $S(a)\equiv ab$, is
\texttt{prime\_sum\_eq\_product}; part~\textup{(d)} is \texttt{plus\_quantity\_mod8} with the
square--residue filter \texttt{plus\_hit\_residue} and the $\omega$--parity split
\texttt{plus\_hit\_residue\_split}, which supplies the $\equiv1$ versus $\equiv0,4$ dichotomy;
part~\textup{(c)} is
\texttt{mixed\_member\_mod16}, \texttt{even\_member\_mod16} and \texttt{mixed\_prime\_sum\_mod8},
the mod--$16$ statements chaining onto the law through its integer form \texttt{mod8\_law\_int}. All
four parts are therefore formal \textup{(\texttt{lean/Erdos307/Mod8.lean})}.
\end{proposition}
\begin{proof}
$q^2\equiv1\pmod8$ for odd $q$, so $m/q\equiv mq\pmod8$; summing over $q\mid m$ gives
$m'\equiv m\,S(m)\pmod8$, and multiplying by $m$ (with $m^2\equiv1$) gives the law. (a) is the law
modulo $2$, since $S(m)\equiv\omega(m)$ for odd primes. (b) is the law at $a$ and at $b$ with $a'=b$,
$b'=a$. In (c), Leibniz gives $b=(2k)'=k+2k'$, and the law at $k$ gives $k'\equiv kS(k)\pmod8$,
whence $b\equiv k+2kS(k)\pmod{16}$; the second congruence is the law at $b$ with $b'=2k$, and the
companion follows from $2k(k+2k')\equiv2k^2+4kk'\pmod{16}$ with $k^2\equiv1\pmod8$ and
$kk'\equiv S(k)\pmod8$. In (d), $s^2\bmod8\in\{0,1,4\}$ and $s\equiv N'\equiv\omega(N)\pmod2$.
All parts were verified exhaustively over the $405{,}285$ odd squarefree $m\le10^6$ and all odd
plus--hits below $10^6$: zero exceptions (\texttt{code/mod8\_plusthin.py}).
\end{proof}

\begin{proposition}[Structure of the plus--layer]\label{prop:plus}
Call squarefree $N$ with $\omega(N)\ge2$ a \emph{plus--hit} if $N'+2N$ is a perfect square; the first
are $130,145,305,689,1745,1870,1970,2353,\dots$ \textup{(}$114$ below $2\times10^5$; OEIS~\cite{OEIS_A396915}\textup{)}.
\textup{(i)} For a semiprime $N=pq$, $2(N'+2N)+1=(2p+1)(2q+1)$, so a plus--hit forces
$(2s)^2\equiv-2\pmod\ell$ for every prime $\ell\mid(2p+1)(2q+1)$, whence $\ell\equiv1,3\pmod8$ and
$p\equiv q\equiv1\pmod4$ (this explains the dominance of $5$ and $13$).
The identity inverts the search: a semiprime plus--hit is exactly a factorisation
$2r^2+1=(2p+1)(2q+1)$, so one scans $r\le\sqrt{7N/3}$ and sieves $2r^2+1$ in $r$, at cost
$O(\sqrt N)$ in place of $O(N)$. This gives $243{,}129$ semiprime plus--hits with $v\le10^{13}$ in
eighty seconds, and $p\equiv q\equiv1\pmod4$ holds for every one of them
\textup{(\texttt{code/plushit\_inverted.rs}, which reproduces the $1{,}787$ below $3\times10^8$
exactly, with no miss and no extra)}. \textup{(ii)} \textup{[}empirical\textup{]} $\#\{\text{plus--hits}\le X\}\approx0.25\sqrt X$
(ratios $0.250,0.253,0.255$ at $X=5\!\times\!10^4,10^5,2\!\times\!10^5$): the plus factor behaves like
a random integer of its size, so (with Proposition~\ref{prop:splitdisc}) the whole integer
obstruction sits in the minus factor. \textup{(iii)} Fixing $p=5$, $q=(s^2-5)/11$ with $s$ even and
$s\equiv\pm4\pmod{11}$ gives $q\equiv1\pmod4$, and each prime value yields the plus--hit $N=5q$
($68$ with $s<2000$, up to $N=1{,}774{,}805$); Bunyakovsky's conjecture for this quadratic thus
implies infinitely many plus--hits. The two discriminant factors separate cleanly: the plus--layer is
Bunyakovsky--class, the minus--layer is the barrier.
\textup{(iv)} [Remark] The full cycle condition states that $(d^2,N',s^2)$ is a three--term
arithmetic progression with common difference $2N$ whose outer terms are squares; a
congruent--number--shaped configuration; the elliptic--curve machinery of that theory does not
transfer because $N'$ is not an algebraic function of $N$. A $251$--term b--file to $N<10^6$
(counting ratio $0.251$) and an OEIS submission draft accompany this note; the plus--hit
sequence was not in OEIS when we searched; it is now A396915, submitted by us and cited above, so the entry is not independent corroboration.
\textup{(v)} [P, C] \emph{General local law (QR tournament).} For every odd prime $\ell\mid N$ of
a plus--hit, $N'+2N\equiv N/\ell\pmod\ell$ (all cofactor terms $N/p$ with $p\ne\ell$ carry a factor
of $\ell$; only $N/\ell$ survives), so $N/\ell$ is a nonzero quadratic residue mod $\ell$. Verified
with zero violations on all $114$ plus--hits below $2\times10^5$ and all $251$ terms to $10^6$;
the semiprime $p\equiv q\equiv1\pmod4$ theorem of~\textup{(i)} is its $\omega=2$ shadow via quadratic
reciprocity.
\end{proposition}

\textup{(vi)} \emph{Parity of $\omega$ on odd plus--hits, and its single failure.} Among the odd
plus--hits the number of prime factors is almost always even: to $2\times10^7$ all $598$ have
$\omega$ even, and to $10^8$ all but one of $1{,}252$ do. The exception is
\[ v=30735705=3\cdot5\cdot7\cdot11\cdot13\cdot23\cdot89,\qquad \omega(v)=7,\qquad
   v'+2v=89094721=9439^2, \]
the least odd plus--hit with $\omega$ odd. It is not isolated: to $3\times10^8$ there are three,
$30735705$, $165064305$ and $230839455$, all with $\omega=7$ and with $t=4,4,5$, against $2{,}143$
with $\omega$ even, a rate of $0.14\%$ \textup{(\texttt{code/plushit\_omega.rs})}. So the parity is a
bias of about $99.9\%$ and not a law,
which is worth recording because the bias is strong enough to look like one over any range short of
$3\times10^7$. What does hold without exception is the genus constraint: writing
$t=\#\{p\mid v:p\equiv3\bmod4\}$, every one of the $2{,}146$ to $3\times10^8$ has $t\equiv0$ or $1\pmod4$, as the
identity $\prod_{p\mid v}\bigl(\tfrac{v/p}p\bigr)=(-1)^{\binom t2}$ requires, the counterexample
included with $t=4$. The observed pairs $(\omega,t)$ are $(2,0)$ with $1{,}787$ instances, $(4,1)$ with $339$, $(6,4)$
with $17$, $(7,4)$ with two and $(7,5)$ with one.

The empirical $\tfrac14\sqrt X$ law of part (ii) can be proved, up to logarithms, on the stratum that
dominates the small--scale census.

\begin{proposition}[The semiprime plus--layer is provably thin]\label{prop:plusthin}
The number of semiprime plus--hits $N=pq\le X$ is $O\bigl(\sqrt X\,(\log X)^2\bigr)$, by a fully
elementary argument; Erd\H{o}s's bound $\sum_{s\le Y}\tau(2s^2+1)\ll Y\log Y$ \cite{Erdos1952}
sharpens this to $O(\sqrt X\log X)$. Under Bunyakovsky's conjecture the $p=5$ family of
Proposition~\ref{prop:plus}(iii) alone contributes $\gg\sqrt X/\log X$, so conditionally
\[ \frac{\sqrt X}{\log X}\;\ll\;\#\{\text{semiprime plus--hits}\le X\}\;\ll\;\sqrt X\,\log X: \]
the $\sqrt X$ order of the empirical law is genuine on the semiprime stratum, not an artefact of the
tested range.
\end{proposition}
\begin{proof}
For a semiprime plus--hit $N=pq\le X$ one has $s^2=N'+2N=p+q+2pq\le3X+1$ (as $p+q\le pq+1$), and
Proposition~\ref{prop:plus}(i) gives $2s^2+1=(2p+1)(2q+1)$: the map
$N\mapsto\bigl(s,\{2p+1,2q+1\}\bigr)$ into pairs (root, unordered nontrivial factorisation) is
injective, so the count is at most $\sum_{s\le Y}\tau(2s^2+1)$ with $Y=\sqrt{3X+1}$. The polynomial
$2s^2+1$ has discriminant $-8$, prime to every odd modulus, so its roots modulo an odd prime are
simple and lift uniquely by Hensel: $2s^2+1\equiv0\pmod u$ has at most $2^{\omega(u)}\le\tau(u)$
solutions modulo $u$ (none for even $u$). Pairing each divisor of $2s^2+1$ with its cofactor and
keeping the smaller member $u\le\sqrt{2s^2+1}\le\sqrt2\,Y$,
\[ \sum_{s\le Y}\tau(2s^2+1)\;\le\;2\!\!\sum_{u\le\sqrt2Y}\!\!\#\{s\le Y:u\mid2s^2+1\}+Y
\;\le\;2\!\!\sum_{u\le\sqrt2Y}\!\!\tau(u)\Bigl(\frac Yu+1\Bigr)+Y\;\ll\;Y(\log Y)^2, \]
by Dirichlet's $\sum_{u\le Z}\tau(u)\sim Z\log Z$ and $\sum_{u\le Z}\tau(u)/u\sim\tfrac12(\log Z)^2$.
For the conditional lower bound, the family of Proposition~\ref{prop:plus}(iii) has $\asymp\sqrt X$
eligible $s\le\sqrt{2.2X+5}$ (two residue classes modulo $22$), of which Bunyakovsky for
$(s^2-5)/11$ makes $\gg\sqrt X/\log X$ prime. The injection, the Hensel root count, and the censuses
were verified numerically ($59$ semiprime hits below $2\times10^5$, all reproduced;
$\sum_{s\le Y}\tau(2s^2+1)/(Y\log Y)$ stable near $1.05$--$1.09$ for $Y=500$--$2000$;
\texttt{code/mod8\_plusthin.py}). The $\omega\ge3$ stratum remains empirical.
\end{proof}

\emph{What of the last argument is formal.} The proof splits into an elementary skeleton and an analytic tail, and the split is worth making
explicit because the two have different status. The skeleton is formalised:
\texttt{Erdos307.plus\_semiprime\_identity} is the factorisation $2s^2+1=(2p+1)(2q+1)$,
\texttt{plus\_range} the bound $s^2\le3N+1$ that confines $s$ to $Y$, \texttt{plus\_recover} the
injectivity, and \texttt{count\_le\_divisor\_sum} the reduction of a count to
$\sum_{s\le Y}\tau(2s^2+1)$ by injection into the disjoint union of the divisor sets. That reduction
is stated abstractly, over any finite set carrying a root and a divisor: the notion ``plus--hit''
is not itself defined in Lean, so the instantiation to this layer is done in the prose above and not
machine--checked. \texttt{card\_large\_le\_card\_small} and \texttt{tau\_le\_two\_small} are the
divisor--pairing step that reorganises that sum over the small member alone, and
\texttt{two\_roots\_quadratic} is the Hensel root count at the prime level, in a form needing only
$2\ne0$ rather than the discriminant. $0$ \texttt{sorry}
\textup{(\texttt{lean/Erdos307/PlusThin.lean})}.

The rest of the elementary chain is formal as well. The count of $s\le Y$ in the admissible residue
classes is \texttt{count\_in\_residue\_classes}, and the divisor sum is bounded in the printed shape
by \texttt{sum\_tau\_div\_le\_log\_sq}, $\sum_{u\le Z}\tau(u)/u\le(1+\log Z)^2$, from Mathlib's
\texttt{harmonic\_le\_one\_add\_log}. Only Dirichlet's \emph{asymptotics} are absent from the
library, and the upper bound does not use them.

The step that was missing is now supplied. The chain needs $|R|\le\tau(u)$ for $R$ the roots of
$2s^2+1\equiv0$ modulo $u$, and the half $2^{\omega(u)}\le\tau(u)$ was already
\texttt{two\_pow\_omega\_le\_tau}. The other half, $|R|\le2^{\omega(u)}$, is the multiplicativity of
the root count, and is now \texttt{Erdos307.rootSet\_card\_le\_two\_pow\_omega}, with the
composition \texttt{rootSet\_card\_le\_tau}. Two steps, both elementary in $\mathbb{Z}/u$: the
Chinese remainder ring equivalence carries a root modulo $ab$ to a pair of roots injectively, so the
count is submultiplicative on coprime factors \textup{(}\texttt{rootSet\_card\_mul\_le}\textup{)};
and modulo an odd prime power there are at most two roots, because $(x-y)(x+y)=0$ while $p$ cannot
divide the representatives of both factors, their sum representing the unit $2x$
\textup{(}\texttt{rootSet\_card\_prime\_pow\_le\_two}\textup{)}. Modulo a power of $2$ there are
none. $0$ \texttt{sorry} \textup{(\texttt{lean/Erdos307/RootCount.lean})}.

What remains outside Lean is the \emph{instantiation}, not a lemma: ``plus--hit'' is not defined in
Lean, so tying the abstract reduction to this layer is done in the prose above. Erd\H{o}s's
sharpening and Bunyakovsky's lower bound stay cited, and neither is used by the upper bound.


The semiprime argument generalises: the slot calculus of Proposition~\ref{prop:slot} counts.

\begin{proposition}[Slot counting: every fixed--$\omega$ stratum of the plus--layer is sparse]\label{prop:strata}
Fix $r\ge2$. The number of plus--hits $N\le X$ with $\omega(N)=r$ is $O\bigl(X^{1-1/r+o(1)}\bigr)$.
\end{proposition}
\begin{proof}
Write $p=P^+(N)$ for the largest prime of $N$ and $M=N/p$, so $\omega(M)=r-1$ and, since $N<p^r$,
$M\le N^{1-1/r}\le X^{1-1/r}$. The slot identity of Proposition~\ref{prop:slot} (base $M$, slot $p$)
gives $N'+2N=p\,E_0(M)+M$ with $E_0(M)=M'+2M$, so a plus--hit $s^2=N'+2N$ forces
\[ s^2\equiv M\ (\mathrm{mod}\ E_0(M)),\qquad p=\frac{s^2-M}{E_0(M)}, \]
and the map $N\mapsto(M,s)$ is injective. For squarefree $M$ one has $\gcd(M,M')=1$ (modulo
$q\mid M$ only the cofactor $M/q$ survives in $M'$), hence $\gcd(M,E_0(M))=1$: at every odd prime
$\ell\mid E_0(M)$ a root of $s^2\equiv M$ has $\ell\nmid s$, is nonsingular, and lifts uniquely
(Hensel), so the congruence has at most $4\cdot2^{\omega(E_0)}=X^{o(1)}$ roots modulo $E_0(M)$.
Since $s^2=N\bigl(\sigma(N)+2\bigr)$ and $\sigma(N)\ll\log\log X$, each base carries at most
$X^{o(1)}\bigl(X^{1/2}/E_0(M)+1\bigr)$ admissible $s$, and with $E_0(M)\ge2M$,
\[ \#\{N\}\;\le\!\!\sum_{\substack{M\le X^{1-1/r}\\ \omega(M)=r-1}}\!\!X^{o(1)}\Bigl(\frac{X^{1/2}}{M}+1\Bigr)
\;\ll\;X^{1/2+o(1)}+X^{1-1/r+o(1)}\;=\;X^{1-1/r+o(1)}. \]
(For $r=2$ this recovers the order of Proposition~\ref{prop:plusthin}, which is sharper by
logarithms.) The injection, the identity, $\gcd(M,M')=1$ and $P^+(N)\ge N^{1/r}$ were verified on
all $251$ plus--hits below $10^6$: zero exceptions (\texttt{code/strata\_count.py}).
\end{proof}

\emph{What separates the strata from the full layer.} The bound is honest but far from the truth: at $X=10^6$ the strata $\omega=2,\dots,5$ hold
$124$, $69$, $50$, $8$ hits (every one below $\sqrt X$) because the congruence
$s^2\equiv M\pmod{E_0(M)}$ is usually \emph{insoluble}: over random odd semiprime bases the average
number of roots is $0.34$, the no--residue branch of the trichotomy of Proposition~\ref{prop:slot}
dominating. Summing the proposition over $r$ does \emph{not} give the full layer density zero,
because the $X^{o(1)}$ root--count factors are not uniform in $r$; an average bound for
$\tau(M'+2M)$ (a Titchmarsh--divisor--type sub--problem) remains the route to the true
($\sqrt X$--order) count. Density zero itself, however, needs less: a smooth/rough dichotomy
bypasses the average entirely, because roughness buys a factor $X^{-1/\log\log X}$ against which
the divisor bound costs only $X^{\log2/\log\log X}$, and $\log2<1$.


\begin{proposition}[The plus--layer has density zero]\label{prop:pluszero}
Unconditionally, with $\log_3=\log\log\log$,
\[ \#\{\text{plus--hits }N\le X\}\;\ll\;\frac{X}{(\log X)^{(1+o(1))\log_3X}}. \]
In particular the plus--layer has density zero. The proof applies verbatim to
$\{N\ \text{squarefree}:\ N'+cN=\square\}$ for any fixed $c\ge1$.
\end{proposition}
\begin{proof}
Split at $y=X^{1/u}$ with $u=\log\log X$. \emph{Smooth part.} The number of $N\le X$ with
$P^+(N)\le y$ is $\Psi(X,y)=X\,u^{-u(1+o(1))}=X(\log X)^{-(1+o(1))\log_3X}$ (Rankin's bound;
\cite{Tenenbaum}, III.5). \emph{Rough part.} If $P^+(N)=p>y$, the injection of
Proposition~\ref{prop:strata} applies with $M=N/p\le X/y=X^{1-1/u}$: the hit determines $(M,s)$
with $s^2\equiv M\pmod{E_0(M)}$, a congruence with at most
$\rho(E_0)\le4\cdot2^{\omega(E_0)}\le\exp\bigl((\log2+o(1))\log X/\log\log X\bigr)$ roots by the
maximal order of the divisor function \cite{HardyWright}; each base therefore carries at most
$\rho(E_0)\bigl(X^{1/2+o(1)}/E_0(M)+1\bigr)$ values of $s$. Summing, with $E_0(M)\ge2M$:
\[ \#\{\text{rough hits}\}\;\ll\;X^{1/2+o(1)}\;+\;X^{1-1/u}
\exp\Bigl((\log2+o(1))\frac{\log X}{\log\log X}\Bigr)
\;=\;X^{1/2+o(1)}+X\exp\Bigl(-(1-\log2-o(1))\frac{\log X}{\log\log X}\Bigr), \]
and $1-\log2=0.3068\ldots>0$: the rough part is smaller than every $X/(\log X)^A$. The smooth part
dominates and gives the stated rate. (At $X=10^6$ the split reads $251=225$ rough $+\,26$ smooth,
the injection collision--free; \texttt{code/pluszero.py}.)
\end{proof}

On a derivative \emph{line} the slot is not merely constrained; it is pinned, and the counting
sharpens to every line at once.

\subsection{Slot and line anatomy}\label{ssec:slotanatomy}

\begin{proposition}[The determined slot: all derivative lines are thin, uniformly]\label{prop:lines}
For integers $e\ge1$ and $c$,
\[ \#\{N\le X\ \textnormal{squarefree}:\ N'=eN+c\}\;\ll\;\frac{X}{(\log X)^{(1+o(1))\log_3X}}, \]
uniformly in $e$ and $c$. In particular the primary pseudoperfect line $n'=n-1$, the Giuga
quotient--one line $n'=n+1$, every line $n'=n\pm d$ of
\cite{GrauOllerMarcenSadornil2021}, and every slope--two line $n'=2n+c$ obey this bound.
(Continuity of the limit law of Proposition~\ref{A-prop:density} already forces each line to have
density zero; the content is the rate, far beyond what concentration inequalities give, and the
uniformity. Classical lines with extra divisibility structure may well admit sharper bounds; we
claim none.)
\end{proposition}
\begin{proof}
Split at $y=X^{1/u}$, $u=\log\log X$; smooth members number
$\Psi(X,y)=X(\log X)^{-(1+o(1))\log_3X}$ (Rankin). For a rough member $N=Mp$, $p=P^+(N)>y$,
Leibniz gives $N'-eN=p\,(M'-eM)+M$, and $M'-eM\ne0$: $\sigma(M)=e$ would make the auto--reduced
fraction of Theorem~\ref{K-thm:rigidity} an integer, forcing $D_M=1$, $M=1$. Hence
\[ p\;=\;\frac{c-M}{M'-eM} \]
is \emph{determined} by $M$ (the line has no free root at all) so $N\mapsto N/P^+(N)$ is
injective on rough members, whose number is at most $X^{1-1/u}$. The smooth part dominates.
Two classical objects fall out of the formula: on the stratum $M'-M=-1$ of the two slope--one
lines it reads $p=M+1$ ($c=-1$) and $p=M-1$ ($c=+1$) (the oblong chains
$6\mapsto42\mapsto1806$ and $30=6\cdot5$, $1722=42\cdot41$) while off--stratum members stay
slot--determined ($858=66\cdot13$ with $M'-M=-5$; $47058=1518\cdot31$ with $M'-M=-49$). Verified:
the formula reproduces $P^+(N)$ on all listed members and identically on the $51{,}201$ squarefree
$N\le10^5$ with $\omega\ge2$ (\texttt{code/lines\_slot.py}).
\end{proof}

The members of a line carry far more structure than thinness: every prime factor is pinned modulo
its cofactor, and the top of the factorisation is capped by the bottom.

\begin{proposition}[Line anatomy: pinning, coupling, caps]\label{prop:lineanatomy}
Let $M$ be squarefree on the line $M'=eM+c$, let $p>q$ be its two largest primes and $R=M/(pq)$.
\begin{enumerate}[label=\textup{(\alph*)},leftmargin=2.2em]
\item \emph{(Pinning.)} $M/t\equiv c\pmod t$ for \emph{every} prime $t\mid M$; the Giuga
congruence of the line, independent of $e$. Equivalently, each prime factor lies in one determined
residue class modulo its cofactor: $t\equiv c\,(M/(tr))^{-1}\pmod r$ for all primes $r\mid M/t$.
\item \emph{(Top cap.)} If $M/p\neq c$ then $p\le\sqrt M+|c|$: line members are
$(\sqrt M+|c|)$--smooth.
\item \emph{(Coupling.)} $pq\mid R(p+q)-c$. Here $\gcd(pq,c)=1$ is automatic: $p\mid c$ together
with $p\mid qR-c$ would give $p\mid qR$, impossible as $p$ exceeds every prime of $qR$.
\item \emph{(Second cap.)} If $R(p+q)\neq c$ then $q<R+\sqrt{R^2+|c|}$; the degenerate stratum
$R(p+q)=c$ forces $c\ge2\sqrt{RM}$, so it is empty on every line with $c<2\sqrt M$; in
particular on all the near--critical minus lines.
\end{enumerate}
\end{proposition}
\begin{proof}
(a) Cofactor survival gives $M'\equiv M/t\pmod t$ while $eM\equiv0$; the second form is CRT.
(b) $M/p-c$ is a multiple of $p$; if nonzero, $M/p\ge p-|c|$. (c) $p\mid qR-c$ and $q\mid pR-c$
give $pq\mid(qR-c)(pR-c)=pqR^2-cR(p+q)+c^2$, so $pq\mid c\,(R(p+q)-c)$, and the displayed
coprimality removes $c$. (d) If $R(p+q)\neq c$ then $pq\le|R(p+q)-c|\le R(p+q)+|c|<2Rp+|c|$,
whence $q^2<2Rq+|c|$; if $R(p+q)=c$ then $c\ge2R\sqrt{pq}=2\sqrt{RM}$ by AM--GM. All parts were
verified on the $7{,}184{,}234$ squarefree $n\le2\times10^7$ with $\omega\ge3$, for both $e=1,2$
with $c:=n'-en$ (zero exceptions, the second cap attained with slack $1$ at $n=70$) and on
the classical members, where the coupling quotient $\bigl(R(p+q)-c\bigr)/pq$ equals $1$ on the
oblong--chain members and $5$ on both off--stratum members $47058$ and $66198$
(\texttt{code/line\_anatomy.py}).
\end{proof}

The lines themselves close into a family under peeling, and the closure explains why no
reorganization of the counting has broken through.

\begin{proposition}[The peeling groupoid]\label{prop:groupoid}
For integers $a\ge1$, $b$, $c$ let $L(a,b,c)=\{n\ \textnormal{squarefree}:\ an'=bn+c\}$.
\begin{enumerate}[label=\textup{(\roman*)},leftmargin=2.2em]
\item \emph{(Closure; the slope is a mass.)} If $n=qm\in L(a,b,c)$ with $q=P^+(n)$, then
$m\in L(aq,\,bq-a,\,c)$, and the slope moves by $b/a\mapsto b/a-1/q$: the full peel chain of $n$
descends the partial sums of $\sigma(n)$; the classical lines, their $\mu$--Sondow shifts, and
the slope--two lines of the minus square are one orbit family.
\item \emph{(Rigidity at every level.)} Along the chain of $n$ on an integer line $(a,b)=(1,e)$,
the level--$j$ denominator vanishes iff $\sigma(m_j)=b_j/a_j$, which telescopes to the single
$j$--independent condition $\sigma(n)=e$; impossible, since an integer mass forces $n=1$ by
Theorem~\ref{K-thm:rigidity}. So every denominator on every genuine chain is nonzero.
\item \emph{(Determined slot at every level.)} Consequently $q=(c-am)/(am'-bm)$ at each step:
a line member is a tower of single determined candidates over its own core.
\end{enumerate}
Every partial--peel reorganization of the line counts (peeling one prime, several, in any
order, from either layer of the split criterion) stays inside this groupoid, where part
\textup{(iii)} reproduces the one--candidate--per--cofactor obstruction verbatim. The obstruction
is therefore \emph{invariant} under the natural changes of viewpoint: a proof must inject input
from outside the groupoid; primality along determined thin sequences, or a decomposition of
integers not indexed by their primes.
\end{proposition}
\begin{proof}
(i) is Leibniz: $a(qm'+m)=bqm+c$ gives $aq\,m'=(bq-a)m+c$, and $(bq-a)/(aq)=b/a-1/q$. In (ii),
$\sigma(m_j)=b_j/a_j=e-\sum_{i<j}1/q_i$ says $\sigma\bigl(m_j\prod_{i<j}q_i\bigr)=\sigma(n)=e$;
by the auto--reduced fraction of Theorem~\ref{K-thm:rigidity} an integer mass forces $D=1$. (iii) is
the display above (i). Verified: closure and slot on $10^5$ random $(n,a,b)$ (the only slot
exceptions being exactly the $\sigma(m)=b/a$ locus, excluded on genuine chains, $73$ of $10^5$
random triples, $0$ of $310{,}788$ genuine chain steps at $e=1,2,3$), and the slope telescope
exact on every chain (\texttt{code/peeling\_groupoid.py}).
\end{proof}

The determined slot bounds more than itself: unwound one level further it constrains the
\emph{second} prime by a mass.

\subsection{A second derivation of the barrier constant}\label{ssec:secondbarrier}

\begin{proposition}[Mass--defect bound on the second prime]\label{prop:massbound}
Let $M$ be squarefree with $\omega(M)\ge3$ on the line $M'=eM+c$, and write $p=P^+(M)$, $m=M/p$,
$P=P^+(m)$, $n=m/P$. If $\sigma(m)<e$ then
\[ P\bigl(e-\sigma(n)\bigr)\;<\;2-\frac cm,\qquad\text{equivalently}\qquad
   P^2\,(en-n')\;<\;2Pn-c. \]
For $e=2$ the hypothesis is automatic below $\Pi_{59}$; that is, throughout the range in which the
minus layer is counted. On the near--critical slope--two lines ($|c|$ small) the bound reads
$P<2/(2-\sigma(n))+O(|c|/n)$: \emph{the second--largest prime of $M$ is bounded by twice the
reciprocal of the mass defect of the remaining cofactor.}
\end{proposition}
\begin{proof}
$M$ squarefree gives $p>P$. The determined slot of Proposition~\ref{prop:groupoid} is
$p=(c-m)/(m'-em)$, and $\sigma(m)<e$ makes the denominator negative, so
$p=(m-c)/\bigl(m(e-\sigma(m))\bigr)$ and $p>P$ reads $P\,m(e-\sigma(m))<m-c$, i.e.\
$P(e-\sigma(m))<1-c/m$. Since $m=Pn$ and $\sigma(m)=\sigma(n)+1/P$, the left side is
$P(e-\sigma(n))-1$, which gives the first display; multiplying by $Pn$ clears denominators.
Verified with $0$ violations on every squarefree $M\le2\times10^6$ with $\omega\ge3$, for
$e=1,2,3$; at $e=1$ the $21{,}146$ exceptions are exactly the $\sigma(m)\ge1$ cases excluded by the
hypothesis, and the bound is within a factor $2$ of equality for $11\%$ of members and within $100$
for $59\%$ (\texttt{code/slot\_massbound.rs}).
\end{proof}

\begin{remark}[The exact form, and a window that is real but inert]\label{rem:masswindow}
Proposition~\ref{prop:massbound} has an exact ancestor. With $N=pm$, $p=P^+(N)$, $k=2m-m'$, Leibniz
gives $N'=pm'+m$, so on the minus layer
\[ d^2\;=\;N'-2N\;=\;m-pk,\qquad\text{whence}\qquad
   \sigma(N)\;=\;2+\frac1{P^+(N)}-\frac km,\qquad k\ge1. \]
Since $k\ge1$ this is a \emph{two--sided} window; the barrier supplies only the lower half:
\[ 2\;<\;\sigma(U)\;<\;2+\frac1{\max U},\qquad\text{equivalently}\qquad \sigma\bigl(U\setminus\{\max U\}\bigr)<2, \]
so a minus--layer support is \emph{mass--minimal}: deleting any one of its primes drops the mass below
$2$ (the largest is the binding case). The upper half is nonetheless \emph{inert} where it can be
tested: all $49{,}961$ admissible level--$59$ supports already satisfy it, the tightest margin being
$1.26\times10^{-3}$, so it prunes nothing there. We record it because it is exact, because it is the
identity behind Proposition~\ref{prop:massbound}, and because it may bind at higher levels, where the
largest admissible prime grows. (The enumeration also reproduces the count $49{,}961$ of
Proposition~\ref{K-prop:close59} independently, from a different traversal;
\texttt{code/masswindow.rs}.)
\end{remark}

\begin{remark}[The minus layer is the anatomy of a quadratic polynomial]\label{rem:quadratic}
The identity $d^2=m-pk$ of Remark~\ref{rem:masswindow} can be read as a parametrisation rather than a
constraint. On the minus layer $d^2=N'-2N\ge0$ is a square, so
\[ m\;=\;pk+d^2, \]
and for \emph{fixed} $(p,k)$ the cofactor $m$ runs over the values of the quadratic polynomial
$f_{p,k}(d)=d^2+pk$. The minus layer is therefore not an arbitrary exact coincidence but a question
about the \emph{anatomy of quadratic--polynomial values}: for which $d$ does $f_{p,k}(d)$ carry
mass $\sigma=2-k/f_{p,k}(d)$, i.e.\ $\omega\approx58$ prime factors concentrated on the small
primes, with $P^+(f_{p,k}(d))<p$?

This is the first place in this circle where a \emph{named} toolkit applies. The anatomy of
polynomial values is a developed subject; Erd\H{o}s--Kac for $\omega(f(n))$, Halberstam--Richert
sieve bounds for almost--primes in polynomial sequences, and the divisor sums of \cite{Erdos1952}
already used in Proposition~\ref{prop:plusthin}; whereas ``primality of a determined value''
(Remark~\ref{rem:onequestion}) had none. What the toolkit must supply is unusual: not the typical
behaviour of $\omega(f(d))$, which is $\log\log$ by Erd\H{o}s--Kac, but a \emph{large deviation} to
$\omega\approx58$ with the prime factors forced onto the bottom of the spectrum, since only the
small primes carry mass. Existing polynomial--anatomy results control typical and almost--prime
behaviour, not deviations of this shape; so the reformulation names the tool without yet applying
it. We record it as the most promising translation of the residual question, and note that the
$p=5$ family of Proposition~\ref{prop:plus}(iii) is the plus--layer shadow of the same quadratic
mechanism. (Identity verified on all $1{,}606{,}956$ squarefree $M\le3\times10^6$ with $\omega\ge2$:
$0$ violations.)
\end{remark}

\begin{remark}[The residue amplifier is real, and exactly paid for]\label{rem:amplifier}
Sieving $f_c(d)=d^2+c$ by its roots exposes a mechanism the union--side reading of
Corollary~\ref{cor:qr} misses. A prime $\ell$ divides some value iff $-c$ is a quadratic residue
modulo $\ell$, which removes half the primes; the $2^{-\omega}$ penalty the corollary records.
But an \emph{admissible} odd $\ell$ has \emph{two} roots, so it divides values with density
$2/\ell$ rather than $1/\ell$: on the primes that survive, the weight is doubled. Only the penalty
has ever been counted, and one might expect the doubling to tilt the heuristic. It does not: the
two effects cancel exactly.

Per fixed $c$ the effect is genuine. Writing $\sigma_B$ for the mass on primes $\le B$, measured
over $d\le2\times10^6$ with $B=2\times10^4$: a generic $c=1009$ gives mean $\sigma_B=0.368$, a
friendlier $c=4966$ (admitting $35$ of the $46$ primes below $200$) gives $0.400$, against $0.452$
for random integers; a \emph{perfectly} friendly $c$, admitting every odd prime, would give
$\tfrac14+2\sum_{\ell>2}\ell^{-2}=0.654$, a factor $1.45$ above random. So friendliness helps, and
in the limit beats a random integer outright.

But friendliness is priced. Each odd $\ell$ is admissible for only half of the $c$, so a $c$
admitting the first $58$ odd primes has density $2^{-58}$, exactly cancelling the $2^{58}$ gained
from the doubled densities:
\[ 2^{-58}\cdot\tfrac12\prod_{i=2}^{59}\frac2{\ell_i}\;=\;\prod_{i=1}^{59}\frac1{\ell_i}, \]
the random--integer probability. Averaging over $c=1,\dots,400$ confirms it numerically: the mean
mass ratio is $0.9990$, and the tail ratios $\Pr[\sigma_B>T]$ are $0.998$, $0.996$, $0.978$ at
$T=0.8,1.0,1.2$; unity to within sampling, while individual $c$ range from $0.257$ to $0.635$
in mean. \emph{The quadratic reformulation is heuristically neutral}: it neither strengthens nor
weakens the \textsc{yes} lean of Remark~\ref{A-rem:killcount}, and the $2^{-59}$ filter of
Corollary~\ref{cor:qr} is correctly priced after all. What
Remark~\ref{rem:quadratic} buys is a toolkit, not a bias (\texttt{code/quad\_anatomy.rs},
\texttt{code/amplifier\_cost.rs}).
\end{remark}

The bound has enough force to regenerate the barrier on its own.

\begin{corollary}[A second, independent derivation of the barrier]\label{cor:secondbarrier}
Let $N$ be a minus--layer hit with $\omega(N)\ge3$; more generally a member of any line
$N'=2N+c$ with $c\ge0$. Write $p=P^+(N)$, $m=N/p$, $P=P^+(m)$, $n=m/P$. Then
$N>10^{112.94}$, without using the arithmetic--geometric mean inequality.
\end{corollary}
\begin{proof}
If $\sigma(m)\ge2$ then $m$ carries at least $59$ primes and $N>m\ge\Pi_{59}$ already. Otherwise
Proposition~\ref{prop:massbound} applies, and $c\ge0$ makes its right side at most $2$:
$P\bigl(2-\sigma(n)\bigr)<2$. Since $P>P^+(n)$ this gives
\[ 2-\sigma(n)\;<\;\frac2{P^+(n)},\qquad\text{i.e.}\qquad \sigma(n)>2-\frac2{P^+(n)}. \]
But $\sigma(n)\le\sum_{q\le P^+(n)}1/q$, so $P^+(n)$ must satisfy
$\sum_{q\le p}1/q>2-2/p$. The least such prime is $p=269$, the $57$th: at $p=269$ the sum is
$1.99505>1.99257$, while at $p=263$ it is $1.99133<1.99240$. Hence $\omega(n)\ge57$,
$P^+(n)\ge269$, and $n\ge\Pi_{57}>10^{108.07}$; with $p>P>P^+(n)$ the two further primes give
$N>\Pi_{57}\cdot271\cdot277>10^{112.94}$.
\end{proof}

\emph{Two mechanisms, one constant.} Corollary~\ref{cor:secondbarrier} reaches $10^{112.94}$, and $\Pi_{59}>8.77\times10^{112}$ is
$10^{112.94}$: the two derivations agree to two decimal places. They share no step.
Theorem~\ref{K-thm:barrier} divides the defining equation by $ab$ and applies AM--GM to
$a/b+b/a$; the present route never forms that quotient, using instead the determined slot of
Proposition~\ref{prop:groupoid}, the trivial inequality $p>P^+(N/p)$, and the additivity
$\sigma(m)=\sigma(n)+1/P$. What they share is the arithmetic input underneath (how many primes
Mertens needs to reach mass~$2$) which is why the constants coincide rather than merely being
comparable. The barrier is thus not an artefact of the mass inequality one happens to use: it is
the same Mertens threshold seen through two independent structures, and the critical set
$S_2=\{n:2-\sigma(n)<2/P^+(n)\}$ of the second structure is empty below $10^{108}$ (verified empty
for $n\le5\times10^6$, \texttt{code/critical\_set.rs}).


\begin{remark}[The minus side: coercive versus degenerate moduli]\label{rem:minusslot}
Proposition~\ref{prop:pluszero} rests on the plus modulus being \emph{coercive}:
$E_0^+(M)=M'+2M\ge2M$. The minus--square obeys the same slot identity with
$E_0^-(M)=M'-2M=M(\sigma(M)-2)$, a modulus that \emph{vanishes exactly on the mass--$2$ wall}: the
degeneration locus of the counting method is the barrier stratum itself; the counting
incarnation of the form degeneration $N_0(2-\sigma)\to0$ of Proposition~\ref{prop:form}. The
minus--layer is the union of the slope--two lines $N'=2N+d^2$ over square heights, and three facts
are unconditional. \textup{(i)} A rough minus--hit ($P^+(N)=p>y$) forces its base upward: if
$\sigma(M)<2-1/y$ then $d^2=p\,E_0^-(M)+M<M(1-p/y)\le0$, impossible; so $\sigma(M)\ge2-1/y$, and
once $y>(2-T_{58})^{-1}=793.7\ldots$ this yields $\omega(M)\ge59$ and $M\ge\Pi_{59}$: \emph{the
$59$--prime wall propagates from the hit to its base}, and $N\ge\Pi_{59}\,y$. \textup{(ii)} Bases
with $\sigma(M)\ge2+1/y$ have moduli $E_0^-\ge M/y$, and contribute as in
Proposition~\ref{prop:pluszero}. \textup{(iii)} The remaining bases are near--critical,
$\sigma(M)\in[2-1/y,2+1/y)$; exactly the low slope--two lines $M'-2M=\pm k$, $k\le M/y$. Each is
thin by Proposition~\ref{prop:lines}, but at a degenerate modulus the $d$--freedom is
$\sqrt X$--sized, and summing the uniform bound does not close the layer. What would:
\emph{square--root bounds on slope--two lines}, $\#\{M\le Z:M'-2M=\pm k\}\ll Z^{1/2+o(1)}$
uniformly in $k$. There is a second, sharper route to that bound, obtained by abandoning the
prime--indexed decomposition altogether. Split $M=uv$ by \emph{size} rather than by largest prime,
$u$ the coprime divisor nearest $\sqrt M$; Leibniz gives the exact bilinearisation
\[ M'-2M\;=\;v\,\delta(u)+u\,\delta(v),\qquad \delta(n)=n'-n, \]
so a slope--two line becomes a bilinear equation in $(u,v)$ with both variables of comparable size
and, for fixed $u$, the exact line $u\,v'-(2u-u')\,v=c$, i.e.\ $v\in L\bigl(u,\,2u-u',\,c\bigr)$
in the notation of Proposition~\ref{prop:groupoid}. Writing $N(u)$ for the number of $v$ on that
line in the balanced window $v\asymp M/u$,
\[ \#\{M\le Z\ \text{on the line}\}\;\le\;\sum_{u\le Z^{1/2}}N(u), \]
so the missing square--root bound follows from $N(u)\ll Z^{o(1)}$ on average; a
\emph{short--window line--multiplicity} hypothesis. This is a strictly sharper target than the
one--class--per--cofactor count: the summation length drops from $Z$ cofactors to $Z^{1/2}$ balanced
divisors (empirically $Z^{0.43}$, the median of $\log u/\log M$ over nearest--to--$\sqrt M$ coprime
splits being $0.432$), and the surviving object is bilinear in a short window; the shape on which
dispersion and Fouvry--Iwaniec methods operate, rather than a worst--case single--class count. The
measured multiplicity is exactly $1$ (mean $1.00$, no sampled line carrying two solutions), so the
hypothesis is not merely plausible but numerically flat; proving it, however, runs back into the
determined slot, since $L(u,2u-u',c)$ is itself a groupoid object. All identities verified
exhaustively on the $591{,}423$ squarefree $M\le10^6$, in exact rationals for the line form
(\texttt{code/bilinear\_balanced.py}). By
Proposition~\ref{prop:lineanatomy}(a) that missing bound is a \emph{one--class--per--cofactor}
prime count: a line member is $M=pm$ with $p$ prime in a single determined class modulo its own
cofactor $m$, so the count is $\sum_m\#\{p\le Z/m:p\equiv\alpha_c(m)\ (\mathrm{mod}\ m)\}$, and the
trivial estimate $Z/m^2+1$ per cofactor fails exactly on its $+1$: one candidate integer per $m$,
thinned only by the primality of that single determined value. Bounding worst--case single--class
counts summed over all moduli lies beyond Bombieri--Vinogradov and beyond Linnik uniformly; the
same species of count guards the true $\sqrt X$ order of the plus layer. The two layers' remaining
questions are one wall. The ledger: plus layer closed by coercivity; minus layer reduced to the
near--critical spectrum. (A small
unconditional companion, from Proposition~\ref{prop:mod8}: an odd member of $N'=2N+c$ has
$S(N)\equiv cN+2\pmod8$.)
\end{remark}

\emph{The primality can be dropped: the missing bound is a divisibility count.} The formulation above asks for a bound on how often a determined value is \emph{prime}, which is
what places it beyond Bombieri--Vinogradov. That requirement can be removed. The slot
$(c-m)/(m'-em)$ is an integer at all only when
\[ (m'-em)\ \big|\ (c-m), \]
and line members inject into the set of $m$ satisfying that divisibility, so
\[ \#\{m\le Z:\ (m'-em)\mid(c-m)\}\ \ll\ Z^{1/2+o(1)}\quad\text{uniformly in }(e,c) \]
already implies the square--root bound, with no prime--distribution input whatever.

The substitution costs almost nothing, because integrality is where the thinness lives.
Counting both conditions separately on the two lines whose members are known, over $m\le2\times10^7$
(so $\log Z=16.8$): the line $n'=n-1$ has $21$ integral slots of which $16$ are prime, and
$n'=n+1$ has $29$ of which $19$ are prime. Integrality alone is already $O(\log Z)$; primality
removes a further factor of about $1.3$ and no more. The prime--counting requirement was carrying
a constant, while the divisibility was carrying the whole saving.

What remains is therefore a purely multiplicative question about the derivative, of the same shape
as Proposition~\ref{prop:groupoid} but with the primality stripped out, and with roughly four orders
of slack against the target at the ranges where the count can be observed. We record the
substitution because it changes which literature the problem faces: divisibility counts for
$m'-em$ rather than primes in determined classes.


\section{The square sieve at the mass-two wall}\label{sec:sqsieve}

Remark~\ref{rem:masstwo} left the residue inside $\{\sigma(m)\ge2-\varepsilon\}$. We now attack that
set directly, by a change of question. The minus layer is $m'-2m=d^2$, so $\sigma(m)=2+d^2/m>2$:
it lives \emph{above} the wall, and the layer is exactly the set on which the arithmetic function
$m\mapsto m'-2m$ takes square values. Atom $\mathrm{A}9$ is therefore an instance of
\emph{how often is an arithmetic function a perfect square}, and for that question there is a tool:
Heath-Brown's square sieve. It has not previously been applied to the arithmetic derivative,
because $m'$ is not a polynomial and no Weil bound is available. The obstruction dissolves because
the relevant symbol factorises.

\begin{lemma}[Factorisation of the symbol]\label{lem:symbolfact}
Let $r$ be odd and let $m$ be squarefree with $\gcd(m,r)=1$ \textup{(}equivalently, no prime factor
of $m$ divides $r$\textup{)}.
Put $\sigma_r(m)=\sum_{p\mid m}p^{-1}\in\mathbb{Z}/r$. Then $m'\equiv m\,\sigma_r(m)\pmod r$ and
\[ \Bigl(\frac{m'-2m}{r}\Bigr)=\Bigl(\frac{m}{r}\Bigr)\Bigl(\frac{\sigma_r(m)-2}{r}\Bigr). \]
\end{lemma}
\begin{proof}
$m'=\sum_{p\mid m}m/p$ and $m/p\equiv m\,p^{-1}\pmod r$ for each $p\mid m$, since $p$ is invertible
modulo $r$. Summing gives $m'\equiv m\sigma_r(m)$, hence $m'-2m\equiv m(\sigma_r(m)-2)$, and the
Jacobi symbol is completely multiplicative in its numerator.
\end{proof}
Verified exhaustively for $r=101,1009,10007$ over all admissible squarefree $m\le3\times10^7$:
$54{,}514{,}790$ checks, $0$ violations \textup{(\texttt{code/sqsieve\_lemmas.rs})}.

The point of Lemma~\ref{lem:symbolfact} is that $\sigma_r$ is \emph{additive} over the distinct
prime divisors, so after a Gauss-sum expansion the second factor becomes a product over $p\mid m$
and the whole sum turns multiplicative. That is what puts it inside Hal\'asz's theorem.

\subsection{The square sieve and Theorem~\ref{thm:a9}}\label{ssec:sieve}

The argument recorded in this section is a square sieve at $P\asymp(\log N)^{1/4}$ whose analytic
input does not close in the form attempted there, and why Remark~\ref{B-rem:whatisleft} keeps this
material off the critical path; the rate is nonetheless proved, by a different route, in
Theorem~\ref{cor:a9rate}.
The density statement itself, Theorem~\ref{thm:a9}, is proved, by a different sieve at a far smaller
range (Proposition~\ref{prop:condrate}). The section occupied eight pages in the main line. It is
moved verbatim to Appendix~\ref{app:sieve}; nothing is omitted and every label is unchanged, so the
cross-references from \S\ref{ssec:pairavg} still resolve. What the argument establishes, what it
assumes, and exactly where the gap sits are stated there.

\subsection{Pair averages and the mass--two wall}\label{ssec:pairavg}

\begin{lemma}[$\sigma$ is injective on squarefree integers]\label{lem:sigmainj}
For squarefree $d,d'$, $\sigma(d)=\sigma(d')$ forces $d=d'$. Consequently
$F(d,d'):=D(d)\,d'-D(d')\,d=dd'(\sigma(d)-\sigma(d'))$ vanishes only on the diagonal.
\end{lemma}
\begin{proof}
By Theorem~\ref{K-thm:structure}, $\sigma(d)=D(d)/d$ is already in lowest terms, since
$\gcd(D(d),d)=1$ for squarefree $d$; so $d$ is recoverable from $\sigma(d)$ as its denominator.
\end{proof}
Verified over all $182{,}377$ squarefree $d\le3\times10^5$: $0$ collisions
\textup{(\texttt{code/sqsieve\_pairavg.rs})}. The lemma is formalised: in the notation of
Section~\ref{K-sec:lean} of the core note, \texttt{Erdos307.dprod\_eq\_of\_mass\_eq} and
\texttt{Erdos307.csum\_eq\_of\_mass\_eq} derive it from \texttt{rigidity\_coprime}, with $0$
\texttt{sorry} and axioms \texttt{propext, Classical.choice, Quot.sound}
\textup{(\texttt{lean/Erdos307/Injective.lean})}. Its label is therefore \textup{VERIFIED}.

\begin{theorem}[The average pair bound]\label{thm:pairavg}
Let $\mathcal{R}=\{pq:p\neq q\in\mathcal{P}\}$ and
$P(D,r)=\#\{(d,d')\sim D:\sigma_r(d)\equiv\sigma_r(d')\}$. Then
\[ \sum_{r\in\mathcal{R}}P(D,r)\ \ll\ |\mathcal{R}|\,D+D^{2+o(1)} . \]
\end{theorem}
\begin{proof}
As $\gcd(d,r)=\gcd(d',r)=1$, the condition $\sigma_r(d)\equiv\sigma_r(d')$ is exactly $r\mid F(d,d')$,
so $\sum_{r\in\mathcal{R}}P(D,r)=\sum_{(d,d')}\#\{r\in\mathcal{R}:r\mid F(d,d')\}$. On the diagonal
$F=0$ is divisible by every $r$, contributing $|\mathcal{R}|\cdot\#\{d\le D:\mu^2(d)=1\}$. Off it
$F\neq0$ by Lemma~\ref{lem:sigmainj}, and $|F|\le2D^2\sigma_{\max}$, so
$\#\{r\in\mathcal{R}:r\mid F\}\le\tau(F)\ll D^{o(1)}$, contributing $D^{2+o(1)}$.
\end{proof}
Measured against $D^2/r+D$ the average has ratio $1.011$ at $|\mathcal{R}|=78$
\textup{(\texttt{code/sqsieve\_pairavg.rs})}. Note the proof uses no equidistribution of $\sigma_r$
and no uniformity in $r$: only the divisor bound and Lemma~\ref{lem:sigmainj}, the latter being the
same rigidity that makes \#307 a two-cycle problem at all. The counting skeleton is formalised as
\texttt{Erdos307.pairavg\_bound}, with the divisor step \texttt{card\_filter\_dvd\_le} and the
double count \texttt{sum\_filter\_comm}; $0$ \texttt{sorry}, standard axioms
\textup{(\texttt{lean/Erdos307/PairAvg.lean})}. The arithmetic input enters only through the
hypothesis that $F$ vanishes exactly on the diagonal, which is Lemma~\ref{lem:sigmainj}, itself
\textup{VERIFIED} and transferred from prime sets to squarefree integers by
\texttt{Erdos307.eq\_of\_cross}, so that the hypothesis consumed by the skeleton is exactly the
statement the paper uses.

The decomposition demanded by Remark~\ref{rem:missingdecomp} can be bypassed entirely. Rather than
splitting $\sum_{m\le N}$ into bilinear pieces, one may reverse the roles of the two arguments of the
symbol and appeal to the quadratic large sieve, which averages over the modulus. This is the more
promising route and we record it, together with the one hypothesis it still needs.

\begin{lemma}[Reversal]\label{lem:reversal}
Write $a_m=m'-2m$ and $a_m=\pm s_m^2k_m$ with $k_m$ squarefree. For odd squarefree $r$,
\[ T_r(N)=\sum_{k}b_k^{+}\Bigl(\frac kr\Bigr)+\Bigl(\frac{-1}r\Bigr)\sum_{k}b_k^{-}\Bigl(\frac kr\Bigr)
   \;+\;O\bigl(\#\{m\le N:\ p^2\mid a_m\ \text{for some }p\mid r\}\bigr), \]
where $b^{\pm}_k=\#\{m\le N:\mu^2(m)=1,\ k_m=k,\ \pm a_m>0\}$ depend on $N$ only. For
$r=pq$ with $p,q\asymp N^{1/2}$ one would want the error term to be $O(1)$, which would need Lemma~\textup{\ref{lem:localsq}} in the form $\#\{m\le N:p^2\mid a_m\}\ll N/p^2$. That form is
the one that lemma explicitly disowns, since its implied constant depends on $p$ while $p$ here grows
with $N$; the error term is not controlled at this range.

\emph{Regime caveat, and its retraction.} That $O(1)$ holds only at $p\asymp N^{1/2}$; at
$P\asymp(\log N)^{1/4}$ the same estimate gives only $N/p^2\asymp N/(\log N)^{1/2}$. Two earlier
revisions read that as a second, independent reason for the \textsc{conjecture} label then carried by
Theorem~\ref{cor:a9rate}. It is not, for two reasons. The comparison there was made against $N$,
whereas the benchmark is that statement's own target, now
$N(\log\log N)^{1/2+\delta}(\log N)^{-1/4}$ \textup{(}Theorem~\ref{cor:a9rate}\textup{)}, against which
$N(\log N)^{-1/2}$ is \emph{smaller}, an error term below the claimed rate cannot obstruct it.
And Theorem~\ref{cor:a9rate} does not travel this route at all: re--derived as in
Proposition~\ref{prop:condrate} it needs no kernel reversal, and Proposition~\ref{prop:circular}
shows the reversal is no reduction in any case. The operative blocker is recorded at the corollary
itself.
\end{lemma}
\begin{proof}
If $\gcd(a_m,r)=1$ then $(a_m/r)=(\pm1/r)(k_m/r)(s_m/r)^2=(\pm1/r)(k_m/r)$. The two sides can differ
only when some $p\mid r$ divides $s_m$ but not $k_m$, that is when $p^2\mid a_m$, which is the error
term. The coefficients $b^{\pm}_k$ carry no dependence on $r$.
\end{proof}
Measured at $N=5\times10^6$: $30.0\%$ of the $a_m$ are non-squarefree, but only $0.035\%$ have
$s_m>10^3$, so for $p\asymp N^{1/2}$ the error term is empty in practice
\textup{(\texttt{code/sqsieve\_holes.rs})}. The kernels $k_m$ are even $22.2\%$ of the time, so
the moduli are split into the four classes mod $8$ before the sieve is applied, at a cost of a
factor $4$.

\begin{hypothesis}[Kernel collisions]\label{hyp:kernel}
$\sum_k\bigl(|b^{+}_k|^2+|b^{-}_k|^2\bigr)\ll N^{1+o(1)}$, equivalently
$\#\{(m,m')\ \text{squarefree},\ \le N:\ a_ma_{m'}\ \text{is a perfect square}\}\ll N^{1+o(1)}$.
\end{hypothesis}

\begin{proposition}[The route is circular]\label{prop:circular}
Hypothesis~\textup{\ref{hyp:kernel}} implies the bound it would be used to prove, so it cannot serve
as a reduction. Indeed $b^{+}_1=\#\{m\le N:\ \mu^2(m)=1,\ a_m\ \text{is a positive perfect square}\}$ is precisely
the minus layer, and $(b^{+}_1)^2\le\sum_k|b^{+}_k|^2$, so
Hypothesis~\textup{\ref{hyp:kernel}} forces $b^{+}_1\ll N^{1/2+o(1)}$, which is the conclusion.
\end{proposition}

\emph{Why computation cannot detect this.} The circularity is invisible to measurement. Below the barrier $\sigma(m)>2$ is impossible, so
$a_m>0$ never occurs and $b^{+}_1=0$ identically: at $N=5\times10^6$ one finds $b^{+}_1=0$ while
$b^{-}_1=1213$, the latter contributing $12.5\%$ of $\sum_k|b_k|^2$
\textup{(\texttt{code/sqsieve\_kernel.rs})}. Every computed check of
Hypothesis~\ref{hyp:kernel} therefore tests a statement from which the load-bearing term has been
deleted. This is the same blindness recorded in Proposition~\ref{B-prop:nearmiss}, appearing here in a
form that would otherwise have passed as supporting evidence.

Accordingly, the large sieve does apply and Lemma~\ref{lem:reversal} is correct, but the resulting
statement is a \emph{reformulation}: $\mathrm{CRUX}\text{-}\mathrm{A}9$ is equivalent to a kernel-collision
bound, with the equivalence carried by the $k=1$ term. The new attack surface is genuine, since the
reformulated statement is unsigned and about products, but it is not a reduction and is not counted
as progress on the crux. The unconditional record is Proposition~\ref{prop:condrate}.


\emph{The hypothesis measures, and its constant identifies itself.} $P(D,r)$ against $D^2/r+D$ has ratio $0.3688$, $0.3688$, $0.3689$ at $r=1009$;
$0.3705$, $0.3696$, $0.3695$ at $r=10007$; and $0.3867$, $0.3727$, $0.3702$ at $r=100003$, for
$D=10^6,5\times10^6,2\times10^7$ \textup{(\texttt{code/sqsieve\_bilinear.rs})}. The ratio is not
merely bounded but constant, and the constant is $(6/\pi^2)^2=0.3696$, the square of the squarefree
density, which is exactly what equidistribution of $\sigma_r$ predicts once both $d$ and $d'$ are
restricted to squarefree integers.

Compare Hypothesis~\ref{hyp:sq}. That one asks for signed square-root
cancellation, and Proposition~\ref{prop:routeclosed} showed the only available route to it fails.
Theorem~\ref{thm:pairavg} instead bounds an \emph{unsigned} second moment, and does so
unconditionally; what it does not do is bound $T_r$, since the passage from bilinear forms to $T_r$
is exactly the gap recorded in Remark~\ref{rem:missingdecomp}.


\emph{What made this work.} Every previous attack on $\mathrm{A}9$ counted members of the minus layer directly, and each ran
into the same obstruction: the layer is a union of lines $m'=2m-k$, each thin by
Proposition~\ref{prop:lines} but not summably so, and the peeling reorganisations that would combine
them are all conjugate under the groupoid of Proposition~\ref{prop:groupoid}. The square sieve never
counts a member. It asks only whether $m'-2m$ is a square, and answers by averaging a Jacobi symbol,
so the groupoid has nothing to act on. The one obstacle, that $m'$ is not a polynomial, is removed by
Lemma~\ref{lem:symbolfact}: the symbol splits into a character in $m$ times a character in
$\sigma_r(m)$, and $\sigma_r$ is additive, which is exactly the hypothesis Hal\'asz needs. The
analytic inputs are all standard: Hal\'asz's theorem, the Gauss sum evaluation
$|\tau(\chi)|=\sqrt r$, the classical zero--free region, the $\log L$ bound for non--principal
characters, and Siegel--Walfisz. No Sali\'e or Weil bound is used: the sum in question is a Gauss
sum, as Lemma~\ref{lem:charcancel} records. Note also that the theorem does not decide \#307, and cannot: it is a statement about
density, and a single two-cycle is a density-zero event.


\begin{proposition}[Slot stratification of the divisibility set]\label{prop:slotstrat}
Fix $c$ and let $D(c)=\{m \text{ squarefree}: (2m-m')\mid(m-c),\ m\neq c\}$. For $m\in D(c)$ put
$p=(m-c)/(2m-m')$, the \emph{slot}. Then
\[ p\,m' \;=\; (2p-1)\,m + c, \qquad\text{equivalently}\qquad \sigma(m)=2-\frac{m-c}{p\,m}, \]
so $D(c)=\bigsqcup_{p\neq0} L(p,2p-1,c)$ is a disjoint union of generalized lines in the sense of
Proposition~\textup{\ref{prop:groupoid}}, the stratum $p=1$ being the slope-\emph{one} line
$m'=m+c$. The identity was verified with $0$ violations on all $2{,}294$ members with
$m\le2\times10^7$, $|c|\le400$ \textup{(\texttt{code/slot\_stratify.rs})}.
\end{proposition}
\begin{proof}
$m-c=p(2m-m')$ rearranges to $pm'=(2p-1)m+c$; dividing by $pm$ gives the mass form. Distinct $p$
give disjoint strata since $p$ is a function of $m$.
\end{proof}

\begin{theorem}[The residual difficulty sits entirely at the mass-two wall]\label{thm:masstwo}
For every $\varepsilon>0$ and every $c$,
\[ \#\bigl(D(c)\cap[1,Z]\bigr)\;\le\;\bigl(\delta(2-\varepsilon)+o_\varepsilon(1)\bigr)\,Z+O(|c|), \]
where $\delta(x)$ is the density of $\{m:\sigma(m)\ge x\}$. Moreover
\[ \delta(x)\;\le\;\inf_{t>0}\;e^{-tx}\prod_{p}\Bigl(1+\frac{e^{t/p}-1}{p}\Bigr), \]
which evaluates to $\delta(1)\le10^{-0.8}$, $\delta(1.5)\le10^{-7.2}$ and
$\delta(2)\le10^{-105.6}$ \textup{(\texttt{code/chernoff\_mass.rs})}.
\end{theorem}
\begin{proof}
Split $D(c)$ at $|p|\le P$ and $|p|>P$. The first part is a union of at most $2P$ generalized lines,
each of size $\ll Z/(\log Z)^{(1+o(1))\log_3 Z}$ by Proposition~\ref{prop:lines}
\textup{(}\emph{caveat}: Proposition~\ref{prop:lines} is stated for \emph{integer} slope $e\ge1$,
while the strata $L(p,2p-1,c)$ of Proposition~\ref{prop:slotstrat} have slope $(2p-1)/p$, which is
non--integral for $p\ne1$. The bound is therefore invoked outside the range in which it is proved,
and the first part of this proof is \textsc{conjecture} pending a version of
Proposition~\ref{prop:lines} uniform over rational slopes\textup{)}; taking
$P=P(Z)\to\infty$ slowly makes the total $o(Z)$. For the second, $|p|>P$ forces
$|2-\sigma(m)|<(1+|c|/m)/P$ by the mass form, so for $m>|c|$ we get $\sigma(m)>2-2/P$, and the count
is at most $\delta(2-2/P)Z(1+o(1))$. For the Chernoff bound, $f(m)=\prod_{p\mid m}e^{t/p}$ is
multiplicative with $f(p)\ge1$, so
$\sum_{m\le Z}f(m)=\sum_{d\le Z}\mu^2(d)\prod_{p\mid d}(f(p)-1)\lfloor Z/d\rfloor
\le Z\prod_p(1+(f(p)-1)/p)$, and $\#\{\sigma\ge x\}\le e^{-tx}\sum_{m\le Z}f(m)$.
The product converges because $(e^{t/p}-1)/p\sim t/p^2$.
\end{proof}

\begin{remark}[What this settles and what it does not]\label{rem:masstwo}
Theorem~\ref{thm:masstwo} replaces the constant $\delta=0.0420$ of
Proposition~\ref{prop:divset}\,(ii) by one below $10^{-105}$, an unconditional gain of more than a
hundred orders of magnitude, and it does so by \emph{using} the divisibility rather than merely the
mass condition it forces. It also localises the residue exactly: away from the mass-two wall the set
is a bounded union of thin lines, and \emph{all} the remaining difficulty lies in
$\{\sigma(m)\ge2-\varepsilon\}$. It does not prove $\mathrm{A}9$. The set $\{\sigma\ge2\}$ has
\emph{positive} density, at least $1/\Pi_{59}=10^{-112.9}$ since every multiple of $\Pi_{59}$ lies
in it, so the method cannot reach $o(Z)$: the limiting constant is positive, not zero. The bound
$10^{-105.6}$ sits within seven orders of that truth, so little is lost in the Chernoff step, and the
obstruction is genuine rather than an artefact of the estimate. The residual question is therefore
sharp: is $D(c)$ thin \emph{within} the mass-two wall, the one place where the barrier
\textup{(Theorem~\ref{K-thm:barrier})} also lives?
\end{remark}

\begin{proposition}[The divisibility set: slope one is dense, slope two is not]\label{prop:divset}
Let $e\ge1$ and $c$ be integers and let $m$ be squarefree with
$(em-m')\mid(m-c)$ and $m\neq c$. Then
\[ \sigma(m)\;\ge\;e-1-\frac{|c|}{m}. \]
Consequently the two slopes behave oppositely.
\begin{enumerate}[label=\textup{(\roman*)},leftmargin=2.2em]
\item For $e=1$ the condition is vacuous and the set has positive density: every prime $m$ satisfies
it when $c=1$, since then $m'=1$ and $em-m'=m-1=m-c$. Over $m\le2\times10^7$ this accounts for
$1{,}270{,}651$ of the members, exactly $\pi(2\times10^7)$.
\item For $e=2$ the condition reads $\sigma(m)\ge1-|c|/m$, so by Proposition~\ref{A-prop:density}
\[ \#\{m\le Z:\ (2m-m')\mid(m-c),\ m\neq c\}\;\le\;(\delta+o(1))\,Z+O(|c|), \]
with $\delta=0.0420\ldots$ the abundance density, certified in Proposition~\ref{A-prop:density} to lie
in $[0.04202,0.04223]$.
\end{enumerate}
\end{proposition}
\begin{proof}
If $|em-m'|>|m-c|$ the divisibility forces $m=c$. So $|em-m'|\le|m-c|\le m+|c|$; below the barrier
$\sigma(m)<e$ for $e\ge2$, so $em-m'=m(e-\sigma(m))>0$ and $m(e-\sigma(m))\le m+|c|$, which is the
display. Part (i) is the computation just given, and (ii) is Proposition~\ref{A-prop:density} applied
to $\{\sigma\ge1-\varepsilon\}$. The mass condition was checked with $0$ violations over all
squarefree $m\le2\times10^7$ and all $|c|\le400$ at $e=1,2,3$, where the nontrivial member counts are
$2{,}197{,}230$, $28$ and $8$ respectively (\texttt{code/a9prime\_falsify.rs}).
\end{proof}

\emph{Consequence for the residual question.} Proposition~\ref{prop:divset} settles the shape of the residual bound. A uniform statement over all
slopes is false, and false for a trivial reason: at $e=1$ the primes sit on the set by an identity.
The correct target is confined to $e\ge2$, where a genuine mass condition appears, and there
Proposition~\ref{prop:divset}(ii) gives the first unconditional bound, $\delta Z$ with
$\delta<0.043$. The gap to the $Z^{1/2}$ needed for Remark~\ref{rem:minusslot} is therefore now a
gap between two \emph{proved} quantities rather than between a proof and a wish, and the observed
counts at $e=2$ ($28$ over the whole range $|c|\le400$, against $Z^{1/2}=4472$) sit far below both.


\begin{remark}[The two open routes are one question, asked in opposite directions]\label{rem:onequestion}
The existence route of Remark~\ref{A-rem:lehmer} and the counting route of
Remark~\ref{rem:minusslot} look unrelated (one is a search for a single solution, the other a
density estimate) but they are the same question about the same kind of object. In each case an
index determines a single integer, and everything turns on whether that integer is prime: in
Remark~\ref{A-rem:lehmer} the Pell index $n$ determines $q_n=(x_n^2-D_S)/A_S$, a term of a ternary
linear recurrence, and \#307 needs it prime \emph{once}; in Remark~\ref{rem:minusslot} the cofactor
$m$ determines the slot $(c-am)/(m'-em)$ of Proposition~\ref{prop:groupoid}, and the density needs
it prime \emph{rarely}. Existence asks for a lower bound on the number of prime values, density for
an upper bound, on determined values along a thin deterministic sequence, and neither estimate is
available: no unconditional infinitude of prime terms is known for any non--degenerate linear
recurrence of order $\ge2$, and no worst--case upper bound is known for single--class prime counts
summed over all moduli. \#307 sits exactly between two missing estimates on one object, which is why
the heuristics can be read either way and why progress on either side would be informative about the
other. The remaining walls are, in this sense, one wall seen from two directions.
\end{remark}

\subsection{Slot structure}\label{ssec:slotstructure}

\begin{corollary}[Quadratic--residue filter for \#307 solutions]\label{cor:qr}
Let $(P,Q)$ solve Erd\H{o}s~\#307 and set $a=D_P$, $b=D_Q$, $N=ab$. Since $N'=a^2+b^2$ and
hence $N'+2N=(a+b)^2$ automatically, every prime $\ell\in P\cup Q$ satisfies $(N/\ell\mid\ell)=+1$:
the cofactor $N/\ell$ is a nonzero quadratic residue modulo $\ell$. Indeed $\ell$ divides exactly
one of $a,b$ (by $P\cap Q=\varnothing$), so $\ell\nmid a+b$ and $(a+b)^2\not\equiv0\pmod\ell$;
combined with the identity $N'+2N\equiv N/\ell\pmod\ell$ of~\textup{(v)}, this forces $N/\ell$ to
be a nonzero square mod $\ell$. As a search filter the condition prunes candidate unions by a factor
${\le}2^{-|P\cup Q|}\le2^{-59}$ (heuristically each of the $\ge59$ primes imposes an independent
QR constraint); for genuine cycles it is automatic and does not alter the existence heuristics.
We note that Corollary~\ref{cor:qr} is \emph{not} independent of Bado: it is the squared form of his
congruence $D_Q\equiv D_P/\ell\pmod\ell$ (itself $N_P$ reduced mod $\ell$ by cofactor survival),
giving $N/\ell\equiv D_Q^2\pmod\ell$. Squaring is lossy ($r$ and $-r$ share a square) so Bado's
linear congruence is strictly stronger, pinning $D_Q\bmod\ell$ exactly where the residue filter
determines it only up to sign; the filter's one added value is being \emph{split--free} (the residue
$N/\ell\bmod\ell$ is computable from the candidate union alone, with no knowledge of the $P\mid Q$
partition). The general law of Proposition~\ref{prop:plus}(v), by contrast, holds for \emph{any}
squarefree plus--hit with no $P\mid Q$ split and is not subsumed by Bado.
\end{corollary}

The plus-- and minus--hits both organise into one--prime ``slot'' families, which makes the structure
of Proposition~\ref{prop:splitdisc} explicit at every $\omega$.

\begin{proposition}[Slot calculus for the discriminant layers]\label{prop:slot}
Let $N_0$ be squarefree, $r\nmid N_0$ a prime, $N=N_0r$, and $E_0\coloneqq N_0'+2N_0$. By Leibniz
($N'=N_0'r+N_0$),
\[ N'+2N=rE_0+N_0,\qquad N'-2N=r(N_0'-2N_0)+N_0, \]
so $N$ is a plus--hit iff $r=(s^2-N_0)/E_0$ for some integer $s$ (equivalently $s^2\equiv N_0\pmod{E_0}$).
\begin{enumerate}[label=\textup{(\alph*)},leftmargin=2.2em]
\item \emph{Every $\omega$--stratum.} A plus--hit $N$ ($\omega\ge2$) factors, through any prime divisor
$r$, as $N=N_0r$ with $N_0=N/r$ squarefree, and lies in the quadratic family
$\{N_0(s^2-N_0)/E_0\}_s$ of its base; generalising the $p=5$ family of
Proposition~\ref{prop:plus}(iii) beyond semiprimes (each step adds one prime,
$\omega(N)=\omega(N_0)+1$). The congruence $s^2\equiv N_0\pmod{E_0}$ is solvable for $669$ of the
$3041$ squarefree $N_0\in[2,5000]$ ($22.0\%$); solvability is necessary but not sufficient for the
family to contain a hit.
\item \emph{Recursive towers: a trichotomy.} When the base is itself a plus--hit, $E_0=s_0^2$ is a
square; below $2\times10^5$ the $114$ plus--hit bases split into three kinds. \emph{No residue} ($84$):
$s^2\equiv N_0\pmod{E_0}$ is unsolvable and the family is empty; a quadratic--residue obstruction,
occurring for both even $E_0$ ($N_0=305$, with $305$ a non--residue mod $13$) and odd $E_0$
($N_0=1870$, $E_0=73^2$). \emph{Fixed divisor} ($13$, each with fixed divisor $2$): the family is
nonempty but every slot is even, so only $r=2$ could serve and never gives a square; sterile (e.g.\
$N_0=145$, $E_0=18^2$). \emph{Live} ($17$): fixed divisor $1$ and a prime slot is exhibited (e.g.\
$N_0=130$, first prime slot $r=83399$, giving the plus--hit $N=10{,}841{,}870$ with $N'+2N=5487^2$).
\item \emph{Minus layer; the barrier in slot form.} For $\sigma(N_0)<2$ the minus identity gives
$N'-2N\ge0$ iff $r\le1/(2-\sigma(N_0))$, i.e.\ iff $\sigma(N)\ge2$: this is the slot--variable form of
Corollary~\ref{K-cor:emptytest}, not an independent proof of it. The bound $1/(2-\sigma(N_0))$ is
unbounded as $\sigma(N_0)\to2^-$; over squarefree $N_0<10^5$ its maximum is only $1.524$ (at
$N_0=30030$), so no such base admits a prime minus--slot, yet for the $58$--prime primorial the bound
exceeds $793$ and the next prime $277$ fits, producing the $59$--prime primorial, the smallest
minus--hit base. A prime minus--slot thus first becomes feasible \emph{exactly} at the $59$--prime
barrier; the minus--empty--below--barrier statement rests on the primorial growth of
Corollary~\ref{K-cor:emptytest}, not on the figure $1.524$.
\end{enumerate}
\end{proposition}
\noindent (The two identities are checked symbolically and on $2\times10^4$ random pairs; the base
trichotomy, slot lists, and the $1.524$ maximum by exhaustive enumeration.)

The two square conditions over a base combine into a single binary--form representation, which places
the obstruction inside Gauss's genus theory.

\begin{proposition}[The Pythagorean condition as a binary--form representation]\label{prop:form}
For squarefree $N_0$ and prime $r\nmid N_0$, $N=N_0r$, the identity
\[ (2N_0-N_0')(N'+2N)+(2N_0+N_0')(N'-2N)=4N_0^2 \]
holds unconditionally (Leibniz; checked on $2\times10^4$ random pairs). Hence if $N$ is a Pythagorean
pair product over $N_0$ (both $N'+2N=s^2$ and $N'-2N=d^2$ squares) then $(s,d)$ represents
$4N_0^2$ by the binary form $Q_{N_0}(s,d)=A\,s^2+B\,d^2$ with $A=2N_0-N_0'=N_0(2-\sigma(N_0))$,
$B=2N_0+N_0'=N_0(2+\sigma(N_0))$, of discriminant $-4AB=4(N_0'^2-4N_0^2)$; conversely such a
representation yields a Pythagorean pair over $N_0$ exactly when the recovered slot $r=(s^2-N_0)/E_0$
is a positive integer, prime, and coprime to $N_0$ (the integrality $s^2\equiv N_0\pmod{E_0}$ holds on
every representation we found, $N_0<2\times10^4$). For $\sigma(N_0)<2$ the form is positive--definite
($A>0$), and below the barrier representations are rare and all ``ghosts'': the only squarefree
$N_0\le3000$ representing $4N_0^2$ are $N_0=30$ (with $(s,d)=(11,1)$) and $N_0=1722$ ($(83,1)$), each
carrying the non--admissible slot $r=1$, and no live prime slot occurs. As $\sigma(N_0)\to2^-$ the
leading coefficient $A=N_0(2-\sigma(N_0))\to0$ and the ellipse $A\,s^2+B\,d^2=4N_0^2$ degenerates to an
arbitrarily thin sliver: the geometric face of the mass--2 barrier. Consequently \#307 is the existence
of a base with $\sigma(N_0)\ge2-1/r$ whose form $Q_{N_0}$ represents $4N_0^2$ with an admissible,
integral, \emph{prime} slot and deviation $k=0$.
\end{proposition}
\noindent (We make no class--number claim: each base yields a form of a \emph{different} discriminant,
and the question is single--value representation, not a class count; only the degeneration $A\to0$ is
used.)

\begin{proposition}[An explicit Pythagorean family, and why it stops at rung $-1$]\label{prop:family}
Let $N_0$ be squarefree with $N_0'=2N_0-1$ \textup{(}the slope--two line $c=-1$, i.e.\ $k_0=1$\textup{)}
and suppose $N_0-1$ is prime. Then $(a,b)=(N_0-1,\,N_0)$ is a Pythagorean pair:
\[ a^2+b^2\;=\;2N_0^2-2N_0+1\;=\;(ab)'. \]
It arises from the form equation of Proposition~\ref{prop:form} at its degenerate end: with $k_0=1$,
$E_0=4N_0-1$, the pair $(s,d)=(2N_0-1,\,1)$ represents $4N_0^2$ exactly, and the slot it yields is
$r=(s^2-N_0)/E_0=N_0-1$. Every member of the family sits on rung $k=-1$, and no member can sit on
rung $0$: $a$ is prime, so $a'=1$, and $a'=b$ would force $N_0=1$.
\end{proposition}
\begin{proof}
$(ab)'=N_0'(N_0-1)+N_0(N_0-1)'=(2N_0-1)(N_0-1)+N_0=2N_0^2-2N_0+1=(N_0-1)^2+N_0^2$, using
$(N_0-1)'=1$ for $N_0-1$ prime. The form identity is
$(2N_0-1)^2+(4N_0-1)=4N_0^2$, and $(s^2-N_0)/E_0=(4N_0^2-5N_0+1)/(4N_0-1)=N_0-1$. The rung is
$k=(a'-b)/a=(1-N_0)/(N_0-1)=-1$.
\end{proof}

\emph{What the family shows.} This is the first explicit construction of Pythagorean pairs in this note: everything earlier
established that they are rare \textup{(}Corollary~\ref{K-cor:emptytest}: none below $10^{112}$\textup{)}
without exhibiting a mechanism. The mechanism is a clean reduction (the pair exists as soon as the
line $n'=2n-1$ carries a member whose predecessor is prime) and it is Bunyakovsky--shaped, matching
the observation in Remark~\ref{A-rem:lehmer} that a single square condition is of that class. It is also
conditional in a way the barrier controls: $\sigma(N_0)=2-1/N_0$ forces $\omega(N_0)\ge58$, so the
family is empty below the scale where the critical set of Remark~\ref{rem:minusslot} switches on.

Two further consequences. First, combining with Corollary~\ref{cor:kdetermined}: for odd $N_0$ the
pair has mixed parity with odd member $N_0$, so $k\equiv\omega(N_0)\pmod2$; since $k=-1$ here,
\emph{any odd member of the line $n'=2n-1$ whose predecessor is prime has $\omega(N_0)$ odd}; a
necessary condition on that line, derived rather than observed. Second, and this is the honest
limitation, the construction is \emph{permanently} off the cycle rung: primality of $a$ is exactly
what makes it work and exactly what forces $a'=1\ne b$. A $k=0$ analogue would need $a$ composite with
$a'=b$, which is the original problem. So the family exhibits the Pythagorean locus without
approaching \#307; useful as the first worked mechanism, not as a route.


\begin{proposition}[The composite mechanism: $|P|=2$ is one line and one primality]\label{prop:compositea}
Write $a=2m$ with $m$ odd squarefree, so that $b=a'=m+2m'$ \textup{(}Theorem~\ref{K-thm:parity}\textup{)}.
The two--cycle condition is then the single equation
\[ b'\;=\;2m,\qquad b=m+2m', \]
in the one unknown $m$. For $m=q$ prime it reads $b=q+2$ and
\[ \boxed{\ b'=2b-4\ }\qquad\text{with } b-2 \text{ prime}, \]
and any such $b$ \emph{is} a solution: $a=2q$, $b=q+2$ satisfy $a'=b$, $b'=a$ and
$\sigma(a)\sigma(b)=\bigl(\tfrac12+\tfrac1q\bigr)\bigl(2-\tfrac4{q+2}\bigr)=1$ identically. Hence
\[ \#307\ \text{with}\ |P|=2 \iff \exists\,b \text{ squarefree}:\ b'=2b-4 \text{ and } b-2 \text{ prime}. \]
\end{proposition}
\begin{proof}
$(2m)'=m+2m'$ by Leibniz, so $a'=b$ by definition of $b$, and the cycle needs only $b'=a=2m$. For
$m=q$ prime, $b=q+2$ and $b'=2q=2(b-2)=2b-4$; the mass identity is the displayed product.
\end{proof}

\emph{Two lines, two rungs.} Comparing with Proposition~\ref{prop:family} isolates exactly what separates the cycle from its
neighbour. Both are ``a member of a slope--two line whose shifted predecessor is prime'':
\[
\begin{array}{lll}
\text{line } n'=2n-1, & N_0-1 \text{ prime} & \Longrightarrow\ \text{Pythagorean pair, rung } k=-1;\\
\text{line } n'=2n-4, & b-2 \text{ prime}   & \Longrightarrow\ \text{genuine two--cycle, rung } k=0.
\end{array}
\]
The difference between the open problem and its solved neighbour is the constant in the line. And the
obstruction changes character with it: Proposition~\ref{prop:family} is barred from rung $0$
\emph{structurally} (its $a$ is prime, so $a'=1\ne b$) whereas Proposition~\ref{prop:compositea}
is barred only by \emph{scale}, since $\sigma(b)=2-4/b$ forces $\omega(b)\ge58$ and $b\ge\Pi_{58}$,
putting the whole family above $10^{108}$ and the resulting $N=2qb$ at the barrier. A direct search of
the reduction over odd squarefree $m\le2\times10^6$ ($810{,}590$ admissible $m$, of which $111{,}464$
prime) finds none, with nearest miss $|b'-2m|=5$ at $m=3$ (\texttt{code/composite\_a.rs}); exactly
as Corollary~\ref{K-cor:emptytest} demands. This is the composite mechanism the rung--$0$ case requires;
what it does not do is evade the barrier, which reappears as the mass of the line.


\begin{proposition}[The barrier is a function of the split]\label{prop:sectorbarrier}
For a solution with $|P|=j$ write $a=\prod P$, $b=\prod Q$. Since $\sigma(a)\sigma(b)=1$ and
$P\cap Q=\varnothing$, fixing $j$ forces $\sigma(b)=1/\sigma(a)$ and leaves $Q$ only the odd primes
outside $P$, which bounds $N$ from below split by split. Taking $a=2\cdot(\text{the }j-1\text{
smallest odd primes})$, the most favourable configuration at that size:
\[
\begin{array}{r|cccc}
|P| & \sigma(a) & \sigma(b) & |Q|\ge & \log_{10}N\ge\\\hline
2 & 0.8333 & 1.2000 & 67 & 137.9\\
3 & 1.0333 & 0.9677 & 56 & 112.9\\
4 & 1.1762 & 0.8502 & 63 & 132.8\\
5 & 1.2671 & 0.7892 & 73 & 161.0\\
6 & 1.3440 & 0.7440 & 85 & 195.3\\
12 & 1.5927 & 0.6279 & 181 & 491.1
\end{array}
\]
The minimum over splits is $10^{112.9}$ at $|P|=3$, recovering the constant of
Theorem~\ref{K-thm:barrier} by an independent route, and every other split is strictly dearer.
\end{proposition}
\begin{proof}
$\sigma(b)=1/\sigma(a)$ is the cross--match. $Q$ is disjoint from $P$, so its mass must be assembled
from odd primes outside $P$; the least number of them reaching $\sigma(b)$, and the product of that
many, are finite computations (\texttt{code/sector\_barrier.rs}). Enlarging any prime of $a$ lowers
$\sigma(a)$, raises $\sigma(b)$ and hence $|Q|$, so the displayed configuration is extremal for each
$j$.
\end{proof}

\emph{Why the composite mechanism is out of reach.} Proposition~\ref{prop:sectorbarrier} settles the standing of the clean reduction of
Proposition~\ref{prop:compositea}. There $|P|=2$, and the mechanism is forced through the dearest
small sector: $b$ is odd, so $Q$ must assemble mass $\sigma(b)=2q/(q+2)$ from odd primes alone, and
the sector cannot begin below $10^{137.9}$; some twenty--five orders of magnitude \emph{above} the
general barrier $10^{112.9}$. The line $n'=2n-4$ is therefore a correct and complete description of
the $|P|=2$ case and simultaneously the most expensive place to look.

This also explains the balance the barrier strikes. The cheapest configuration is the near--balanced
one, $\sigma(a)\approx\sigma(b)\approx1$ at $|P|=3$; pushing the split either way (towards a small
$P$ with $\sigma(b)\to2$, or towards a large $P$ with $\sigma(b)\to0$) costs orders rapidly, because
mass near $2$ needs the odd primes to $\approx10^4$ and mass near $0$ needs many large ones. The
$59$--prime barrier is thus not merely a statement about $|P\cup Q|$ but the \emph{minimum of a
convex profile over splits}, and Theorem~\ref{K-thm:barrier}'s constant is the value at that minimum.
(Rows with $|P|+|Q|<60$ are further excluded by Proposition~\ref{K-prop:close59}, which sharpens the
$|P|=3$ entry without moving the minimum's location.)

\begin{proposition}[The cycle condition does not improve the barrier]\label{prop:nogain}
A two--cycle satisfies more than the mass inequality: it forces $b=a'$, that is $b=a\,\sigma(a)$
exactly. Since $\sigma(b)=1/\sigma(a)$ and $P\cap Q=\varnothing$, writing $S$ for the small primes of
$a$ and $\Pi_{\min}(S)$ for the least product of primes outside $S$ with reciprocal sum
$\ge1/\sigma(S)$, one gets
\[ a\;\ge\;\max\Bigl(\textstyle\prod S,\ \Pi_{\min}(S)/\sigma(S)\Bigr),\qquad
   N\;=\;a^2\sigma(a)\;\ge\;\max\Bigl(\textstyle\prod S,\ \Pi_{\min}(S)/\sigma(S)\Bigr)^2\sigma(S). \]
Minimised over $S$, this yields $\log_{10}N\ge113.2$, attained at
$S=\{\text{odd }p\le151\}\setminus\{3,97\}$ with $\sigma(S)=1.0495$ and $|Q|\ge26$. The barrier of
Theorem~\ref{K-thm:barrier} is $\log_{10}N\ge112.9$: \emph{the cycle condition buys essentially nothing
beyond the mass condition}, and Theorem~\ref{K-thm:barrier} is already close to optimal for arguments
of this class.

\textup{Formalised.} The deductive content is \texttt{Erdos307.cycle\_bound\_max}: two independent
lower bounds on $a$ combine through the maximum, and $N=a^2\sigma(a)$ is monotone in both arguments,
so the constraints give a bound on $N$. The conclusion drawn, that the gain over
Theorem~\ref{K-thm:barrier} is under one order, is an \emph{upper} bound on a minimum, and an upper
bound on a minimum needs one witness, not the search. Since $\Pi_{\min}(S)$ is itself a minimum,
bounding the expression at $S$ from above needs a witness for it too: any $T$ disjoint from $S$ with
$\sigma(T)\ge1/\sigma(S)$ has $\Pi_{\min}(S)\le\prod T$. Taking $S$ to be the odd primes to $151$
less $\{3,97\}$ and $T$ its greedy completion, $26$ primes ending at $277$, one checks
$\prod T\le\prod S\cdot\sigma(S)$, so the maximum is $\prod S$, and
\[ \Bigl(\bigl(\textstyle\prod S\bigr)^2\sigma(S)\Bigr)^{2}\;<\;10^{227}, \]
squared so the exponent is an integer: the bound on $\log_{10}N$ at this $S$ is at most $113.5$, and
the gain over $112.9$ is under $0.6$ of one order. Note $\bigl(\prod S\bigr)^2\sigma(S)$ is an
integer, $\sigma(S)$ having denominator exactly $\prod S$, so the whole check is exact integer
arithmetic. $0$ \texttt{sorry} \textup{(\texttt{lean/Erdos307/NoGain.lean},
\texttt{lean/Erdos307/NoGainWitness.lean})}. The exhaustive minimisation certifies the sharper claim
that $113.2$ is exactly the minimum; that is not needed here, and it was the only reason this
proposition was carried as a formalisation refusal.
\end{proposition}

\begin{remark}[Why restricted families overstated the gain]\label{rem:sectorsquared}
Compare the conclusion of Proposition~\ref{prop:nogain} with what restricted families
suggest, because the discrepancy is large and instructive. Confining $S$ to an initial segment
$\{2,3,5,\dots\}$ gives a minimum of $10^{222.9}$; enumerating all subsets of the first $14$, $16$,
$18$ primes gives $10^{195.9}$, $10^{188.9}$, $10^{181.7}$; still falling, hence not converged.
Each of these looks like a large improvement of the barrier and none of them is one. The optimum sits
outside all of them: it hands $2$ to $Q$ (half of $Q$'s mass budget in a single prime), keeps the odd
primes to $151$ in $S$, and drops two of them, a configuration no initial segment and no small subset
can express. Only after searching that structure does the bound settle at $113.2$.

The lesson is about the shape of the search rather than the arithmetic. A bound minimised over a
family that cannot express the optimum reports the family's best, not the problem's, and here that
error was worth eighty orders of magnitude. We record the negative result as the substantive one:
imposing $b=a'$, the strongest structural constraint available beyond the masses, does not move the
barrier, which is evidence that $10^{112.9}$ is not an artefact of the AM--GM route but the true
cost of the mass condition (\texttt{code/full\_minimisation.rs},
\texttt{code/minimise\_structured.rs}).
\end{remark}

\begin{remark}[Decoupling the two squares: what it buys, and what it costs]\label{rem:decouple}
Remark~\ref{A-rem:lehmer} locates the exponential difficulty in the \emph{conjunction} of the two
squares, so it is worth asking for a family on which one of them is free. The slot calculus supplies
one: on
\[ \mathcal F(N_0)\;=\;\bigl\{\,N_0r\ :\ r=(s^2-N_0)/E_0\,\bigr\} \]
the plus--square $N'+2N=s^2$ holds \emph{identically}, by construction. The minus--square then costs
exactly the form equation of Proposition~\ref{prop:form}, $k_0s^2+E_0d^2=4N_0^2$, and the class of
the residual question genuinely changes. On a coupled tail family the solutions run along a Pell
orbit and the candidate slots grow like $\varepsilon^n$; here, for $\sigma(N_0)<2$ the form is
positive definite, so $4N_0^2$ has only \emph{finitely} many representations and each base carries
finitely many candidate slots. Decoupling therefore trades a Lehmer sequence for a finite
per--base check.

What it does not do is help. Over the $30{,}397$ squarefree bases $N_0\le5\times10^4$ with
$\sigma(N_0)<2$ the form represents $4N_0^2$ exactly \emph{twice} (at most once per base, largest
$s=83$) and both representations are the ghost slot $r=1$; there is no live prime slot at all
(\texttt{code/plus\_automatic.rs}, reproducing the sub--barrier census of
Proposition~\ref{prop:form} from an independent enumeration). The reason is the degeneration already
noted: the barrier forces $\sigma(N_0)\uparrow2$, so $k_0=N_0(2-\sigma)$ collapses and the form
becomes ever more lopsided, while the target stays at $4N_0^2$. So freeing one square does not free
the problem; it relocates the difficulty from the primality of an exponential sequence to the
solubility of a degenerating quadratic form. The pattern is the circularity of
Section~\ref{K-sec:frame} of the core note once more, but the relocation is informative, since a finite, degenerating
representation problem is a different kind of object from a Lehmer sequence, and no no--go here
forbids attacking it.
\end{remark}

The emptiness of that census is structural, and the reason fixes the range in which the question
lives.

\begin{proposition}[The live--slot threshold]\label{prop:liveslot}
Let $N_0$ be squarefree and suppose $Q_{N_0}(s,d)=A\,s^2+B\,d^2$ represents $4N_0^2$ with a slot
$r=(s^2-N_0)/B$ satisfying $r\ge2$. Then
\[ \sigma(N_0)\;\ge\;\tfrac32, \]
hence $\omega(N_0)\ge10$ and $N_0\ge\Pi_{10}=6{,}469{,}693{,}230$.
\end{proposition}
\begin{proof}
Assume first $\sigma=\sigma(N_0)<2$, so $A=N_0(2-\sigma)>0$ and the form is positive definite. From
$r\ge2$ and $B=N_0(2+\sigma)$,
\[ s^2\;\ge\;N_0+2B\;=\;N_0(5+2\sigma), \]
while $A\,s^2\le A\,s^2+B\,d^2=4N_0^2$ gives $s^2\le 4N_0^2/A=4N_0/(2-\sigma)$. The two are
compatible only if $(5+2\sigma)(2-\sigma)\le4$, that is $2\sigma^2+\sigma-6\ge0$, whose positive root
is exactly $\tfrac32$. Since among $k$--element prime sets the reciprocal sum is maximised by the
first $k$ primes (Lemma~\ref{K-lem:59}) and
$\sigma(\Pi_9)=1.498956\ldots<\tfrac32\le1.533439\ldots=\sigma(\Pi_{10})$, the mass bound forces
$\omega(N_0)\ge10$ and $N_0\ge\Pi_{10}$. If instead $\sigma(N_0)\ge2$ then $N_0\ge\Pi_{59}>\Pi_{10}$
by Theorem~\ref{K-thm:barrier}, so the conclusion holds a fortiori.
\end{proof}

The threshold is formalised. \texttt{Erdos307.live\_slot\_threshold} proves the analytic core, that
a representation of $4N_0^2$ with $s^2\ge N_0+2B$ forces $\sigma\ge\tfrac32$, by dividing out $N_0^2$
and factoring $4-(2-\sigma)(5+2\sigma)=(2\sigma-3)(\sigma+2)$ against $\sigma+2>0$; and
\texttt{threshold\_bracket} checks the two numerals
$\sigma(\Pi_9)<\tfrac32\le\sigma(\Pi_{10})$ that place the entry point at ten primes. Both carry $0$
\texttt{sorry} and the three standard axioms \textup{(\texttt{lean/Erdos307/LiveSlot.lean})}. The
passage from the mass bound to $\omega(N_0)\ge10$ is the step already formalised as
\texttt{card\_ge\_59\_of\_recipSum\_ge\_two} for the barrier, and is cited rather than repeated.

The census reported above ran to $N_0\le5\times10^4$, which is below $\Pi_{10}$ by a factor
$1.3\times10^5$: it could not have met a live slot, and its emptiness is therefore not evidence
about the family. The two representations it did find, at $N_0=30$ and $N_0=1722$, have
$\sigma=1.0333$ and $\sigma=1.0006$, both far below $\tfrac32$, which is why both are ghosts.

Three mass thresholds now organise the problem, and each is a Mertens count: $\sigma=1$ at three
primes, $\sigma=\tfrac32$ at ten, and $\sigma=2$ at fifty--nine. The first bounds the trivial regime,
the third is Theorem~\ref{K-thm:barrier}, and the second is the entry point of the decoupled family. The
entry region has now been swept: over the $105{,}104{,}561$ squarefree $N_0\le10^{18}$ with
$\sigma(N_0)\ge\tfrac32$, the form $Q_{N_0}$ represents $4N_0^2$ not once, so there is no live slot
and no ghost either \textup{(\texttt{code/live\_slot\_census2.rs})}. Two checks accompany it. The
same program without the mass gate returns the two representations of Proposition~\ref{prop:form}
exactly, at $N_0=30$ with $(s,d)=(11,1)$ and $N_0=1722$ with $(83,1)$, both carrying $r=1$; and the
primorial bases $\Pi_{10},\dots,\Pi_{13}$ were re-checked by an independent implementation, each
admitting no representation. The sweep uses the reduction $A\mid N_0'(N_0'-2d^2)$, which follows
from $2N_0\equiv N_0'\pmod A$ and lets a single quadratic non-residue retire a base without any
search; it retires $45\%$ of them. The same program run without the
mass gate returns the two representations of Proposition~\ref{prop:form} exactly, at $N_0=30$ with
$(s,d)=(11,1)$ and $N_0=1722$ with $(83,1)$, both with $r=1$, which is the control that the search
sees what it should.

\subsection{Classical layers: genus, Giuga, parity}\label{ssec:classical}

\begin{proposition}[A principal--genus necessary condition, and a cascade on solutions]\label{prop:genus}
As $\gcd(N_0,N_0')=1$, the value $4N_0^2$ is a square coprime to the odd part of $\operatorname{disc}
Q_{N_0}$, so every genus character is $+1$ on it and a primitive form represents it only from the
principal genus \cite{Cox}. Reading the odd assigned characters off $Q_{N_0}(1,1)=4N_0$ gives the
computable necessary condition
\[ (N_0\mid\ell)=+1 \quad\text{for every odd prime }\ell\mid (N_0'^2-4N_0^2)\text{ to odd multiplicity} \]
(a $2$--adic condition accompanies it; for odd $N_0$ of even $\omega$ the form is imprimitive and must be
primitivised first; so we state the criterion on the odd part, with a $2$--adic caveat). Both ghost
bases $30,1722$ satisfy it, and about $16\%$ of squarefree bases pass in the ranges tested (an
\emph{upper bound}, using only the odd characters, and mildly range--dependent). \emph{Cascade.} In a
solution of \#307 with $N=D_PD_Q$, each $\ell\in P\cup Q$ gives a cofactor $N/\ell$ realising the
Pythagorean structure of $N$ over base $N/\ell$; every cofactor with $\sigma(N/\ell)<2$ (all of them
for a minimal first--$59$--primes solution) must then satisfy the principal--genus criterion, a
cascade of $\ge59$ computable genus conditions complementing the quadratic--residue filter of
Corollary~\ref{cor:qr}. \textup{(The further class--group and slot--admissibility stages of the
fertility gap we observe only numerically at toy scale.)}
\end{proposition}

Below the barrier the form has exactly one source of solutions, and it is entirely classical.

\begin{proposition}[The sub--barrier ghosts are pronic Giuga numbers]\label{prop:giuga}
Every representation of $4N_0^2$ by $Q_{N_0}$ over squarefree $N_0\le3\times10^4$ is a ghost
$(s,d)=(\sqrt{4N_0+1},1)$, occurring exactly when $N_0'=N_0+1$ with $N_0$ pronic; in range only
$30=5\cdot6$ and $1722=41\cdot42$. The two lines through such bases are classical: since $n'=n\,\sigma(n)$
for squarefree $n$,
\[ \{n:n'=n-1\}=\text{primary pseudoperfect numbers }(\textstyle\sum_{p\mid n}\tfrac1p+\tfrac1n=1),\qquad
   \{n:n'=n+1\}=\text{Giuga numbers (quotient one)}. \]
The derivative characterisations are due to Lava (OEIS \cite{OEIS_A054377,OEIS_A007850}, 2009; the
Giuga form $n'=n+1$ being Lava's \emph{conjecture}, with the generalisation $n'=an+1$, $a\in\NN$, the
theorem of Grau--Oller-Marc\'en \cite{GrauOllerMarcen2012}); the underlying number classes are
Butske--Jaje--Mayernik \cite{BJM2000} (exactly one primary pseudoperfect number per $\omega\le8$, eight
known) and Giuga \cite{Giuga1950,BBBG1996}, and the oblong chain ($N$ with $N+1$ prime breeds the next
term $N(N+1)$) is classical in both directions (Forgues and Sondow, OEIS \cite{OEIS_A054377}). To
$2\times10^6$ the lines are $\{n'=n-1\}=\{2,6,42,1806,47058\}$ and $\{n'=n+1\}=\{30,858,1722,66198\}$,
and the chain correctly predicts the next terms: $47059$ is prime, so $47058\cdot47059=2{,}214{,}502{,}422$
is the sixth primary pseudoperfect number, and $47057$ is prime, so the pronic
$47057\cdot47058=2{,}214{,}408{,}306$ is the fifth Giuga number; both already known, recovered here by
structure. Our contribution is the reframing: the chain step is the one--line Leibniz identity
$N'=N-1\Rightarrow(N(N+1))'=N'(N+1)+N^2=N(N+1)-1$, exhibiting the chain as orbits of the Artin--Schreier
map $x\mapsto x(x+1)=x^2+x$; the same map, with the same halt ($2\to6\to42\to1806$, stopping at
$1807=13\cdot139$), as the $\mathbb{F}_2[t]$ tower of Proposition~\ref{prop:ff-pyth}: characteristic $2$
and $\mathbb{Z}$ run identical breeders on the two lines.
\end{proposition}

\begin{remark}[Converse, and the bridge]
Every pronic solution of $n'=n+1$ is $m(m+1)$ with $m$ prime and $m+1$ primary pseudoperfect: Leibniz
gives $m'(m+1)+m(m+1)'=m^2+m+1$, and writing $m'=km+1$ forces $(m+1)'=(1-k)m-k$, which is $\ge0$ only for
$k=0$, whence $m'=1$ ($m$ prime) and $(m+1)'=m$. (The congruence $m'\equiv1\pmod m$ alone does \emph{not}
force $m$ prime; its composite solutions are precisely the Giuga numbers $30,858,1722,66198,\dots$,
so primality comes from the sign elimination $k=0$.) Thus the only sub--barrier life in the form equation
is the classical Giuga/primary--pseudoperfect world, known to sit in a web with Sylvester's sequence,
Zn\'am's problem and the Erd\H{o}s--Moser equation \cite{SondowMacMillan2017} and with the arithmetic
derivative \cite{GrauOllerMarcen2012,GrauOllerMarcenSadornil2021}; we claim none of this web, only that
\#307 meets it through $n'=n\,\sigma(n)$, and that its $a\ge2$ (equivalently $\sigma>2$) branch lands on
the same $\ge59$--prime barrier \cite{BorweinWong1997} that bars genuine cycles
(Corollary~\ref{K-cor:emptytest}).
\end{remark}

Writing both members of a cycle as members of these lines stratifies \#307 by the gap.

\begin{proposition}[Gap stratification: \#307 as adjacency of $\mu$--Sondow lines]\label{prop:gap}
For a squarefree two--cycle $a'=b$, $b'=a$ with $a>b$ and gap $d=a-b$, the smaller member lies on the line
$n'=n+d$ and the larger on $n'=n-d$ (as $b'-b=d$ and $a'-a=-d$), with $\gcd(d,ab)=1$ (a prime dividing $d$
and $a$ would divide $b=a-d$); conversely $x'=x+d$ with $(x+d)'=x$ is a two--cycle. These opposite--sign
lines are the quotient--one $\mu$--Sondow numbers of parameter $\mp d$ \cite{GrauOllerMarcenSadornil2021}.
Hence \#307 holds iff for some $d\ge1$ the lines $n'=n+d$ and $n'=n-d$ carry members exactly $d$ apart; for
$d=1$ this is \emph{a Giuga number adjacent to a primary pseudoperfect number}. By the parity dichotomy
(Theorem~\ref{K-thm:parity}) the mixed--parity case forces $d$ odd and the both--odd case ($d$ even) carries
the $10^{2519}$ bound, so every cycle below that bound has odd gap and $d=1$ is the lead stratum. A $d=1$
pair $(k,k+1)$ has exactly one odd member, hence requires either an odd Giuga number (known to have at least
$9$ prime factors \cite{BBBG1996,BorweinWong1997}) or an odd primary pseudoperfect number (none is known):
the lead stratum inherits two classical open problems. Tabulating cycle--eligible members (squarefree,
coprime to $d$) to $2\times10^6$ for all $d\le50$ gives $21$ on the plus lines (none of even gap, consistent
with the parity remark) and $81$ on the minus lines, with \emph{no} adjacent pair at any $d\le50$; the
barrier in gap--stratified form. \textup{(}For coprime $M$ on $n'=n-\alpha$ and $N$ on $n'=n+\beta$, Leibniz
gives $(MN)'-2MN=\beta M-\alpha N$, so $\beta M-\alpha N=\square$ would manufacture a minus--side square in
the sense of Proposition~\ref{prop:splitdisc}; the case $(\alpha,\beta)=(1,1)$, $M-N=1$ is the $d=1$ cycle
itself, and the known finite Giuga/primary--pseudoperfect lists contain no coprime pair, all members being
even.\textup{)} To our knowledge this gap--stratified equivalence and the pairing of opposite--sign lines at
fixed additive distance are new; the lines themselves are \cite{GrauOllerMarcenSadornil2021}.
\end{proposition}

The line parameter $\delta(n)=n'-n$ obeys an affine recursion that organises the classical members into a
genealogy.

\begin{proposition}[The $\delta$--recursion and the genealogy of the classical lines]\label{prop:delta}
With $\delta(n)=n'-n$, for squarefree $M$ and a prime $p\nmid M$,
\[ \delta(Mp)=p\,\delta(M)+M \qquad(\text{Leibniz; checked on }10^3\text{ random pairs}). \]
Solving $p\,\delta(M)+M=\pm d$ gives the child rule $p=(M\mp d)/e$ from a parent $M$ on the minus line
$\delta(M)=-e$. Every known Giuga number arises this way from a small--parameter parent: $30=6\cdot5$
($e=1$), $858=66\cdot13$ ($e=5$), $1722=42\cdot41$ ($e=1$), $66198=1122\cdot59$ ($e=19$),
$2{,}214{,}408{,}306=47058\cdot47057$ ($e=1$); the seed $47058$ itself resolves as $6\to66\to1518\to47058$
(with $\delta=-1,-5,-49,-1$), so the previously unexplained members $858$ and $66198$ acquire genealogies.
Two members of a cycle can never share a parent: siblings $M(M\mp d)/e$ differ by $2Md/e>d$, since
$e=M-M'<2M$. The oblong case $e=1$ is the classical chain of Proposition~\ref{prop:giuga} (its
$\mu$--Sondow form is \cite{GrauOllerMarcenSadornil2021}); the recursion, the general--$e$ genealogy, and the
resulting small--$|\delta|$ tree do not appear there.
\end{proposition}

\begin{corollary}[Nonemptiness families for the $\mu$--Sondow conjecture]\label{cor:gos}
In the notation of \cite{GrauOllerMarcenSadornil2021} (their $S_\mu$; the quotient--one slice is the line
$n'=n-\mu$), the recursion produces explicit members exceeding $|\mu|$ in three infinite parameter families:
$2p\in S_{p-2}$ for every odd prime $p$; $Mp\in S_{p-M}$ for every primary pseudoperfect number $M$ and
prime $p\nmid M$, and $Gp\in S_{-(p+G)}$ for every Giuga number $G$ and prime $p\nmid G$; in each case
$\mu/n+\sum_{q\mid n}1/q=1$ exactly (checked symbolically and on $151$ instances). This settles their
nonemptiness conjecture ($S_\mu\cap(|\mu|,\infty)\ne\varnothing$, their Conjecture~1(ii)) for infinitely
many parameters, including infinitely many outside their verified range $-145<\mu<673$; these families do
not appear among their constructions, which scale by $\operatorname{rad}(\mu)$ rather than extend by a new
prime. Their two open instances resist precisely because every quotient--one parent route hits a composite:
for $\mu=-145$, $M-145$ is composite for all eight known primary pseudoperfect numbers
($1661=11\cdot151$, $46913=43\cdot1091$, and the three larger cases, including the recently found eighth
member) while $p+G=145$ has no prime solution; for $\mu=673$, $M+673$ is composite for all eight
($679=7\cdot97$, $2479=37\cdot67$, $47731=59\cdot809$, the rest ending in $5$) and $p=675$ fails. A prime
in either family at the next generation would settle a published open instance; an exhaustive search found
none. Concretely: no quotient--one member of $S_{-145}$ or $S_{673}$ exists with \emph{any} prime--factor
cofactor $\le10^{11}$ (every squarefree parent $M\le10^{11}$ was scanned against the child rules, a region
whose members reach ${\sim}10^{22}$, together with a depth--two descent), the quotient--zero slice of
$S_{-145}$ is empty outright ($n'=145$ forces $n\le5256$ by $n'\ge2\sqrt n$, and the finite check is
empty), and quotients $\ge2$ force $\sigma\gtrsim2$, hence at least $59$ prime factors. Both instances
therefore remain open, with the quotient--one frontier moved from the $10^{10}$ of
\cite{GrauOllerMarcenSadornil2021} to cofactor scale $10^{11}$.
\end{corollary}

\begin{remark}[The look--back]
Along small--$|\delta|$ paths the forced prime is $p\approx M/e$, so members grow doubly exponentially and
the $10^{56}$--$10^{113}$ wall zone lies only a few generations from the seeds: a sparse, recursively
enumerable tree (a child needs $(M-j)/e$ prime for small $j$). The adjacency of two such trees at a fixed
additive distance places the final condition in the Pillai/$S$--unit genre (two multiplicatively
generated sequences meeting at a fixed additive distance) where fixed--support finiteness is classical but
growing--support statements are conjectural; a lens, not a transfer. Re--expressed in $\delta$--coordinates
the obstruction is once more $\sigma>1$ (positivity of $\delta$) plus integer reachability $\delta=\pm d$
under forced primes: the same additive--multiplicative exactness reappears in every coordinate tried (mass,
the deviation $k$, the form $Q_{N_0}$, genus theory, $\mu$--Sondow lines, $\delta$--dynamics). The reframings
buy structure, not a reduction.
\end{remark}

\begin{proposition}[Even--gap plus lines are empty below $3.23\times10^9$]\label{prop:oddthr}
An even squarefree $x=2u$ has $\delta(x)=2u'-u$ odd, so every line member of even $\delta$ is odd (with
$\omega$ odd). An odd member of a plus line has $\sigma>1$ over odd primes, and the least odd squarefree
integer of mass exceeding $1$ is $3\cdot5\cdots29=3{,}234{,}846{,}615$ (the $8$--prime product
$3\cdots23=111{,}546{,}435$ has mass $0.99896$). Hence every plus line of even gap $d$ is empty below
$3{,}234{,}846{,}615$, proving the census observation (no even--gap plus member to $2\times10^7$) and placing
an explicit floor under the smaller member of any both--odd configuration.
\end{proposition}

The lines also admit parity, infinitude, and congruence refinements.

\begin{proposition}[Parity floor for odd line members]\label{prop:parityfloor}
For odd squarefree $n$ every cofactor $n/p$ is odd, so $n'\equiv\omega(n)\pmod2$. Hence an odd primary
pseudoperfect number ($n'=n-1$, even) and an odd Giuga number ($n'=n+1$, even) both have $\omega(n)$ even;
combined with the classical bound $\omega\ge9$ \cite{BBBG1996,BorweinWong1997} this forces $\omega\ge10$, a
parity sharpening (for the primary pseudoperfect side, the eight smallest odd primes give $\sum1/p=0.99896<1$,
so $\omega\ge9$ directly). The floor cascades when small primes are avoided: for an odd squarefree member of
any plus line, or of a quotient--one line (where $\sigma(n)\ge1-1/n$), exact greedy/exchange computation
gives $\omega\ge27$ if $3\nmid n$; $\omega\ge66$ if $3,5\nmid n$; $\omega\ge144$ if $3,5,7\nmid n$; and
$\omega\ge16$ if $3\mid n$, $5\nmid n$; each avoided small prime multiplies the required complexity, and
for odd Giuga or primary pseudoperfect numbers the parity constraint rounds odd floors up (e.g.\ $3\nmid n$
forces $\omega\ge28$). (The restriction to plus or quotient--one lines is essential: on a general minus line
$\sigma$ can be small (every prime $p$ trivially satisfies $p'=p-(p-1)$) so no such floor exists
there.)
\end{proposition}

\begin{proposition}[Infinitely many gaps carry a nonempty plus line]\label{prop:plusinf}
By the recursion of Proposition~\ref{prop:delta}, $\delta(30p)=p\,\delta(30)+30=p+30$ for every prime $p>5$
(as $\delta(30)=1$), so the line $n'=n+(p+30)$ is nonempty; by Euclid infinitely many distinct gaps $d$ carry
a nonempty plus line, and the same runs from any Giuga number $G$ ($\delta(Gp)=p+G$). This fixes the census
tail: $330,390,510,570=30\cdot\{11,13,17,19\}$ lie on the gaps $d=41,43,47,49$. (Per--gap infinitude remains
open, of Bunyakovsky type.)
\end{proposition}

\begin{proposition}[A difference--$36$ congruence modulo $72$]\label{prop:mod72}
Let $n=6m$ be a Giuga or primary pseudoperfect number with $\gcd(m,6)=1$, $n'=n+\varepsilon$
($\varepsilon=+1$ Giuga, $-1$ primary pseudoperfect). Then $6m'=m+\varepsilon$, and with
$m'\equiv\omega(m)\pmod2$ this yields $n\equiv36\,\omega(n)-78\pmod{72}$ on the Giuga side and
$n\equiv36\,\omega(n)-66\pmod{72}$ on the primary pseudoperfect side; in particular equal $\omega$ forces
equal residue. The congruence is verified on all twelve Giuga numbers and all seven $6$--divisible primary
pseudoperfect numbers listed in the OEIS (we recomputed the factorisations of the five largest Giuga
members; all are squarefree, with $\omega=7,7,8,8,8$), and gives, at the modulus $72$, the
difference--$36$ arithmetic progression observed by Sondow--MacMillan \cite{SondowMacMillan2017}; their
finer modulus--$288$ pattern remains empirical.
\end{proposition}

The modulus--$72$ argument is the $D=6$ case of a general transport of the line equation across a fixed
divisor.

\begin{proposition}[Divisor transport]\label{prop:transport}
Let $n=Dm$ with $D,m$ coprime squarefree and $n'=n+\varepsilon$ ($\varepsilon=\pm1$) a quotient--one member.
Leibniz transports the line equation to
\[ D\,m'=(D-D')\,m+\varepsilon. \]
\begin{enumerate}[label=\textup{(\alph*)},leftmargin=2.2em]
\item If $\sigma(D)>1$ (so $D'>D$; least case $D=30$), the right side is negative for $m>1$ while the left
is nonnegative; the only escape is $m=1$, $\varepsilon=+1$. Hence \emph{no primary pseudoperfect number is
divisible by $30$, and the only Giuga number divisible by $30$ is $30$ itself} (elementary; we have not
located it in the literature), consistent with the full known corpus, of which only $30$ carries the
factor~$5$.
\item At a primary pseudoperfect block $D$ (where $D-D'=1$) the slice $m'=1$ reads $m=D-\varepsilon$ prime:
the two classical chain rules ($D(D+1)$ primary pseudoperfect when $D+1$ is prime, $D(D-1)$ Giuga when
$D-1$ is prime \cite{GrauOllerMarcenSadornil2021,SondowMacMillan2017}) are the two signs of one
transported equation ($1806=42\cdot43$ and $1722=42\cdot41$ sit on the same slice $D=42$).
\item For any gap $d$ the same transport stratifies the line $n'=n+d$ by divisor blocks $D$ into the steeper
lines $D\,m'=(D-D')m+d$, a further enumeration coordinate for the genealogy of
Proposition~\ref{prop:delta}.
\end{enumerate}
\end{proposition}

Finally, the form obstruction can be averaged: it eliminates almost every integer, though not the thin
stratum where a solution could live.

\begin{proposition}[Average rarity of form--barrier survivors]\label{prop:survivors}
Call a squarefree $N$ \emph{blocked} at an odd prime $\ell$ if $\ell\nmid N$, $N'\equiv\pm2N\pmod\ell$, and
$(N\mid\ell)=-1$, and a \emph{survivor} if it is blocked at no odd prime. A blocked $N$ admits no
representation $E\,s^2+D\,d^2=4N^2$ (reducing mod $\ell$ forces $N$ to be a quadratic residue, against
$(N\mid\ell)=-1$; this is the genus fragment of Proposition~\ref{prop:genus}), so every union of a \#307
solution is a survivor. Now $N'\equiv N\,S_\ell(N)\pmod\ell$ with $S_\ell(N)=\sum_{p\mid N}p^{-1}$ an
\emph{additive} function, so blocking at $\ell$ is an additive--function event ($S_\ell\equiv\pm2$) twisted by
the quadratic character, of limiting density $1/(\ell+1)$ \textup{[}\textsc{heuristic}: this constant
is asserted, not derived. Equidistribution of $S_\ell$ over $\mathbb{Z}/\ell$ together with density
$\tfrac12$ for $(N\mid\ell)=-1$ would give $2/\ell\cdot\tfrac12=1/\ell$ instead. The two agree to
leading order and $\sum_\ell1/\ell$ diverges as well, so the shape of the conclusion is unaffected;
the exact constant is not established here\textup{]}. As $\sum_\ell1/(\ell+1)$ diverges (Mertens), the
Selberg--Delange/Halász theory of additive functions in residue classes \cite{Tenenbaum} gives that the
survivors have natural density zero among the squarefree integers (unconditional), with
$\#\{\text{survivors}\le X\}=O\!\bigl(X/\log\log\log X\bigr)$ under a standard, but here unproved,
uniformity estimate. Thus the single--prime form obstruction removes \emph{almost every} squarefree integer.
We stress the scope: survivors remain \emph{infinite}, and a \#307 solution; one integer beyond
$2.09\times10^{56}$, in the density--zero stratum $\sigma(N)\ge2$; is untouched by such an average
statement, which speaks only to the typical integer, not to existence.
\end{proposition}

\appendix

\section{The minus--layer density theorem, and the square sieve that does not reach its rate}\label{app:sieve}

Relegated from \S\ref{ssec:sieve}. This appendix is the historical route to
Theorem~\ref{thm:a9}: a square sieve at $P\asymp(\log N)^{1/4}$ whose analytic input was never
supplied, since the uniformity it needs is exactly what Remark~\ref{rem:charcancelerratum}
forecloses at that range. The theorem is now proved instead by
Proposition~\ref{prop:condrate}, at the much smaller $P=\log\log N/(\log\log\log N)^{1+\varepsilon}$, where the same
uniformity \emph{is} available (Lemma~\ref{lem:charcancelunif}) (at the cost of a far weaker rate.
The sieve is kept in full because its algebraic layer is correct and reusable, because the located
gap is what identified the workable range, and because Theorem~\ref{cor:a9rate}) now proved, via
Lemma~\ref{lem:swdirect}, is what this route gives once its input is supplied.

\begin{lemma}[Cancellation in the character sums]\label{lem:charcancel}
Fix an odd squarefree $r>1$. Then
\[ T_r(N)\;:=\!\!\sum_{\substack{m\le N,\ \mu^2(m)=1\\ \gcd(m,r)=1}}\!\!\Bigl(\frac{m'-2m}{r}\Bigr)\;=\;o_r(N)
\qquad (N\to\infty). \]
The statement is qualitative and per fixed $r$: the implied constant depends on $r$, and no
uniformity is claimed. Remark~\ref{rem:charcancelerratum} records why the uniform version is not available; Remark~\ref{rem:charcancelsharp}
records what the constant is and is not.
\end{lemma}
\begin{proof}
Five standard facts are used, all at fixed $r$: Hal\'asz's theorem; $|\tau(\chi)|=\sqrt r$ for a
primitive character; the classical zero--free region; $|\log L(s,\psi)|\ll_r\log\log(3+|\tau|)$ for
$\psi$ non--principal at $s=1+1/\log N+i\tau$; and Siegel--Walfisz. None is formalised.

\emph{The expansion.} Write $\chi=(\cdot\mid r)$. For $r$ odd squarefree the Jacobi symbol is a
primitive real character modulo $r$, so
$\chi(x)=\tau(\chi)^{-1}\sum_{t\bmod r}\chi(t)e_r(tx)$ holds for \emph{every} $x$, including
$\gcd(x,r)>1$, where both sides vanish. That case occurs: at $r=15$ already $m=2$ gives
$\sigma_r(m)-2=6$. By Lemma~\ref{lem:symbolfact},
$(m'-2m\mid r)=\chi(m)\,\chi(\sigma_r(m)-2)$, and expanding the second factor,
\[ T_r(N)=\tau(\chi)^{-1}\sum_{t\bmod r}\chi(t)\,e_r(-2t)\;S_t(N),\qquad
   S_t(N)=\!\!\sum_{\substack{m\le N,\ \mu^2(m)=1\\ \gcd(m,r)=1}}\!\!h_t(m), \]
where $h_t(m)=\chi(m)\prod_{p\mid m}e_r(tp^{-1})$, using that $\sigma_r$ is a sum over $p\mid m$ and
$e_r$ carries sums to products. Extend $h_t$ to $\mathbb{N}$ by $h_t(p^k)=0$ for $k\ge2$ and
$h_t(p)=0$ for $p\mid r$, the only extension preserving $S_t$. It is multiplicative with
$|h_t|\le1$, which is Hal\'asz's hypothesis; $|h_t|=1$ on the support only.

\emph{The distance.} Hal\'asz needs
$M=\min_{|\tau|\le\log N}\sum_{p\le N}\bigl(1-\operatorname{Re}h_t(p)p^{-i\tau}\bigr)/p\to\infty$.
Now $h_t$ is not a function of $m\bmod r$, being built from the factorisation: at $r=7$,
$h_1(3)=0.2225+0.9749i$ while $h_1(10)=-1$, though $3\equiv10$. Its restriction to primes is,
which is all that is needed. Put $g_t(a)=\chi(a)e_r(ta^{-1})$ on $(\mathbb{Z}/r)^\times$, so
$h_t(p)=g_t(p\bmod r)$ for $p\nmid r$, and expand $g_t=\sum_\psi c_\psi\psi$. Reindexing by the
involution $a\mapsto ta^{-1}$ of $(\mathbb{Z}/r)^\times$ and using $\chi(a^{-1})=\chi(a)$,
\[ c_{\psi_0}=\frac1{\varphi(r)}\sum_{a\in(\mathbb{Z}/r)^\times}\chi(a)e_r(ta^{-1})
   =\frac{\chi(t)\,\tau(\chi)}{\varphi(r)},\qquad |c_{\psi_0}|=\frac{\sqrt r}{\varphi(r)}. \]
This identity uses no primitivity: it holds for any quadratic character on any finite abelian group,
which is what the formalisation records. Primitivity enters twice elsewhere, for
$|\tau(\chi)|=\sqrt r$ and for the inversion at $\gcd(x,r)>1$ above.

Splitting the prime sum by $\psi$,
\[ \sum_{p\le N}h_t(p)p^{-1-i\tau}
   =c_{\psi_0}\!\!\sum_{p\le N,\,p\nmid r}\!\!p^{-1-i\tau}
    +\sum_{\psi\ne\psi_0}c_\psi\sum_{p\le N}\psi(p)p^{-1-i\tau}. \]
The first sum is at most $\sum_{p\le N}1/p=\log\log N+O(1)$ in modulus, trivially and uniformly in
$\tau$. For the coefficients, $|g_t|=1$ gives $\sum_\psi|c_\psi|^2=1$ by Parseval, hence
$\sum_{\psi\ne\psi_0}|c_\psi|\le\sqrt{\varphi(r)}$ by Cauchy--Schwarz; they are \emph{not} small
individually. Indeed $c_\psi=\varphi(r)^{-1}\sum_b(\chi\psi)(b)e_r(tb)$, so
$|c_\psi|=\sqrt{\operatorname{cond}(\chi\psi)}/\varphi(r)$, and the maximum $\sqrt r/\varphi(r)$ is
attained exactly on the $\prod_{p\mid r}(p-2)$ characters with $\chi\psi$ primitive: three of the
eight at $r=15$, five of the twelve at $r=21$. Their contribution is small
because the $L$--functions have no pole, and reaching that needs care: the cited bound governs
$\log L$, an Euler product over all $p$, while what is required is the truncation. Both routes go by partial summation against Siegel--Walfisz and both shed the same factor
$1+|\tau|$; what differs is the \emph{range}, and it is the range and not any shift that does the
work. On $[2,N]$ one is left with $\int_{\log2}^{\log N}u^{-A}du=O_A(1)$, so $(1+|\tau|)\cdot O(1)$
is $\ll\log N$, worse than trivial. On $[N,\infty)$, writing $A(v)=\sum_{p\le v}\psi(p)\ll_{A,r}
v(\log v)^{-A}$,
\[ \sum_{p>N}\psi(p)p^{-s}\;=\;-A(N)N^{-s}+s\!\int_N^\infty\!A(v)v^{-s-1}\,dv
   \;\ll\;(\log N)^{-A}+(1+|\tau|)\frac{(\log N)^{1-A}}{A-1}\;\ll\;(\log N)^{2-A}, \]
which is $o(1)$ for $A>2$ and $|\tau|\le\log N$. Note what is \emph{not} used: the damping
$p^{-1/\log N}$ attached to the shifted point $s=1+1/\log N+i\tau$ is bounded by $1$ and is simply
discarded above. It contributes nothing to the tail, and indeed the shift alone leaves the tail of
size $\int_1^\infty e^{-w}w^{-1}dw=E_1(1)=0.2194$. The shift costs $O(1)$ by Mertens, since
$\sum_{p\le N}\log p/(p\log N)\ll1$, and it is convenient for keeping the Euler logarithm absolutely
convergent, but it is not what makes the tail small. The zero--free region then gives
$|\log L(s,\psi)|\ll_r\log\log(3+|\tau|)$, the argument of the double logarithm being taken above
$e$ so that the bound is positive at $\tau=0$; the non--principal terms are
$O_r(\log\log(3+|\tau|))=O_r(\log\log\log N)$ uniformly for $|\tau|\le\log N$, since at the top of
that range the cited input gives $\log\log(3+\log N)\asymp\log\log\log N$ and not $O(1)$. Hence
\[ M\;\ge\;c_r\log\log N-O_r(\log\log\log N),\qquad c_r:=1-\frac{\sqrt r}{\varphi(r)}\;\ge\;1-\frac{\sqrt3}2
   \;=\;0.1339\ldots, \]
the extremum being at $r=3$: each local factor of $\sqrt r/\varphi(r)$ is $\sqrt p/(p-1)$, and
$(p-1)/\sqrt p\ge2/\sqrt3$ for every odd prime, so the supremum is the single factor at $r=3$
\textup{(\texttt{code/charcancel\_check.gp})}.

\emph{Conclusion.} With $T=\log N$, Hal\'asz gives
$S_t(N)=\sum_{m\le N}h_t(m)\ll N\bigl((1+M)e^{-M}+(\log N)^{-1/2}\bigr)$. Note $M\to\infty$ alone would not suffice:
at fixed $T$ the bound is only $\ll NT^{-1/2}$. Both limits are in play, and
$S_t(N)\ll_rN(\log N)^{-\min(c_r,1/2)+o(1)}=o(N)$, the $\tfrac12$ being the floor the
$(\log N)^{-1/2}$ term imposes however large $M$ becomes. Only $r=3$ and $r=5$ have $c_r\le\tfrac12$,
so for every other modulus, and in particular for every $r=pq$ arising in Theorem~\ref{thm:a9}, the
floor is the binding constraint. Only the $\varphi(r)$ values of $t$ coprime to $r$
contribute, since $\chi(t)=0$ otherwise, so
\[ |T_r(N)|\;\le\;r^{-1/2}\!\!\sum_{\gcd(t,r)=1}\!\!|S_t(N)|\;\le\;\frac{\varphi(r)}{\sqrt r}\,
   \max_t|S_t(N)|\;\le\;\sqrt r\,\max_t|S_t(N)|\;=\;o_r(N). \]
The net factor is a \emph{loss} of $\sqrt r$: the $\varphi(r)$ summands outweigh
$|\tau(\chi)|^{-1}$. At fixed $r$ that is a constant.
\end{proof}

\begin{remark}[What the constant $c_r$ is and is not]\label{rem:charcancelsharp}
Three things about $c_r=1-\sqrt r/\varphi(r)$ are worth recording, since each is easily overstated.

\emph{Sharpness is conditional.} The bound $M\ge c_r\log\log N-O_r(1)$ is attained only when the
main term $\chi(t)\tau(\chi)$ is real and positive. Since $\overline{\tau(\chi)}=\chi(-1)\tau(\chi)$
with $\chi(-1)=(-1)^{(r-1)/2}$, the Gauss sum is purely imaginary for $r\equiv3\pmod4$, and there
the true $M$ has constant $1$ rather than $c_r$. For $\chi(t)=-1$ the main term has the wrong sign, and no
admissible $\tau$ rotates it back, since $|\sum_{p\le X}p^{-1-i\tau}|$ collapses to
$O(\log\log(3+|\tau|))$ away from $\tau=0$; there the constant is $1+\sqrt r/\varphi(r)$. Measured
at $r=5$: the growth rate $\Delta M/\Delta\sum_p1/p$ tends to $0.440887\to1-\sqrt5/4=0.440983$ for
$t\equiv\pm1$, and to $1.558921\to1+\sqrt5/4=1.559017$ for the other two residues, whose $M(10^7)$
is $4.378$ against the leading term $1.341$. So the constant is attained for
$r\equiv1\pmod4$ and $\chi(t)=+1$, and is not sharp otherwise.

\emph{The finite-$X$ deficit is a constant, though the proved error term is not.} The bound carries
$-O_r(\log\log\log N)$, which is what the $L$--bound gives at the top of the $\tau$--range; what is
\emph{measured} at a single modulus is smaller and flat. At $r=5$, $t=1$ the deficit
$c_r\sum_{p\le X}1/p-M$ is $0.1399,0.1401,0.1402,0.1402$ at $X=10^4,\dots,10^7$, flat to three
places and still rising in the fourth, while $\log\log\log X$ rises by $28\%$ over the same range.
Both sums here run over \emph{all} $p\le X$: $h_t$ is extended by $h_t(p)=0$ at $p\mid r$, so those
primes contribute $1/p$ to $M$, and omitting them shifts every number in this remark
\textup{(\texttt{code/deficit\_check.gp})}. The $\log\log\log$ enters only
through uniformity in $\tau$.

\emph{The rate that follows, and the one that does not.} Hal\'asz gives
$T_r(N)\ll_rN(\log N)^{-c_r+o(1)}$. Since $c_3=0.134$ and $c_5=0.441$ are below $\tfrac12$, this does \emph{not} yield $T_r(N)\ll\sqrt r\,N(\log N)^{-1/2}$ at those moduli. Quantitatively the lemma is asymptotic in the strong sense: at $r=5$, with the
measured constants and an implied Hal\'asz constant of $1$, the bound first beats the trivial
$|T_r(N)|\le N$ at $\log N=143.0$, that is $N=10^{62.1}$. Measured against the honest trivial bound,
which is $0.5066N$ because $T_r$ has only that many terms, the crossing moves to $\log N=1013.7$,
that is $N=10^{440}$.
\end{remark}

\begin{remark}[No uniform bound on the character sum near the threshold]\label{rem:charcancelerratum}
For $|\tau|\ge1/\log N$ the bound $\bigl|\sum_{p\le N}p^{-1-i\tau}\bigr|\ll\log\log(3+|\tau|)$ fails in the range that matters. Near the threshold the sum has size about $\log(1/|\tau|)$, which is of order $\log\log N$,
not $\log\log(3+|\tau|)$, which is of order $1$. Measured at $N=2\times10^{7}$ with
$\tau=2/\log N=0.11897$: the left side is $2.6797$ against a claimed bound of $0.12883$, a violation
by a factor $20.8$; at $\tau=1/\log N=0.05948$ the left side is $2.9773$ against $0.11176$, a factor
$26.6$ \textup{(\texttt{code/charcancel\_check.gp})}. A second defect is that the inequality bounds the
untwisted sum $\sum p^{-1-i\tau}$, whereas the Hal\'asz distance $M$ requires the $h_t$--twisted sum.

Consequences, stated exactly. Neither defect touches the lemma as now proved, because that proof
does not use the offending step: it bounds the \emph{twisted} sum through the character expansion of
$h_t$, where the principal coefficient contributes at most $\sqrt r/\varphi(r)\le0.8661$ of
$\sum_{p\le N}1/p$ and the non--principal ones contribute $O_r(\log\log\log N)$, both uniformly in
$\tau$. The
trivial bound $\bigl|\sum_{p\le N}p^{-1-i\tau}\bigr|\le\sum_{p\le N}1/p$ is all that is needed of the
untwisted sum, so the false inequality is not merely repaired but unnecessary.

Two claims must be distinguished. What is shown false
above is the $r$--\emph{uniform form of one input}, the $\log L$ bound, whose lower half needs
$L(1,\psi)$ away from $0$ and has only $L(1,\psi)\gg r^{-1/2}$ effectively or $\gg_\varepsilon
r^{-\varepsilon}$ by Siegel, ineffectively. That is \emph{not} the same as the uniform form of the
lemma being false, and the distinction turned out to be the whole opening: the obstruction is
\emph{range--dependent}. At $r\le(\log\log N)^2$ one has $\log r\le2\log\log\log N$, so even a
Siegel--worst--case $|\log L(1,\psi)|$ costs $\eta\log r+O_\eta(1)$, which is negligible against the
main term $\log\log N$; that is Lemma~\ref{lem:charcancelunif}, and with Lemma~\ref{lem:deficit} it
proves Theorem~\ref{thm:a9} with a rate. One caveat on how much this explains: at the far larger
moduli of Theorem~\ref{cor:a9rate} the $L$--function bound does also fail, but it is not what
binds there. Even an $O(1)$ per--character bound would leave the requirement
$\sqrt{\varphi(r)}=o(\log\log N)$, from the \emph{coefficient mass} rather than from any
$L$--function, a limitation of the character expansion, not of the arithmetic; see the blocker
recorded at that corollary and the question posed in Remark~\ref{rem:a9ceiling}. What this forecloses is uniformity at the far larger $r$ that Theorem~\ref{cor:a9rate} needs, where $\log r$ is comparable
to the main term, and that foreclosure binds only the route that expands in characters. The rate is proved in Theorem~\ref{cor:a9rate} by a route that does not expand \textup{(}Lemma~\ref{lem:swdirect}\textup{)}.
\end{remark}

\begin{lemma}[Cancellation, uniform at polyloglog moduli]\label{lem:charcancelunif}
\textup{(Label discipline, as for Theorem~\ref{thm:a9}: this statement went through six rounds of
adversarial review, each by a session that did not write the text under review. Rounds one to five
each found defects, twice in a repair made by the round before, which were repaired; round
six, auditing statement and proof line by line and recomputing every printed constant, found
nothing, which is what the label records. Proposition~\ref{prop:condrate} passed clean at round
three.)}
Fix $\varepsilon>0$ and write $L=\log\log N$. For every odd squarefree $r$ with
$7\le r\le L^2/(\log L)^{2+\varepsilon}$,
\[ T_r(N)\;\ll\;N\,(\log N)^{-2/5}, \]
with an implied constant that depends only on $\varepsilon$ and is ineffective; the same bound, by the
same proof with $e_r(+2t)$ in place of $e_r(-2t)$, holds for the sum with $(m'+2m\mid r)$ in place
of $(m'-2m\mid r)$. The moduli $r=3,5$ are excluded: there $c_r<1/2$, and they are
covered instead at fixed $r$ by Lemma~\ref{lem:charcancel} with the exponents of
Remark~\ref{rem:charcancelsharp}.
\end{lemma}
\begin{proof}
The proof of Lemma~\ref{lem:charcancel} is followed step by step; what is new is that every
implied constant is tracked in $r$, and each is absolute in the stated range. Here
$r\le L^2/(\log L)^{2+\varepsilon}\le L^2$, so $\log r\le2\log L$. The range is exactly the one
the argument supports: the binding requirement below is $\sqrt r\,\log L=o(L)$, that is
$r=o\bigl((L/\log L)^2\bigr)$.

\emph{Expansion.} Unchanged: $|T_r(N)|\le r^{-1/2}\sum_{t}|S_t(N)|\le r^{1/2}\max_t|S_t(N)|$,
the sum over the $\varphi(r)$ values of $t$ with $\gcd(t,r)=1$, since $\chi(t)=0$ otherwise; the
loss $r^{1/2}\le L$ is absorbed below.

\emph{The distance.} As before,
$\sum_{p\le N}h_t(p)p^{-1-i\tau}=c_{\psi_0}\sum_{p\le N,\,p\nmid r}p^{-1-i\tau}
+\sum_{\psi\ne\psi_0}c_\psi S_\psi(\tau)$ with $S_\psi(\tau)=\sum_{p\le N}\psi(p)p^{-1-i\tau}$
and $|c_{\psi_0}|=\sqrt r/\varphi(r)$. The principal part is at most
$(\sqrt r/\varphi(r))\log\log N$ in modulus, trivially and uniformly in $\tau$; dropping the
primes $p\mid r$ from the untwisted sum costs $\sum_{p\mid r}1/p\le\omega(r)/3+O(1)=O(\log L)$,
absorbed into the error below.

\emph{The non--principal sums, uniformly.} Fix $\psi\ne\psi_0$ mod $r$ and $|\tau|\le\log N$, and
put $s=1+1/\log N+i\tau$. For composite $r$ the character $\psi$ need not be primitive; let
$\psi^{*}$ of conductor $d\mid r$, $d>1$, induce it. Then $S_\psi$ and $S_{\psi^{*}}$ differ only at
the primes $p\mid r$, by at most $\sum_{p\mid r}1/p=O(\log L)$, absorbed into the bound below, and
Siegel--Walfisz, the zero--free region and Siegel's theorem are applied to $\psi^{*}$. Partial summation against Siegel--Walfisz, exactly as in
Lemma~\ref{lem:charcancel} with $A=3$, replaces the truncated sum by
$\log L(s,\psi)+O(1)$ at a cost $O((\log N)^{-1})$; Siegel--Walfisz is uniform, with an absolute
ineffective constant, for conductors up to $(\log N)^{3}$, and $r\le L^{2}\le(\log N)^{3}$ holds
for all $N\ge16$. For $\log L(s,\psi)$ two cases arise. If $L(\cdot,\psi)$ has no exceptional
zero, the classical zero--free region gives
$|\log L(s,\psi)|\ll\log\log\bigl(r(3+|\tau|)\bigr)\ll\log L$ with an absolute constant. If
$\psi$ is real with an exceptional zero $\beta$, split the $\tau$--range. For $|\tau|\le1$ the local
decomposition $\log L(s,\psi)=\log(s-\beta)+O\bigl(\log\log(r(3+|\tau|))\bigr)$ applies, and
\[ |s-\beta|\;\ge\;\operatorname{Re}(s)-\beta\;>\;1-\beta\;\ge\;c(\eta)\,r^{-\eta} \]
by Siegel's theorem in zero--location form, applied with any fixed $\eta>0$, a parameter
unrelated to the $\varepsilon$ of the range, so
$|\log(s-\beta)|\le\log\bigl(1/(1-\beta)\bigr)+O(1)\le\eta\log r+O_\eta(1)
\ll\eta\log L+O_\eta(1)$; the decomposition is \emph{not} asserted beyond $|\tau|\le1$,
where $|s-\beta|$ can be as large as $\log N$ and $\log|s-\beta|$ as large as $L$. For $|\tau|>1$
the zero $\beta$ lies at distance greater than $1$ from $s$, outside the local zero--density disc,
and the unconditional bound $|\log L(s,\psi)|\ll\log\log(r(3+|\tau|))$ applies with no exceptional
case. In both regimes $|S_\psi(\tau)|\ll\log L$, absolutely and ineffectively, uniformly in
$|\tau|\le\log N$.

\emph{Assembly.} By Cauchy--Schwarz on the coefficients,
$\sum_{\psi\ne\psi_0}|c_\psi|\le\sqrt{\varphi(r)}\le\sqrt r\le L/(\log L)^{1+\varepsilon/2}$, so the
non--principal contribution to the distance is $\ll L(\log L)^{-\varepsilon/2}=o(L)$, and
\[ M\;\ge\;\Bigl(1-\frac{\sqrt r}{\varphi(r)}\Bigr)L\;-\;O\bigl(L(\log L)^{-\varepsilon/2}\bigr)
   \;\ge\;\Bigl(1-\frac{\sqrt r}{\varphi(r)}-o(1)\Bigr)L . \]
The constant $1-\sqrt r/\varphi(r)$ does \emph{not} tend to $1$ over the whole range: it is
$0.133\ldots$ at $r=3$ and $0.441\ldots$ at $r=5$, and exceeds $1/2$ for every odd squarefree
$r\ge7$, with minimum $1-\sqrt{15}/8=0.515\ldots$ at $r=15$. Indeed
$\sqrt r/\varphi(r)=\prod_{p\mid r}\sqrt p/(p-1)$, every factor is $<1$, and $p\mapsto\sqrt p/(p-1)$
is decreasing on $p\ge3$; so adjoining a prime only shrinks the product, and over $\omega(r)\ge2$
the supremum is attained at $r=15$, giving $\sqrt r/\varphi(r)\le\sqrt{15}/8=0.4841\ldots$, while
for $\omega(r)=1$ with $p\ge7$ it is at most $\sqrt7/6=0.4409\ldots$. Only $r=3,5$ exceed $1/2$
\textup{(}values confirmed for $r\le20001$ in \texttt{code/uniform\_min\_sweep.gp}\textup{)}. Hal\'asz
(Montgomery's form; the constants are absolute and the function enters only through $|h_t|\le1$)
then gives $\max_t|S_t(N)|\ll N(1+M)e^{-M}+N(\log N)^{-1/2}\ll N(\log N)^{-1/2+o(1)}$ for
$r\ge7$, and the expansion loss $r^{1/2}\le L$ leaves
$T_r(N)\ll N(\log N)^{-1/2+o(1)}\ll N(\log N)^{-2/5}$.

Two honesty notes, as in Remark~\ref{rem:charcancelsharp}. The $o(1)$ terms above are vacuous at
every computable $N$: the crux inequality $\sqrt r\,\log L\ll L$ first bites at the top of the
range around $\log\log N\approx e^{36}$, and the widening over the earlier $r\le L^{2-\delta}$ is
asymptotic only, $L^{\delta}$ overtaking $(\log L)^{2+\varepsilon}$ far beyond even that. And the constant is ineffective twice over,
through Siegel--Walfisz and through Siegel's theorem; no effective form is claimed.
\end{proof}

\begin{lemma}[The distance, without expanding: uniform to power-of-log moduli]\label{lem:swdirect}
Fix $A>0$ and write $L=\log\log N$. Let $r>1$, let $g:(\mathbb Z/r)^\times\to\mathbb C$ satisfy
$|g|\le1$, and let $h$ be the multiplicative function with $h(\ell)=g(\ell\bmod r)$ for primes
$\ell\nmid r$, $h(\ell)=0$ for $\ell\mid r$, and $h(\ell^k)=0$ for $k\ge2$. Put
$\widehat g_0=\varphi(r)^{-1}\sum_{a}g(a)$, the mean of $g$, and let
\[ M\;=\;\min_{|\tau|\le\log N}\ \sum_{\ell\le N}\frac{1-\operatorname{Re}h(\ell)\ell^{-i\tau}}{\ell} \]
be Hal\'asz's distance for $h$. Then
\[ M\;\ge\;\bigl(1-|\widehat g_0|\bigr)\log\log N\;-\;o(\log\log N), \]
uniformly for $r\le(\log N)^A$ and for all such $g$, the $o(\cdot)$ being ineffective and of
unspecified rate.
\end{lemma}
\begin{proof}
Only the values of $h$ at primes are used. Fix $\varepsilon\in(0,1)$ and put
$T_0=\exp\bigl((\log N)^{\varepsilon}\bigr)$; $\varepsilon$ is an auxiliary of the proof and is
removed at the end. Since $\sum_{\ell\le N}1/\ell=L+O(1)$, it suffices to bound
$\bigl|\sum_{\ell\le N}h(\ell)\ell^{-1-i\tau}\bigr|$ uniformly in $|\tau|\le\log N$.

\emph{Small primes.} Trivially
$\bigl|\sum_{\ell\le T_0}h(\ell)\ell^{-1-i\tau}\bigr|\le\sum_{\ell\le T_0}1/\ell
=\log\log T_0+O(1)=\varepsilon L+O(1)$, with no $\tau$--dependence.

\emph{Large primes.} Every $\ell>T_0$ is coprime to $r\le(\log N)^A$, so $h(\ell)=g(\ell\bmod r)$
there. For $a\in(\mathbb Z/r)^\times$ write $E_a(x)=\pi(x;r,a)-\operatorname{li}x/\varphi(r)$ and
$E(x)=\sum_a g(a)E_a(x)$, and set
\[ S(x)=\!\!\sum_{T_0<\ell\le x}\!\!h(\ell)
      =\widehat g_0\bigl(\operatorname{li}x-\operatorname{li}T_0\bigr)+R(x),\qquad
   R(x)=E(x)-E(T_0), \]
using $\sum_a g(a)=\varphi(r)\widehat g_0$; in particular $R(T_0)=0$ identically. Siegel--Walfisz
with the exponent $B=A/\varepsilon$ gives $|E_a(x)|\ll_B x\exp(-c_B\sqrt{\log x})$ uniformly in $a$,
legitimately for every $x\ge T_0$, since then
$(\log x)^{A/\varepsilon}\ge(\log T_0)^{A/\varepsilon}=(\log N)^{A}\ge r$. Summing over the
$\varphi(r)\le r$ classes with $|g|\le1$,
\[ |R(x)|\;\ll\;r\Bigl(x\,e^{-c\sqrt{\log x}}+T_0\,e^{-c\sqrt{\log T_0}}\Bigr)\;\ll\;r\,x\,e^{-c\sqrt{\log x}}, \]
the second term being dominated because $y\mapsto y\,e^{-c\sqrt{\log y}}$ increases once
$\log y>c^2/4$. The factor $r$ comes from that class sum, not from Siegel--Walfisz, whose error
carries no modulus.

Now integrate. Writing the large--prime sum as a Stieltjes integral against $dS$,
\[ \sum_{T_0<\ell\le N}\!\!h(\ell)\ell^{-1-i\tau}
   =\widehat g_0\!\int_{T_0}^{N}\!\frac{x^{-i\tau}}{x\log x}\,dx
    +\Bigl[x^{-1-i\tau}R(x)\Bigr]_{T_0}^{N}+(1+i\tau)\!\int_{T_0}^{N}\!\!R(x)\,x^{-2-i\tau}dx. \]
The first term has modulus at most
$|\widehat g_0|\int_{T_0}^N dx/(x\log x)=|\widehat g_0|(1-\varepsilon)L$. The bracket vanishes at
$T_0$ and is $\ll r\,e^{-c\sqrt{\log N}}=o(1)$ at $N$. For the last, substituting $u=\log x$ and
using $\int_U^\infty e^{-c\sqrt u}\,du=2c^{-2}(1+c\sqrt U)e^{-c\sqrt U}$,
\[ (1+|\tau|)\!\int_{T_0}^{N}\!|R(x)|x^{-2}dx
   \;\ll\;(1+\log N)\,r\!\int_{(\log N)^{\varepsilon}}^{\log N}\!\!e^{-c\sqrt u}du
   \;\ll\;(1+\log N)(\log N)^{A}(\log N)^{\varepsilon/2}e^{-c(\log N)^{\varepsilon/2}}, \]
which is $o(1)$ uniformly in $|\tau|\le\log N$, since
$c(\log N)^{\varepsilon/2}$ eventually exceeds $(A+1+\tfrac\varepsilon2)\log\log N$. Collecting,
\[ \Bigl|\sum_{\ell\le N}h(\ell)\ell^{-1-i\tau}\Bigr|
   \;\le\;\varepsilon L+|\widehat g_0|(1-\varepsilon)L+O(1) \]
uniformly in $|\tau|\le\log N$, whence
$M\ge L-\varepsilon L-|\widehat g_0|(1-\varepsilon)L-O(1)=(1-\varepsilon)\bigl(1-|\widehat g_0|\bigr)L-O(1)$.

\emph{Removing $\varepsilon$.} For each $j\ge1$ take $\varepsilon_j=1/(2j)$ and $B_j=2Aj$, both
\emph{fixed}, so that $c_{B_j}$ and the implied constant $C_j$ of Siegel--Walfisz are constants
(ineffective, but constants). The display above then holds beyond a threshold $N_j$ depending only
on $(A,j)$, and
\[ \Bigl(1-\tfrac1{2j}\Bigr)\bigl(1-|\widehat g_0|\bigr)L-C_j
   \;\ge\;\bigl(1-|\widehat g_0|\bigr)L-\tfrac1jL
   \iff \tfrac{L}{2j}\bigl(1+|\widehat g_0|\bigr)\ge C_j, \]
which holds once $L\ge2jC_j$; enlarge $N_j$ so that it does, and so that $N_j$ increases. Setting
$j(N)=\max\{j:N\ge N_j\}\to\infty$ gives
$M\ge(1-|\widehat g_0|)L-L/j(N)=(1-|\widehat g_0|)L-o(L)$, uniformly in $r$ and $g$ as claimed.
The ineffectivity is exactly Siegel's, entering through Siegel--Walfisz; no exceptional--zero case
split is needed, because no $L$--function is bounded anywhere above.
\end{proof}

\emph{What this changes, and what it does not.} The constant is \emph{not} an improvement: at $g_t(a)=\chi(a)e_r(ta^{-1})$ one has
$\widehat g_0=\chi(t)\tau(\chi)/\varphi(r)$, so $|\widehat g_0|=\sqrt r/\varphi(r)$ and
Lemma~\ref{lem:swdirect} returns the same $1-\sqrt r/\varphi(r)$ that
Lemma~\ref{lem:charcancelunif} already gives; and in the error term it is worse, an $o(L)$ of
unspecified rate against that lemma's explicit $O\bigl(L(\log L)^{-\varepsilon/2}\bigr)$. Both are
ineffective, so ineffectivity is not the difference; the rate is.

What changes is the \emph{range}, and by an exponential. Lemma~\ref{lem:charcancelunif} is capped at
$r\le L^2/(\log L)^{2+\varepsilon}$ because expanding $g$ in multiplicative characters pays
$\sum_{\psi\ne\psi_0}|c_\psi|\asymp\sqrt r$ against a per--character bound $O(\log L)$. The proof
above never expands, so never pays it; the price is instead a factor $r$ against
$\exp(-c\sqrt{\log x})$, which is free at $r\le(\log N)^A$. Remark~\ref{rem:a9ceiling} posed exactly
this substitution as its open question, and the measurement recorded there, that
$\max_s|\sum_{\ell\le10^6}e_p(s\ell^{-1})/\ell|$ stays near $1.3$ while $\sqrt p$ grows from $3.6$
to $44.8$, is the numerical shadow of the same fact.

Two honesty notes. The $o(L)$ is ineffective and of unspecified rate, so nothing here is
quantitative at any particular $N$; and the error estimate is vacuous at every computable $N$, the
bound $(1+\log N)(\log N)^{A}(\log N)^{\varepsilon/2}e^{-c(\log N)^{\varepsilon/2}}$, at
$A=1,\varepsilon=0.2,c=\tfrac12$, \emph{peaking} at $\log N=10^{16.23}$ with value
$9.29\times10^{24}$, and falling below $1$ only past $\log N=10^{23.58}$
\textup{(\texttt{code/swdirect\_check.gp})}. A reading of ``turning over near $10^{18}$'' is contradicted by the certificate itself in both directions. The uniform Mertens estimate underlying the untwisted
case is classical (Norton, \emph{Illinois J. Math.} \textbf{20} (1976), Lemma 6.3, gives
$\sum_{p\le x,\,p\equiv\ell\,(k)}1/p=\log\log x/\varphi(k)+\delta_\ell/\ell+O(\log 3k/\varphi(k))$
with an absolute but ineffective constant) but we are not aware of a reference carrying the
$\tau$--uniformity used here, and deduce it from Siegel--Walfisz \textup{(}Montgomery--Vaughan,
\emph{Multiplicative Number Theory I}, \S11.3\textup{)} by the partial summation above.


\begin{lemma}[Two uniform deficit bounds]\label{lem:deficit}
\textup{(Label discipline, as for Lemma~\ref{lem:charcancelunif}: three rounds of adversarial review,
each by a session that did not write the text. The first two found five defects (including a
principal--coefficient value false in one of three cases, and a numerical certificate that tested the
wrong congruence with a silently inoperative coprimality filter) all repaired; the third found
nothing, which is what the label records.)}
Let $c\in\{2,-2\}$, $L=\log\log N$, and $P=L/(\log L\,w(L))$ with $w(L)\to\infty$. Uniformly
for primes $p\ne q$ in $(\tfrac12P,P]$, with implied constants that depend only on $w$ and are
ineffective,
\[ \#\{m\le N:\ \mu^2(m)=1,\ \gcd(m,p)=1,\ \sigma_p(m)\equiv c\ (\mathrm{mod}\ p)\}\ \ll\ N/p, \]
\[ \#\{m\le N:\ \mu^2(m)=1,\ \gcd(m,pq)=1,\ \sigma_p(m)\equiv c\ (\mathrm{mod}\ p)\ \text{and}\
   \sigma_q(m)\equiv c\ (\mathrm{mod}\ q)\}\ \ll\ N/(pq). \]
\end{lemma}
\begin{proof}
Here $p,q\in(P/2,P]$ and $pq\le P^2=L^2/(\log L\,w(L))^2$, inside
Lemma~\ref{lem:charcancelunif}'s range. Recall
$\sigma_p(m)=\sum_{\ell\mid m}\ell^{-1}\bmod p$, so that $e_p\bigl(t\sigma_p(m)\bigr)
=\prod_{\ell\mid m}e_p(t\ell^{-1})$ is multiplicative in squarefree $m$.

\emph{The second bound; the first is the same argument with one prime.} Detect both congruences
additively:
\[ \#\{\cdots\}\;=\;\frac1{pq}\sum_{s\bmod p}\sum_{t\bmod q}e_p(-sc)\,e_q(-tc)\,S_{s,t},\qquad
   S_{s,t}=\!\!\sum_{\substack{m\le N,\ \mu^2(m)=1\\ \gcd(m,pq)=1}}\!\!h_{s,t}(m), \]
where $h_{s,t}(m)=\prod_{\ell\mid m}e_p(s\ell^{-1})e_q(t\ell^{-1})$, extended by
$h_{s,t}(\ell^k)=0$ for $k\ge2$ and $h_{s,t}(\ell)=0$ for $\ell\mid pq$. It is multiplicative with
$|h_{s,t}|\le1$, which is Hal\'asz's hypothesis. The term $(s,t)=(0,0)$ contributes
$(pq)^{-1}\#\{m\le N:\mu^2(m)=1,\gcd(m,pq)=1\}\le N/(pq)$, which is the asserted bound; the phases
$e_p(-sc)e_q(-tc)$ are unimodular and play no further role, so the two values of $c$ are on the same
footing. It therefore suffices to prove
\[ \max_{(s,t)\ne(0,0)}|S_{s,t}(N)|\;\ll\;N/(pq), \tag{$\ast$} \]
since the $pq-1$ remaining terms then contribute $\ll N/(pq)$ in total.

\emph{The distance.} Put $g_{s,t}(a)=e_p(sa^{-1})e_q(ta^{-1})$ on $(\mathbb Z/pq)^\times$, so
$h_{s,t}(\ell)=g_{s,t}(\ell\bmod pq)$ for $\ell\nmid pq$, and expand
$g_{s,t}=\sum_\psi c_\psi\psi$ over the characters mod $pq$. By the Chinese remainder theorem
$g_{s,t}$ factors and so do its coefficients, $c_\psi=c^{(p)}_{\psi_1}c^{(q)}_{\psi_2}$ with
$\psi=\psi_1\psi_2$. At a single prime $p$ and $s\not\equiv0$, reindexing by $b=sa^{-1}$,
\[ c^{(p)}_{\psi_1}=\frac1{p-1}\sum_{a}e_p(sa^{-1})\overline{\psi_1}(a)
   =\frac{\overline{\psi_1}(s)}{p-1}\sum_{b\ne0}e_p(b)\psi_1(b), \]
so $c^{(p)}_{\psi_{1,0}}=-1/(p-1)$ by the Ramanujan sum $\sum_{b\ne0}e_p(b)=-1$, while
$|c^{(p)}_{\psi_1}|=\sqrt p/(p-1)$ for $\psi_1\ne\psi_{1,0}$ by $|\tau(\psi_1)|=\sqrt p$. For
$s\equiv0$ the factor $g$ is identically $1$ and $c^{(p)}_{\psi_{1,0}}=1$, the other coefficients
vanishing. Hence for $(s,t)\ne(0,0)$ at least one factor is a Ramanujan sum: the three cases
$s,t\ne0$; $s\ne0,t=0$; $s=0,t\ne0$ give $|c_{\psi_0}|$ equal to $1/((p-1)(q-1))$, $1/(p-1)$ and
$1/(q-1)$ respectively, so in every case, both $p$ and $q$ exceeding $P/2$,
\[ |c_{\psi_0}|\;\le\;\max\Bigl(\frac1{p-1},\frac1{q-1}\Bigr)\;\le\;\frac{2}{P-2}\;=\;o(1), \]
whereas Parseval $\sum_\psi|c_\psi|^2=1$ gives, by Cauchy--Schwarz,
$\sum_{\psi\ne\psi_0}|c_\psi|\le\sqrt{\varphi(pq)}\le\sqrt{pq}\le L/(\log L\,w(L))$.

\emph{The non--principal sums.} This is verbatim the corresponding step of
Lemma~\ref{lem:charcancelunif}, at modulus $pq\le L^2/(\log L\,w(L))^2$:
for $\psi\ne\psi_0$ and $|\tau|\le\log N$, partial summation against Siegel--Walfisz (legitimate
since $pq\le L^2\le(\log N)^3$) replaces the truncated sum by $\log L(s,\psi^*)+O(1)$ at the
inducing character, and the classical zero--free region gives
$|\log L(s,\psi^*)|\ll\log\log(pq(3+|\tau|))\ll\log L$; in the exceptional case the same
$\tau$--split applies, the branch $|\tau|\le1$ using
$|s-\beta|>1-\beta\ge c(\eta)(pq)^{-\eta}$ for any fixed $\eta>0$, unrelated to the
$w$ of the range, and hence
$|\log(s-\beta)|\le\eta\log(pq)+O_\eta(1)\ll\eta\log L+O_\eta(1)$, and
the branch $|\tau|>1$ needing no exceptional case. So $|S_\psi(\tau)|\ll\log L$ uniformly.

\emph{Assembly and Hal\'asz.} Therefore
\[ M\;=\;\min_{|\tau|\le\log N}\sum_{\ell\le N}\frac{1-\operatorname{Re}h_{s,t}(\ell)\ell^{-i\tau}}{\ell}
   \;\ge\;\bigl(1-o(1)\bigr)L\;-\;O\bigl(L/w(L)\bigr)\;=\;\bigl(1-o(1)\bigr)L, \]
the principal part costing at most $|c_{\psi_0}|L=o(L)$ and the non--principal part
$O\bigl(L/w(L)\bigr)=o(L)$. Hal\'asz's theorem (Montgomery's form; absolute constants, and
$h_{s,t}$ enters only through $|h_{s,t}|\le1$) with $T=\log N$ gives
\[ |S_{s,t}(N)|\;\ll\;N(1+M)e^{-M}+N(\log N)^{-1/2}\;\ll\;N(\log N)^{-1/2}, \]
the floor dominating since $(1+M)e^{-M}\ll(\log N)^{-1+o(1)}$. Finally
$pq\le L^2$, so $pq=o\bigl((\log N)^{1/2}\bigr)$ and hence
$N(\log N)^{-1/2}=o\bigl(N/(pq)\bigr)$, which is $(\ast)$. The
one--prime bound is the same computation with the single detector $p^{-1}\sum_s e_p(-sc)S_s$,
principal coefficient $-1/(p-1)$, $\sum_{\psi\ne\psi_0}|c_\psi|\le\sqrt p\le\sqrt P$, and
the requirement $\max_{s\ne0}|S_s|\ll N/p$, which is weaker than $(\ast)$.

Two notes, as for Lemma~\ref{lem:charcancelunif}. The constants are ineffective, through
Siegel--Walfisz and Siegel. And the distance constant here is $1-1/(p-1)\to1$, larger than the
$1-\sqrt r/\varphi(r)$ of Lemma~\ref{lem:charcancelunif}: the deficit bounds have more room than the
cancellation bound they are used beside, which is why no analogue of the $r\ge7$ restriction is
needed. Numerically, at $N=2\times10^6$ and $c=2$, the ratio of the
counted set to $N_{\mathrm{sf}}(p)/p$, where $N_{\mathrm{sf}}(p)$ counts the squarefree $m\le N$
coprime to $p$, lies in $[0.988,1.023]$ for $3\le p\le31$, and the two--prime ratio against
$N_{\mathrm{sf}}(pq)/(pq)$ lies in $[0.997,1.057]$ for pairs up to $(23,29)$: the densities are not
merely bounded but asymptotically exact, with no growth in $p$
\textup{(\texttt{code/deficit\_numeric.gp})}. The condition measured is $\sigma_p(m)\equiv2$, not
$m'\equiv2$; these differ by the unit $m$.
\end{proof}

\begin{proposition}[A rate on the minus layer]\label{prop:condrate}
Write $L=\log\log N$. For every function $w$ with $w(L)\to\infty$, however slowly,
\[ \#\{m\le N:\ \mu^2(m)=1,\ m'-2m\ \text{a positive perfect square}\}
   \;\ll_w\;N\,\frac{(\log\log\log N)^2\,w(L)}{\log\log N}, \]
which is the ceiling of the method \textup{(}Remark~\textup{\ref{rem:a9ceiling}}\textup{)}; in
particular the bound $\ll_\varepsilon N(\log\log\log N)^{2+\varepsilon}/\log\log N$ holds for every
fixed $\varepsilon>0$,
and the same for $m'+2m$. In particular Theorem~\ref{thm:a9} follows with a rate, by an
argument with no diagonalisation over $r$.
\end{proposition}
\begin{proof}
Let $P=L/(\log L\,w(L))$, $\mathcal P$ the primes in $(P/2,P]$, $\Pi=|\mathcal P|\gg
P/\log P$, and $a_m=m'-2m$. First, $a_m\ne0$ for every squarefree $m$: $a_m=0$ forces
$\sigma(m)=2$ exactly, impossible for $m>1$ since $\gcd(D(m),m)=1$ (Theorem~\ref{K-thm:structure})
and $a_1=-2$. If $a_m$ is a positive square then $(a_m\mid p)\ne-1$ for every $p$, and
$(a_m\mid p)=0$ exactly when $p\mid a_m$, so
$\sum_{p\in\mathcal P}(a_m\mid p)\ \ge\ \Pi-\#\{p\in\mathcal P:p\mid a_m\}$. Hence
$\Pi^2\,\#\{\text{squares}\}\le\sum_m\bigl(\sum_p(a_m\mid p)+\#\{p\in\mathcal P:p\mid a_m\}\bigr)^2$,
the sum over all squarefree $m\le N$; expanding the square, the cross term is bounded through
$\bigl|\sum_p(a_m\mid p)\bigr|\le\Pi$, which gives
\[ \Pi^2\,\#\{\text{squares}\}\ \le\ \sum_{p\ne q}\bigl(T_{pq}(N)+E_{pq}\bigr)\;+\;\Pi N
   \;+\;2\Pi\sum_{p}\#\{m:p\mid a_m\}\;+\;\sum_{p\ne q}\#\{m:\ p\mid a_m,\ q\mid a_m\}
   \;+\;\sum_p\#\{m:p\mid a_m\}, \]
where $|E_{pq}|\le N/p+N/q\le4N/P$ collects the $m$ with $\gcd(m,pq)>1$, dropped from $T_{pq}$'s summation
range exactly as in the argument for Theorem~\ref{thm:a9}. By
Lemma~\ref{lem:charcancelunif} at $r=pq\le L^2/(\log L\,w(L))^2$, inside its range once
$w(L)\ge(\log L)^{\varepsilon/2}$ or by enlarging $w$ (each $r$ odd squarefree,
$r\ge7$), $\sum_{p\ne q}|T_{pq}|\ll\Pi^2N(\log N)^{-2/5}$. By Lemma~\ref{lem:deficit} (applicable because $p\mid m$ forces
$a_m\equiv m'\equiv m/p\not\equiv0\pmod p$ for squarefree $m$, so the $m$ counted here are
automatically coprime to $p$, and likewise to $pq$ in the two--prime count)
$\sum_p\#\{m:p\mid a_m\}\ll N\Pi/P\ll N/\log P$ and the two--prime sum is $\ll\Pi^2N/P^2\ll
N/(\log P)^2$. Dividing by $\Pi^2$: the diagonal $\Pi N$ gives the main term
$N/\Pi\ll N\log P/P$, and every other term is smaller. Since $\log P\asymp\log\log\log N$, the
bound follows: $\log P\asymp\log L$, so $N\log P/P\asymp N(\log L)^2w(L)/L$. The fixed--$\varepsilon$
form is the case $w=(\log L)^{\varepsilon}$. Note $\mu^2(m)$ and the coprimality conditions are carried inside every count;
no M\"obius expansion is performed.
\end{proof}

Measured at $r=1009$, $N=3\times10^7$: $|\sum_m h_t(m)|/N$ equals $3.6\times10^{-6}$, $1.3\times10^{-3}$,
$2.1\times10^{-3}$, $5.7\times10^{-4}$, $1.5\times10^{-3}$ at $t=0,1,2,17,500$, all of size $\sqrt N$
rather than merely $o(N)$. The sums $T_r(N)$ themselves were measured at $N=5\times10^7$ for
$r=101,1009,10007,100003$ and $r=pq=101\cdot1009,\,1009\cdot10007$, giving
$|T_r|/\sqrt N=8.1,\,0.1,\,1.2,\,1.0,\,1.0,\,0.2$ \textup{(\texttt{code/sqsieve\_test.rs})}.

\begin{lemma}[Local densities at $p^2$]\label{lem:localsq}
For each odd prime $p$,
$\ \limsup_{N\to\infty}N^{-1}\#\{m\le N:\mu^2(m)=1,\ p^2\mid(m'-2m)\}\le p^{-2}$, with absolute
constant. \textup{(}The density form is the one Theorem~\ref{thm:a9} needs, because it sends
$P\to\infty$; the form ``$\ll N/p^2$ for each fixed $p$'' advertises a $p$--dependent implied
constant and would not license that limit.\textup{)}
\end{lemma}
\begin{proof}
For $p\nmid m$ the condition is $\sigma_{p^2}(m)\equiv2\pmod{p^2}$ by Lemma~\ref{lem:symbolfact}
applied with modulus $p^2$. Detecting it by additive characters, the non-trivial frequencies give
multiplicative sums $\prod_{\ell\mid m}e_{p^2}(t/\ell)$ whose values at primes depend only on
$\ell\bmod p^2$ and which are not Dirichlet characters, because $u\mapsto e_{p^2}(t/u)$ is not
multiplicative; the Hal\'asz argument of Lemma~\ref{lem:charcancel} applies, but \emph{not}
verbatim, and the difference is worth naming. That lemma uses primitivity of $\chi=(\cdot\mid r)$
twice, for $|\tau(\chi)|=\sqrt r$ and for the principal coefficient $\chi(t)\tau(\chi)/\varphi(r)$;
here $(\cdot\mid p^2)$ is trivial on $(\mathbb{Z}/p^2)^\times$, so there is no character twist at
all. The principal coefficient is instead a Ramanujan sum $c_{p^2}(t)/\varphi(p^2)$, which vanishes
for $\gcd(t,p^2)=1$ and equals $-1/(p-1)$ when $p\,\|\,t$; the distance bound follows more easily
than in Lemma~\ref{lem:charcancel}, not by the same computation. Hence the density is
$\tfrac6{\pi^2}\cdot\tfrac p{p+1}\cdot p^{-2}+o(1)\le p^{-2}$. The two factors in front are not optional: the count runs over \emph{squarefree} $m$ coprime to $p$, so it carries the
squarefree density $6/\pi^2$ and the coprimality factor $p/(p+1)$. The error was in the safe
direction, the stated $p^{-2}$ being an over-estimate, so the inequality the lemma asserts survives
with room to spare. The case $p\mid m$ does not arise at all: writing $m=pn$ with $p\nmid n$
gives $m'=pn'+n$, so $m'-2m=n+p(n'-2n)\equiv n\not\equiv0\pmod p$, and $p$ does not even divide
$m'-2m$, let alone $p^2$.
\end{proof}
Measured at $N=3\times10^7$ for $p=5,7,11,13,101$, the count divided by
$\tfrac6{\pi^2}\tfrac p{p+1}N/p^2$ is $0.842$, $0.887$, $0.930$, $0.925$, $0.993$, approaching $1$ as
the corrected density predicts. Against the bare $N/p^2$ the same data read $0.512$, $0.539$,
$0.566$, $0.562$, $0.604$, approaching $6/\pi^2=0.6079$.

\begin{theorem}[The minus layer has density zero on the squarefree stratum]\label{thm:a9}
\textup{Label: \textsc{proved}. Two independent routes reach this statement, and the label rests on
the second. The first is the square sieve at fixed $r$ recorded in
Remark~\textup{\ref{rem:a9original}} below: it needs no uniformity, since it sends $P\to\infty$ only
after $N$, but it yields no rate, and its \emph{text} failed three consecutive reviews, a load--bearing display silently shedding
a hypothesis each time, so it has never had a clean round and the label does not rest on it. The
second is Proposition~\textup{\ref{prop:condrate}}, which proves the statement \emph{with a rate}
from Lemma~\textup{\ref{lem:charcancelunif}} and Lemma~\textup{\ref{lem:deficit}}; each of those, and
the proposition itself, earned its label under this project's standard, a review by a session that
did not write it, finding nothing.}

$\#\{m\le N:\ \mu^2(m)=1,\ m'-2m\ \text{is a positive perfect square}\}=o(N)$; indeed, by
Proposition~\ref{prop:condrate}, the count is
$\ll_\varepsilon N(\log\log\log N)^{2+\varepsilon}/\log\log N$ for every $\varepsilon>0$. Equivalently, atom
$\mathrm{A}9$ holds on the squarefree stratum: the squarefree minus layer has density zero. The same
argument with $m'+2m$ gives the squarefree plus--analogue, so the squarefree Pythagorean layer has
density zero unconditionally. That is \emph{not} a reproof of
Proposition~\textup{\ref{prop:plusthin}}, which is a sharper $O(\sqrt X(\log X)^2)$ on the narrower
semiprime stratum; the two are incomparable, and that proposition rests on its own elementary
divisor--sum proof.

The squarefree restriction is not cosmetic. Every tool
in the proof, Rigidity, the cofactor identity, and Lemma~\ref{lem:charcancel}, is stated for
squarefree $m$; the non-squarefree stratum, of density $1-6/\pi^2\approx0.392$, is untreated. Since
a \#307 solution is squarefree by Rigidity, the restriction costs nothing for the application, but
the unrestricted statement is not proved here.
\end{theorem}
\begin{proof}
Immediate from Proposition~\ref{prop:condrate}, which gives the minus--layer count
$\ll_\varepsilon N(\log\log\log N)^{2+\varepsilon}/\log\log N=o(N)$ on the squarefree stratum for
any fixed $\varepsilon>0$, and the plus--layer count by the same statement. Remark~\ref{rem:a9original}
records an independent route to the qualitative form.
\end{proof}

\begin{remark}[The original route, and what it does and does not give]\label{rem:a9original}
This is the argument that carried Theorem~\ref{thm:a9} before
Proposition~\ref{prop:condrate} existed. It is recorded because it is independent of that
proposition, because it needs no uniformity in $r$ whatsoever, and because the errata inside it are
the ones that cost the most to find. It is \emph{not} what the label rests on: its text failed three
consecutive reviews and has never had a clean one, and it yields no rate, since $P\to\infty$ only
after $N$.

Write $a_m=m'-2m$ and fix $P$, letting $\mathcal{P}$ be the primes in $(P,2P]$ and
$\Pi=|\mathcal{P}|\sim P/\log P$. If $a_m$ is a perfect square coprime to every $p\in\mathcal{P}$
then $(a_m/p)=1$ for all such $p$, so $\sum_{p\in\mathcal{P}}(a_m/p)=\Pi$. Hence
\[ \Pi^2\,\#\{m\le N:\ \mu^2(m)=1,\ a_m=\square,\ \gcd(a_m,\textstyle\prod\mathcal{P})=1\}
   \;\le\!\!\sum_{\substack{m\le N\\ \mu^2(m)=1}}\!\!\Bigl(\sum_{p\in\mathcal{P}}\Bigl(\frac{a_m}{p}\Bigr)\Bigr)^2
   \;=\;\sum_{p,q\in\mathcal{P}}\ \sum_{\substack{m\le N\\ \mu^2(m)=1}}\Bigl(\frac{a_m}{pq}\Bigr). \]
Both $m$--sums carry $\mu^2(m)=1$ throughout, and so does the counting set on the left. This is not
a convenience: every tool below is squarefree--only, and writing these sums unrestricted makes the identity with $T_{pq}+E_{pq}$ false by the non--squarefree terms. Those
are not negligible: measured at $N=3\times10^7$ the omitted sum is $119{,}473$ against a $T_{pq}$ of
$2{,}699$ at $r=101\cdot1009$, and $169{,}949$ against $314$ at $r=1009\cdot10007$, exceeding the
term the lemma controls by a factor $44$ to $540$. The paper bounds it by nothing better than the
trivial $(1-6/\pi^2)N$.
The diagonal $p=q$ contributes $(a_m/p^2)=(a_m/p)^2\in\{0,1\}$, hence at most $\Pi N$. The
off-diagonal is $\sum_{p\ne q}\bigl(T_{pq}(N)+E_{pq}(N)\bigr)$, where $T_{pq}$ carries the
coprimality condition $\gcd(m,pq)=1$ of Lemma~\ref{lem:charcancel} and
$E_{pq}(N)=\sum_{m\le N,\ \mu^2(m)=1,\ \gcd(m,pq)>1}(a_m/pq)$ collects the terms it omits. That
correction is harmless but is not empty and must be carried:
$|E_{pq}(N)|\le\#\{m\le N:p\mid m\text{ or }q\mid m\}\le N(1/p+1/q)\le2N/P$ (the primes here
lie in $(P,2P]$, whereas Proposition~\ref{prop:condrate} takes them in $(P/2,P]$ and so carries
$4N/P$) so
$\Pi^{-2}\sum_{p\ne q}|E_{pq}(N)|\le2N/P$, which is dominated by the diagonal's own $N/\Pi\asymp
N\log P/P$ and changes nothing below. Each $T_{pq}(N)$ is $o(N)$, for each of the finitely many
pairs, by Lemma~\ref{lem:charcancel}. A square divisible by $p$ is divisible by $p^2$, so the excluded $m$ are counted by
Lemma~\ref{lem:localsq} in its \emph{density} form: $\limsup_NN^{-1}\#\{m\le N:\mu^2(m)=1,\
p^2\mid a_m\}\le p^{-2}$, with absolute constant $1$. That form is what the endgame needs, since
sending $P\to\infty$ against a $p$--dependent implied constant would not be legitimate; the
statement ``$\ll N/p^2$ for each fixed $p$'' advertises exactly such a constant. Summing, $\sum_{p\in\mathcal{P}}p^{-2}\le\Pi/P^2\asymp1/(P\log P)$. Here
$p\mid m$ is impossible, since $m=pn$ with $p\nmid n$ gives $a_m=n+p(n'-2n)\equiv n\not\equiv0$,
a computation that itself needs $\mu^2(m)=1$. The case $a_m=0$ does not arise: the theorem counts $a_m=\square>0$, and in any case
$m'=2m$ has \emph{no} squarefree solution, since $\sigma(m)=2\in\mathbb{Z}$ forces $D_m=1$ by
rigidity. Altogether
\[ \limsup_{N\to\infty}\frac1N\#\{m\le N:\ \mu^2(m)=1,\ a_m=\square>0\}\;\ll\;\frac1\Pi+\frac1{P\log P}. \]
The left side does not depend on $P$; letting $P\to\infty$ gives $0$. The restriction $\mu^2(m)=1$
must be carried here as everywhere above: without it the display would assert density zero for the
\emph{full} minus layer, which is exactly the statement this theorem disclaims, and three successive
revisions of this proof dropped it from one display or another.
\end{remark}

The proof of Theorem~\ref{thm:a9} sends $P\to\infty$ only after $N$, which costs all quantitative
information. Keeping $P$ a function of $N$ would turn it into an explicit rate, and what that
requires is uniformity in $r$ as $P\to\infty$.

\emph{That uniformity is exactly what is not available}. Lemma~\ref{lem:charcancel} gives
$M\ge\bigl(1-\sqrt r/\varphi(r)\bigr)\log\log N-O_r(\log\log\log N)$, with the error depending on
$r$ through the non--principal $L$--functions, and its statement and the remark following it both say so. The constant is $\sqrt r/\varphi(r)$, not the $3/\sqrt r$ printed here previously; the two
are not interchangeable, since $\varphi(r)\ge r/3$ fails for odd squarefree $r$, first at
$r=3\cdot5\cdots23=111{,}546{,}435$ with $\varphi(r)/r=0.3272$. The statement below was therefore
stated as a \textsc{conjecture} when this remark was written, and its argument here is a sketch
rather than a proof. The uniform input it wanted is Lemma~\ref{lem:swdirect}, and the rate itself is
Theorem~\ref{cor:a9rate}, proved by a different route.

\begin{remark}[The ceiling of the method, and the one input that would break it]\label{rem:a9ceiling}
Two things are worth recording about how far Proposition~\ref{prop:condrate} can be pushed, since
the same accounting error was made twice in this paper's history.

\emph{The ceiling.} The binding constraint there is again the coefficient mass:
$\sqrt{pq}\cdot O(\log\log\log N)=o(\log\log N)$, that is $P=o(L/\log L)$ with $L=\log\log N$. The
proposition takes $P=L/(\log L)^{1+\varepsilon}$, which is short of that ceiling by $(\log L)^{\varepsilon}$,  Taking $P=L/(\log L\,w(L))$ for any $w\to\infty$ gives
\[ \#\{m\le N:\ \mu^2(m)=1,\ m'-2m=\square>0\}\;\ll\;N\,\frac{(\log\log\log N)^2\,w(\log\log N)}{\log\log N}, \]
so the true ceiling of the method is $N(\log\log\log N)^{2+o(1)}/\log\log N$ with the $o(1)$
positive, unremovable, and arbitrarily slow. We state the proposition with a fixed $\varepsilon$
because it is the usable form; the sharpening is asymptotic and worth nothing numerically. What
matters is the shape: no choice of $P$ puts a power of $\log N$ in the denominator \emph{so long as
the distance is estimated through the character expansion}. That proviso is not decoration: see
the next paragraph, where it is exactly what is in question.

\emph{Where the loss actually is, and what it is not.} Both walls, this one and the corollary's, are the same $\sqrt{\varphi(r)}$ paid when the phase $\ell\mapsto e_r(t\ell^{-1})$ is expanded
in multiplicative characters. It is tempting to ask for a Weil--type bound for that phase over
primes, and an earlier draft of this remark did. That framing is \emph{wrong}, and the error is
worth recording. The phase depends only on $\ell\bmod p$, so
\[ \sum_{\ell\le X}e_p(s\ell^{-1})\ell^{-1-i\tau}
   \;=\;\sum_{a\not\equiv0}e_p(sa^{-1})\!\!\sum_{\substack{\ell\le X\\ \ell\equiv a\,(p)}}\!\!\ell^{-1-i\tau}, \]
an exact reorganisation with no loss; its discrete Fourier transform is a Gauss sum, so a ``direct''
estimate is the character expansion restated rather than an independent route. The Kloosterman
literature over primes (Fouvry--Michel, Fouvry--Kowalski--Michel, Garaev, Korolev, Irving) is
\emph{not} the relevant technology: every bound there is nontrivial only for modulus at least a fixed
power of $X$, and ours is at most a power of $\log X$. Note also that the sum is not $o(\log\log X)$
at fixed $p$: since $\sum_{a\not\equiv0}e_p(sa^{-1})=-1$, Mertens in progressions gives
$-\log\log X/(p-1)+O(1)$, measured as $-0.2502$, $-0.1667$, $-0.0830$, $-0.0105$ at
$p=5,7,13,101$ against $-1/(p-1)$
\textup{(\texttt{code/inverse\_phase\_check.gp})}. That failure is confined to bounded $p$ and is
harmless, $1/(p-1)$ being itself $o(1)$.

What the same computation does show is that the \emph{accounting} is lossy: measured,
$\max_s|\sum_{\ell\le10^6}e_p(s\ell^{-1})/\ell|$ is $1.04,1.17,1.35,1.32,1.33$ at
$p=13,101,503,1009,2003$, flat, while $\sqrt p$ grows from $3.6$ to $44.8$. So the triangle
inequality over characters discards almost everything, and the wall above is an artefact of that
step rather than an arithmetic obstruction. It can be avoided, and that is
Lemma~\ref{lem:swdirect}: applying Siegel--Walfisz to the residue classes directly pays a factor $r$
against an error $\exp(-c\sqrt{\log x})$ rather than $\sqrt r$ against a bound on $\log L$, which is
free at $r\le(\log N)^A$. The ceiling computed in this remark is therefore the ceiling of the
\emph{character--expansion} route only; Theorem~\ref{cor:a9rate} passes it by a quarter--power of
$\log N$. The measurement above, $\max_s|\cdot|$ flat near $1.3$ while $\sqrt p$ grows to $44.8$, was the numerical shadow of exactly this.
\end{remark}

\begin{theorem}[A rate for $\mathrm{A}9$]\label{cor:a9rate}
\textup{(Label: \textsc{proved}. Two blockers recorded against this statement are both retracted: Lemma~\ref{lem:localsq}'s
non--uniformity, an artefact of a superseded route that needs a $p^2$ count where
Proposition~\ref{prop:condrate} needs only first--power divisibility; and the regime mismatch after
Lemma~\ref{lem:reversal}, which was measured against $N$ rather than against this statement's own
target, and is smaller than it. The operative blocker was the pretentious distance, and
Lemma~\ref{lem:swdirect} supplies it.)}
For every $\delta>0$,
\[ \#\{m\le N:\ \mu^2(m)=1,\ m'-2m\ \text{is a positive perfect square}\}
    \;\ll_\delta\;N\,\frac{(\log\log N)^{1/2+\delta}}{(\log N)^{1/4}}, \]
and the same for $m'+2m$. The squarefree restriction is inherited from Theorem~\ref{thm:a9} and from Lemma~\ref{lem:charcancel}. The implied constant is
ineffective, through Siegel.
\end{theorem}
\begin{proof}
Run the sieve of Proposition~\ref{prop:condrate} at
$P=(\log N)^{1/4}(\log\log N)^{1/2}$, with $\mathcal P$ the primes in $(P/2,P]$ and
$\Pi=|\mathcal P|\gg P/\log P$; here $\log P\asymp\log\log N$. The expanded square and its five
terms are exactly as there, and we bound each after dividing by $\Pi^2$.

\emph{The distance, at these moduli.} Every modulus arising is $r=pq\le P^2\ll(\log N)^{1/2}\log\log N$,
so Lemma~\ref{lem:swdirect} applies with $A=1$. For the character sums it gives, at
$g_t(a)=\chi(a)e_r(ta^{-1})$ and hence $|\widehat g_0|=\sqrt r/\varphi(r)=o(1)$,
$M\ge(1-o(1))\log\log N$; for the deficit sums, at $g_{s,t}(a)=e_p(sa^{-1})e_q(ta^{-1})$ and hence
$|\widehat g_0|\le\max\bigl(\tfrac1{p-1},\tfrac1{q-1}\bigr)\le2/(P-2)=o(1)$, the same. In both cases
Hal\'asz (Montgomery's form, $T=\log N$) gives
$N(1+M)e^{-M}+N(\log N)^{-1/2}\ll N(\log N)^{-1+o(1)}+N(\log N)^{-1/2}\ll N(\log N)^{-1/2}$: the
floor binds.

\emph{The five terms.} The diagonal contributes the main term $N/\Pi\asymp N\log P/P\asymp
N(\log\log N)^{1/2}(\log N)^{-1/4}$, which is the stated bound. The character term is
$\Pi^{-2}\sum_{p\ne q}|T_{pq}|\le\max_{p\ne q}\sqrt{pq}\cdot\max_t|S_t|\ll 2P\cdot N(\log N)^{-1/2}
\asymp N(\log\log N)^{1/2}(\log N)^{-1/4}$, the same order, this is what fixes $P$, the balance
$N\log P/P\asymp PN(\log N)^{-1/2}$ giving $P\asymp(\log N)^{1/4}(\log\log N)^{1/2}$. The correction
$|E_{pq}|\le4N/P$ contributes $\ll N(\log N)^{-1/4}(\log\log N)^{-1/2}$, smaller by $\log\log N$. For
the two divisibility terms, note first that no $p^2$--count is needed and Lemma~\ref{lem:localsq}
does not enter: what the sieve requires is $\#\{m:p\mid a_m\}\ll N/\Pi$, weaker than $N/p$ by the
factor $\log P$, and the two--prime count needs nothing at all, since
$\#\{m:p\mid a_m,\ q\mid a_m\}\le\#\{m:p\mid a_m\}$ makes
$\sum_{p\ne q}\#\{\cdots\}\le\Pi\sum_p\#\{\cdots\}$. By the detector identity of
Lemma~\ref{lem:deficit} (whose $1/(pq)$ prefactor cancels the $pq-1$ non--principal terms exactly,
with no $\sqrt r$ amplification) and the distance bound above,
$\#\{m:p\mid a_m\}\le N/p+\max_{s\ne0}|S_s|\ll N/p+N(\log N)^{-1/2}\ll N/\Pi$, as required. Every
non--diagonal term is therefore at most the main term, and dividing by $\Pi^2$ gives the bound.
The proof in fact delivers the sharper exponent $\tfrac12$ on $\log\log N$; the $\delta$ in the
statement is a deliberate weakening, so that no reader has to track which constants are absolute.
The plus case is identical with $c=-2$.
\end{proof}

\emph{Why the exponent is $-1/2$ and not $-1/4$.} A rate of $\ll N\log\log N(\log N)^{-1/4}$, balancing at $P\asymp(\log N)^{1/4}$. Carrying the balance properly against the Hal\'asz floor gives
$P\asymp(\log N)^{1/4}(\log\log N)^{1/2}$ and a bound better by $(\log\log N)^{1/2}$. Compared with
Proposition~\ref{prop:condrate}, which remains the statement proved without
Lemma~\ref{lem:swdirect}, the gain is a full quarter--power of $\log N$: that proposition's ceiling
is $N(\log\log\log N)^{2+o(1)}/\log\log N$, and no choice of $P$ improves it while the distance is
estimated through the character expansion.

\begin{proof}[Sketch, conditional on a uniform-in-$r$ form of Lemma~\ref{lem:charcancel}]
Take $\mathcal{P}$ to be the primes in $(P,2P]$ with $P=P(N)$ to be chosen, so
$\Pi\sim P/\log P$. The three terms of Theorem~\ref{thm:a9} are then
\[ \frac{N}{\Pi}\ \ll\ \frac{N\log P}{P},\qquad
   \frac1{\Pi^2}\sum_{p\neq q}|T_{pq}(N)|\ \le\ \max_{p\neq q}|T_{pq}(N)|,\qquad
   \sum_{p\in\mathcal{P}}O\!\left(\frac{N}{p^2}\right)\ \ll\ \frac{N}{P\log P}. \]
Granting a uniform-in-$r$ distance bound, which Lemma~\ref{lem:charcancel} does \emph{not}
supply, Hal\'asz would give
$T_r(N)\ll\sqrt r\,N(\log N)^{-1/2}$, and $r=pq\le4P^2$ gives $\sqrt r\le2P$, so the middle term is
$\ll PN(\log N)^{-1/2}$. Balancing it against the first,
$N\log P/P=PN(\log N)^{-1/2}$, gives $P\asymp(\log N)^{1/4}$, at which both equal
$N\log\log N(\log N)^{-1/4}$ up to constants; the third term is then $O\bigl(N(\log N)^{-1/4}/\log\log N\bigr)$
and is absorbed.
\end{proof}

\begin{remark}[Where the exponent comes from]\label{rem:rateexponent}
The exponent $\tfrac14$ is not intrinsic to the arithmetic. It is what remains after balancing
against the floor $N(\log N)^{-1/2}$ in the Montgomery form of Hal\'asz's theorem, which holds
however large the pretentious distance $M$ becomes. Any improvement of that floor transfers here
without change: a floor of $N(\log N)^{-\alpha}$ yields the exponent $\alpha/2$. The measured
character sums are of size $\sqrt N$ rather than $N(\log N)^{-1/2}$
(Section~\ref{sec:sqsieve}), so the truth is far stronger than what the general theorem can deliver,
and the gap is entirely in the analytic input rather than in the sieve.
\end{remark}

The exponent $\tfrac14$ is an artefact of the general theorem, not of the arithmetic, and the gap
can be closed off exactly. Everything in the sieve except the character-sum input is elementary, so
the reachable exponent is a function of one measurable quantity: how $T_r(N)$ grows with $r$.

\begin{hypothesis}[Square-root cancellation, uniform in the modulus]\label{hyp:sq}
There is an absolute $A\ge0$ with $|T_r(N)|\ll r^{A}N^{1/2+o(1)}$ for every odd squarefree $r$.
\end{hypothesis}

\begin{theorem}[Conditional power saving, and the crux exactly]\label{thm:condpower}
Assume Hypothesis~\textup{\ref{hyp:sq}}. Then
\[ \#\{m\le N:\ \mu^2(m)=1,\ m'-2m\ \text{is a positive perfect square}\}\;\ll\;N^{1-\delta},
   \qquad \delta=\frac1{2(2A+1)}. \]
In particular $A=0$ gives $N^{1/2+o(1)}$, which is exactly the bound
$\mathrm{CRUX}\text{-}\mathrm{A}9$ asks for.
\end{theorem}
\begin{proof}
Run the sieve of Theorem~\ref{thm:a9} with $\mathcal{P}$ the primes in $(P,2P]$, $\Pi\sim P/\log P$,
and dispose of the $m$ with $p\mid a_m$ by the divisor bound rather than
Lemma~\ref{lem:localsq}: a perfect square $a_m$ has at most $\log(2N)/\log P$ prime factors in
$\mathcal{P}$, so $\sum_{p\in\mathcal{P}}(a_m/p)\ge\Pi/2$ once $\Pi\ge2\log(2N)/\log P$. Then
\[ \#\{\,\cdot\,\}\ \ll\ \frac{N}{\Pi}+\frac1{\Pi^2}\sum_{p\neq q}|T_{pq}(N)|
   \ \ll\ \frac{N\log P}{P}+P^{2A}N^{1/2+o(1)}, \]
since $pq\le4P^2$. Balancing the two terms gives $P=N^{1/(2(2A+1))}$ and the stated exponent; the
side condition $\Pi\ge2\log(2N)/\log P$ holds comfortably at that $P$. For $A=0$ take
$P=N^{1/2}$, whence both terms are $N^{1/2+o(1)}$.
\end{proof}

\begin{remark}[The hypothesis is measured, and $A$ reads as zero]\label{rem:Ameasured}
Hypothesis~\ref{hyp:sq} is not a wish. At $N=4\times10^7$, so $\sqrt N=6.325\times10^3$, the ratio
$|T_r(N)|/\sqrt N$ across twelve moduli spanning ten orders of magnitude is
\[
\begin{array}{lrrrrrr}
r\ \text{prime}: & 101 & 1009 & 10007 & 100003 & 10^6{+}3 & 10^7{+}19\\
|T_r|/\sqrt N:   & 7.38 & 0.17 & 0.84 & 1.07 & 0.05 & 0.15\\[2pt]
r=pq:            & 101909 & 10097063 & 1000730021 & 100003300009 & & \\
|T_r|/\sqrt N:   & 0.95 & 0.21 & 1.30 & 0.64 & &
\end{array}
\]
together with $2.34$ and $1.00$ at the squarefree $1155$ and $1113121$. The implied exponent
$A=\log(|T_r|/\sqrt N)/\log r$ is $+0.006$, $-0.004$, $+0.013$, $-0.018$ and $-0.000$ at the five
largest moduli; the outlying $+0.433$ occurs at $r=101$, where $\log r$ is too small for the ratio to
be meaningful. Note that the last two moduli \emph{exceed} $N$, and the sums are still of size
$\sqrt N$ \textup{(\texttt{code/sqsieve\_uniform.rs})}.

So the sharp form of $\mathrm{A}9$ is not blocked by the sieve, which is elementary, nor by the
arithmetic, which cooperates: it is blocked by the absence of a Weil-type bound for a character sum
whose argument is not a polynomial. That is a single, isolated, and precisely stated analytic
deficit, and Lemma~\ref{lem:symbolfact} has already reduced it to a sum over a multiplicative
function of $\sigma_r$.
\end{remark}

The obvious route to Hypothesis~\ref{hyp:sq} is the one already used in
Lemma~\ref{lem:charcancel}: expand by Gauss sums and bound each multiplicative sum
$\sum_{m\le N}h_t(m)$ separately, where square-root cancellation would follow from a Weil-type
input or, in the analytic form, from \textup{GRH} applied to
$\prod_\chi L(s,\chi)^{\widehat\alpha(\chi)}$. That route is closed. The reason is not the one a first look suggests.

\begin{proposition}[The decomposition destroys the cancellation]\label{prop:routeclosed}
Write $\alpha(u)=(u/r)e_r(t u^{-1})$, so that $h_t(p)=\alpha(p\bmod r)$. The principal Fourier
coefficient $\widehat\alpha(\chi_0)=\varphi(r)^{-1}\sum_u\alpha(u)$ is a Sali\'e sum; measured, it has
modulus exactly $\sqrt r$ before normalisation, at $r=1009$ and $r=10007$ giving $31.8$ and $100.0$
against $\sqrt r=31.76$ and $100.03$. It is therefore nonzero, and $\prod_\chi
L(s,\chi)^{\widehat\alpha(\chi)}$ retains a $\zeta$-type branch point at $s=1$. This does \emph{not}
by itself decide the matter, since the measured sums fall $50$ to $350$ times below what that branch
point predicts. What decides it is the growth in $N$: at $r=10007$, $t=1$,
\[
\begin{array}{lrrrrr}
N: & 10^6 & 3\times10^6 & 10^7 & 3\times10^7 & 6\times10^7\\
|\sum h_t|/\sqrt N: & 0.67 & 0.48 & 0.69 & 2.08 & 2.29
\end{array}
\]
a fitted exponent of $0.800$ over that range, with $|\sum h_t|/(N/\log N)$ flat near $0.005$. The
individual multiplicative sums thus grow strictly faster than $\sqrt N$
\textup{(\texttt{code/sqsieve\_route.rs})}.
\end{proposition}

\begin{remark}[Where the cancellation actually lives]\label{rem:whereitlives}
Since $T_r(N)$ itself measures at $\sqrt N$ across ten orders of magnitude in $r$
(Remark~\ref{rem:Ameasured}) while each $\sum_m h_t(m)$ measured here grows like $N^{0.8}$, the
cancellation in $T_r$ is not present frequency by frequency: it lies in the sum over $t$. The
magnitude is not marginal. At $r=1009$, $N=4\times10^7$, bounding each frequency separately and
summing trivially gives $\sqrt r\,N/\log N\approx7\times10^{7}$, against a measured
$|T_r|=1050$: a factor of about $7\times10^{4}$ is lost the moment the Gauss expansion is taken.

The deficit behind Hypothesis~\ref{hyp:sq} is therefore not a missing Weil bound for a single
character sum, which is how Remark~\ref{rem:Ameasured} first framed it. It is that the only
available decomposition of $T_r$ scatters the cancellation across frequencies, and no current
technique grips cancellation in that position. Stating it this way is more useful than the
Weil framing, because it says what would have to be new: an estimate for $T_r$ that never passes
through the multiplicative decomposition at all.
\end{remark}

Remark~\ref{rem:whereitlives} asked for an estimate for $T_r$ that never passes through the
multiplicative decomposition. The following is the mechanism such an estimate would use, isolated
and proved; the reduction of $T_r$ to it is \emph{not} established here, and we say precisely what
is missing.

\begin{proposition}[Kernel norm]\label{prop:bilinear}
Since $\sigma_r$ is additive, for coprime $d,e$ Lemma~\ref{lem:symbolfact} gives
\[ \Bigl(\frac{(de)'-2de}{r}\Bigr)=\Bigl(\frac{d}{r}\Bigr)\Bigl(\frac{e}{r}\Bigr)\,
   K\bigl(\sigma_r(d),\sigma_r(e)\bigr),\qquad K(u,v)=\Bigl(\frac{u+v-2}{r}\Bigr). \]
The matrix $K$ is Toeplitz in $u+v$, so $K(u,v)=\sum_\xi\widehat f(\xi)e_r(\xi u)e_r(\xi v)$
diagonalises it against the orthogonal system $(e_r(\xi u))_u$ of norm $\sqrt r$. As $\widehat f(\xi)$
is a normalised Gauss sum, $|\widehat f(\xi)|=r^{-1/2}$ for $\xi\neq0$ and $\widehat f(0)=0$, so every
singular value of $K$ equals $\sqrt r$ and
\[ \|K\|_{\mathrm{op}}=\sqrt r . \]
Consequently \emph{any} bilinear form $\sum_{d\sim D}\sum_{e\sim E}\alpha_d\beta_e\,
K(\sigma_r(d),\sigma_r(e))$ is at most $\sqrt r\,\|A\|_2\|B\|_2$, where
$A(u)=\sum_{d\sim D,\ \sigma_r(d)\equiv u}\alpha_d$ and $B$ is defined likewise. No multiplicative
decomposition is used; the only input is the modulus of a Gauss sum.
\end{proposition}

\begin{remark}[What is missing, stated exactly]\label{rem:missingdecomp}
Proposition~\ref{prop:bilinear} bounds bilinear forms, not $T_r$. Passing from one to the other
requires a decomposition of $\sum_{m\le N}$ into bilinear pieces, and two obstacles are real rather
than cosmetic.
\begin{enumerate}[label=\textup{(\roman*)},leftmargin=2.2em]
\item Summing over coprime pairs $(d,e)$ with $de=m$ weights each $m$ by $2^{\omega(m)}$: at
$N=2\times10^6$ the coprime-pair count is $10{,}014{,}190$ against $1{,}215{,}876$ squarefree $m$, a
factor $8.24$. Since the summand is \emph{signed}, no inequality relates $\sum_m F(m)$ to
$\sum_m 2^{\omega(m)}F(m)$ in either direction.
\item A single dyadic block $d\sim D$ sees only a minority of $m$: for $D=\sqrt N$ at
$N=2\times10^6$, $24.7\%$ \textup{(\texttt{code/sqsieve\_bilinear.rs})}.
\end{enumerate}
The natural repair is the largest-prime-factor split $m=pn$ with $p=P^{+}(m)$, which is unique and so
avoids \textup{(i)}; the $A$-side is then supported on primes, where $\|A\|_2^2\ll\pi(D)^2/\varphi(r)+\pi(D)$
follows from Brun--Titchmarsh uniformly in $r$, and the $B$-side is governed by
Theorem~\ref{thm:pairavg}. What remains is the coupling $P^{+}(n)<p$ and the dyadic bookkeeping, a
Type I/II analysis that we have not carried out. Until it is, the power saving below is
\textbf{CONJECTURAL}. The best unconditional rate proved here is
Theorem~\ref{cor:a9rate}'s $N(\log\log N)^{1/2+\delta}(\log N)^{-1/4}$; Proposition~\ref{prop:condrate}
is what the character--expansion route gives without Lemma~\ref{lem:swdirect}.
\end{remark}

\begin{conjecture}[Power saving, pending the decomposition]\label{conj:power34}
If the decomposition of Remark~\textup{\ref{rem:missingdecomp}} is carried out with the stated
bounds, then combining Proposition~\ref{prop:bilinear} with Theorem~\ref{thm:pairavg} and
Theorem~\ref{thm:condpower} at $P=N^{1/4}$ gives
\[ \#\{m\le N:\ \mu^2(m)=1,\ m'-2m\ \text{is a positive perfect square}\}\ \ll\ N^{3/4+o(1)} . \]
\end{conjecture}

The two ingredients that would feed it are themselves unconditional, and are recorded separately
because they do not depend on the missing step.


\begin{thebibliography}{9}
\bibitem{Barbeau1961} E.~J.~Barbeau, \emph{Remarks on an arithmetic derivative}, Canad.\ Math.\
Bull.\ \textbf{4} (1961), 117--122.
\bibitem{Barbeau1976} E.~J.~Barbeau, \emph{Computer challenge corner: Problem 477}, J.\ Recreational
Math.\ \textbf{9} (1976/77), 30.
\bibitem{ErdosGraham1980} P.~Erd\H{o}s and R.~L.~Graham, \emph{Old and New Problems and Results in
Combinatorial Number Theory}, Monogr.\ Enseign.\ Math.\ \textbf{28}, Geneva, 1980, p.~38.
\bibitem{GrahamEgyptian} R.~L.~Graham, \emph{Paul Erd\H{o}s and Egyptian fractions}, in \emph{Erd\H{o}s
Centennial}, Bolyai Soc.\ Math.\ Stud.\ \textbf{25}, Springer, 2013, pp.~289--309.
\bibitem{Watanabe2020} T.~Watanabe, \emph{New examples of the representation of $1$ by the sum of
reciprocals of semiprime numbers}, arXiv:2009.03275 (2020).
\bibitem{Bado2026b} I.-O.~Bado, \emph{Primary-pseudoperfect sectors in two-cycles of the
arithmetic derivative}, preprint (2026), communicated to the author.
\bibitem{BloomShifted} T.~F.~Bloom, \emph{Unit fractions with shifted prime denominators}, Proc.\
Roy.\ Soc.\ Edinburgh Sect.\ A (2024); arXiv:2305.02689.
\bibitem{UfnarovskiAhlander2003} V.~Ufnarovski and B.~\AA hlander, \emph{How to differentiate a
number}, J.\ Integer Seq.\ \textbf{6} (2003), Article 03.3.4.
\bibitem{Kovic2012} J.~Kovi\v{c}, \emph{The arithmetic derivative and antiderivative}, J.\ Integer
Seq.\ \textbf{15} (2012), Article 12.3.8.
\bibitem{MO320838} MathOverflow, \emph{Can the product of two sums of prime reciprocals be $1$?},
question 320838 (comments by R.~Israel and J.~Rosen, answer by ``Woett''), 2019.
\url{https://mathoverflow.net/questions/320838}.
\bibitem{ErdosProblems307} T.~F.~Bloom, \emph{Erd\H{o}s problem \#307},
\url{https://www.erdosproblems.com/307}.
\bibitem{OEIS_A003415} OEIS Foundation, \emph{Sequence A003415 (arithmetic derivative)},
\url{https://oeis.org/A003415}.
\bibitem{Pollack2018} P.~Pollack, \emph{On amicable tuples}, Illinois J. Math. \textbf{62} (2018),
  no.~1--4, 225--240; arXiv:1711.06847. (Explicit bounds for relatively prime amicable pairs,
  refining the finiteness of Artjuhov and Borho.)

\bibitem{teRiele1984} H.~J.~J.~te~Riele, \emph{Computation of all the amicable pairs below
$10^{10}$}, Math.\ Comp.\ \textbf{47} (1986), 361--368.
\bibitem{BHP2001} R.~C.~Baker, G.~Harman, and J.~Pintz, \emph{The difference between consecutive
primes, II}, Proc.\ London Math.\ Soc.\ (3) \textbf{83} (2001), 532--562.
\bibitem{Iwaniec1978} H.~Iwaniec, \emph{Almost-primes represented by quadratic polynomials},
Invent.\ Math.\ \textbf{47} (1978), 171--188.
\bibitem{LemkeOliver2012} R.~J.~Lemke Oliver, \emph{Almost-primes represented by quadratic
polynomials}, Acta Arith.\ \textbf{151} (2012), 241--261.
\bibitem{HardyWright} G.~H.~Hardy and E.~M.~Wright, \emph{An Introduction to the Theory of Numbers},
6th ed., Oxford Univ.\ Press, 2008.
\bibitem{Tenenbaum} G.~Tenenbaum, \emph{Introduction to Analytic and Probabilistic Number Theory},
3rd ed., Grad.\ Stud.\ Math.\ \textbf{163}, Amer.\ Math.\ Soc., 2015 (Erd\H{o}s--Wintner theorem).
\bibitem{Stothers} W.~W.~Stothers, \emph{Polynomial identities and Hauptmoduln}, Quart.\ J.\ Math.\
Oxford (2) \textbf{32} (1981), 349--370; R.~C.~Mason, \emph{Diophantine Equations over Function
Fields}, Cambridge Univ.\ Press, 1984.
\bibitem{Seva323194} Seva (\texttt{mathoverflow.net} user), \emph{A recursion with a number--theoretic
function}, MathOverflow question 323194 (2019). \url{https://mathoverflow.net/questions/323194}.
\bibitem{FewPrimes2025} \emph{Egyptian fractions for few primes}, arXiv:2507.03727 (July 2025).
\bibitem{Steinerberger2025} S.~Steinerberger, \emph{On an iterated arithmetic function problem of
Erd\H{o}s and Graham}, arXiv:2504.08023 (April 2025).
\bibitem{StretchedEG2026} \emph{A stretched--exponential bound for an Erd\H{o}s--Graham
unit--fraction problem}, arXiv:2607.04157 (July 2026).
\bibitem{PortFillings2026} \emph{Port fillings for primary pseudoperfect numbers}, arXiv:2605.21518
(May 2026).
\bibitem{Semiprime2026} \emph{Every natural number is a sum of distinct semiprime unit fractions},
arXiv:2606.15159 (June 2026).
\bibitem{ButskeJajeMayernik} W.~Butske, L.~M.~Jaje and D.~R.~Mayernik, \emph{On the equation
$\sum_{p\mid N}1/p+1/N=1$, pseudoperfect numbers, and perfectly weighted graphs}, Math.\ Comp.\
\textbf{69} (2000), 407--420.
\bibitem{Bado2026} I.~O.~Bado, \emph{Structural constraints for an Erd\H{o}s unit--fraction problem
over primes}, preprint, Aix--Marseille University (May 2026).
\url{https://doi.org/10.13140/RG.2.2.32245.54245}.
\bibitem{Giuga1950} G.~Giuga, \emph{Su una presumibile propriet\`a caratteristica dei numeri primi},
Ist.\ Lombardo Sci.\ Lett.\ Rend.\ Cl.\ Sci.\ Mat.\ Nat.\ (3) \textbf{14} (1950), 511--528.
\bibitem{BBBG1996} D.~Borwein, J.~M.~Borwein, P.~B.~Borwein, and R.~Girgensohn, \emph{Giuga's conjecture
on primality}, Amer.\ Math.\ Monthly \textbf{103} (1996), no.~1, 40--50.
\bibitem{BorweinWong1997} J.~M.~Borwein and E.~Wong, \emph{A survey of results relating to Giuga's
conjecture on primality}, CRM Proc.\ Lecture Notes \textbf{11} (1997), 13--27.
\bibitem{BJM2000} W.~Butske, L.~M.~Jaje, and D.~R.~Mayernik, \emph{On the equation
$\sum_{p\mid N}1/p+1/N=1$, pseudoperfect numbers, and perfectly weighted graphs}, Math.\ Comp.\
\textbf{69} (2000), no.~229, 407--420.
\bibitem{GrauOllerMarcen2012} J.~M.~Grau and A.~M.~Oller-Marc\'en, \emph{Giuga numbers and the
arithmetic derivative}, J.\ Integer Seq.\ \textbf{15} (2012), Article 12.4.1; arXiv:1103.2298.
\bibitem{SondowMacMillan2017} J.~Sondow and K.~MacMillan, \emph{Primary pseudoperfect numbers,
arithmetic progressions, and the Erd\H{o}s--Moser equation}, Amer.\ Math.\ Monthly \textbf{124} (2017),
no.~3, 232--240; arXiv:1812.06566.
\bibitem{GrauOllerMarcenSadornil2021} J.~M.~Grau, A.~M.~Oller-Marc\'en, and D.~Sadornil,
\emph{On $\mu$-Sondow numbers}, arXiv:2111.14211 (2021); Acta Math.\ Hungar.\ (2022).
\bibitem{Cox} D.~A.~Cox, \emph{Primes of the Form $x^2+ny^2$: Fermat, Class Field Theory, and Complex
Multiplication}, 2nd ed., Wiley, 2013.
\bibitem{OEIS_A054377} OEIS Foundation, \emph{Sequence A054377 (primary pseudoperfect numbers)},
\url{https://oeis.org/A054377}.
\bibitem{OEIS_A007850} OEIS Foundation, \emph{Sequence A007850 (Giuga numbers)},
\url{https://oeis.org/A007850}.
\bibitem{OEIS_A396915} OEIS Foundation, \emph{Sequence A396915 (composite squarefree $k$ with
$k'+2k$ a perfect square)}, \url{https://oeis.org/A396915}.
\bibitem{Evertse1984} J.-H.~Evertse, \emph{On sums of $S$-units and linear recurrences}, Compositio
Math.\ \textbf{53} (1984), 225--244.
\bibitem{FI1998} J.~Friedlander and H.~Iwaniec, \emph{The polynomial $X^2+Y^4$ captures its
primes}, Ann. of Math. \textbf{148} (1998), 945--1040.
\bibitem{HB2001} D.~R.~Heath-Brown, \emph{Primes represented by $x^3+2y^3$}, Acta Math.
\textbf{186} (2001), 1--84.
\bibitem{Stewart2013} C.~L.~Stewart, \emph{On divisors of Lucas and Lehmer numbers}, Acta Math.
\textbf{211} (2013), 291--314.
\bibitem{BCZ2003} Y.~Bugeaud, P.~Corvaja, and U.~Zannier, \emph{An upper bound for the G.C.D. of
$a^n-1$ and $b^n-1$}, Math. Z. \textbf{243} (2003), 79--84.
\bibitem{Merikoski2019} J.~Merikoski, \emph{On the largest prime factor of $n^2+1$}, J. Eur. Math.
Soc. \textbf{25} (2023), 1253--1284.
\bibitem{BGS2010} J.~Bourgain, A.~Gamburd, and P.~Sarnak, \emph{Affine linear sieve, expanders,
and sundry applications}, Invent. Math. \textbf{179} (2010), 559--644.
\bibitem{ESS2002} J.-H.~Evertse, H.~P.~Schlickewei, and W.~M.~Schmidt, \emph{Linear equations in
variables which lie in a multiplicative group}, Ann.\ of Math.\ (2) \textbf{155} (2002), 807--836.
\bibitem{Erdos1952} P.~Erd\H{o}s, \emph{On the sum $\sum_{k=1}^x d(f(k))$}, J.\ London Math.\ Soc.\
\textbf{27} (1952), 7--15.
\bibitem{BHV2001} Y.~Bilu, G.~Hanrot, and P.~M.~Voutier, \emph{Existence of primitive divisors of
Lucas and Lehmer numbers} (with an appendix by M.~Mignotte), J.\ Reine Angew.\ Math.\ \textbf{539}
(2001), 75--122.
\end{thebibliography}
\end{document}
